\documentclass[11pt,twoside,openright]{book}

\usepackage[T1]{fontenc}
\usepackage[utf8]{inputenc}
\usepackage{amsmath,amssymb,amsthm,mathtools,bm}
\usepackage{mathpazo}
\usepackage{microtype}
\usepackage[inner=1.35in,outer=1.1in,top=1.25in,bottom=1.35in]{geometry}
\usepackage{tikz}
\usepackage{booktabs}
\usepackage{tikz-cd}
\usetikzlibrary{arrows.meta,positioning,calc}
\usepackage{fancyhdr}
\usepackage[hidelinks]{hyperref}
\usepackage{cleveref}

\setlength{\parindent}{1.4em}
\setlength{\parskip}{0pt}
\linespread{1.07}
\emergencystretch=3em
\makeatletter
\def\cleardoublepage{\clearpage\if@twoside\ifodd\c@page\else\hbox{}\thispagestyle{empty}\newpage\fi\fi}
\makeatother

% ---------- headers ----------
\pagestyle{fancy}
\fancyhf{}
\fancyhead[LE]{\small\thepage\quad\textsc{Generative Film}}
\fancyhead[RO]{\small\nouppercase{\leftmark}\quad\thepage}
\renewcommand{\headrulewidth}{0.3pt}
\fancypagestyle{plain}{\fancyhf{}\fancyfoot[C]{\small\thepage}\renewcommand{\headrulewidth}{0pt}}

% ---------- theorem environments ----------
\theoremstyle{plain}
\newtheorem{theorem}{Theorem}[chapter]
\newtheorem{proposition}[theorem]{Proposition}
\newtheorem{lemma}[theorem]{Lemma}
\newtheorem{corollary}[theorem]{Corollary}
\theoremstyle{definition}
\newtheorem{definition}[theorem]{Definition}
\newtheorem{example}[theorem]{Example}
\newtheorem{principle}[theorem]{Principle}
\newtheorem{exercise}{Exercise}[chapter]
\theoremstyle{remark}
\newtheorem{remark}[theorem]{Remark}

% ---------- macros ----------
\newcommand{\R}{\mathbb{R}}
\newcommand{\N}{\mathbb{N}}
\newcommand{\Z}{\mathbb{Z}}
\newcommand{\Sc}{\mathcal{S}}       % scenario space
\newcommand{\Ren}{\mathsf{R}}       % rendering
\newcommand{\Img}{\mathcal{I}}      % image space
\newcommand{\Hist}{\mathcal{H}}     % histories
\newcommand{\Adm}{\mathcal{A}}      % admissible histories
\newcommand{\Path}{\mathbf{Path}}
\newcommand{\Set}{\mathbf{Set}}
\newcommand{\Ab}{\mathbf{Ab}}
\newcommand{\Open}{\mathrm{Open}}
\newcommand{\op}{\mathrm{op}}
\newcommand{\id}{\mathrm{id}}
\newcommand{\KL}{D_{\mathrm{KL}}}
\DeclareMathOperator{\dur}{dur}
\DeclareMathOperator{\Hom}{Hom}
\DeclareMathOperator{\supp}{supp}
\DeclareMathOperator{\Th}{Th}
\DeclareMathOperator{\Mod}{Mod}
\DeclareMathOperator{\Cl}{Cl}
\newcommand{\shuffle}{\mathbin{\sqcup\!\sqcup}}
\newcommand{\restr}[2]{#1|_{#2}}

\title{\Huge\bfseries Generative Film\\[0.6ex]
\Large Categories, Sheaves, and Types for Cinema}
\author{Flyxion\\[0.3ex]\small Independent Researcher}
\date{October 2026}

\begin{document}

\frontmatter
\maketitle

\chapter*{Preface}
\addcontentsline{toc}{chapter}{Preface}

Film theory has always been a theory of composition. Eisenstein asked what happens when two shots meet. Bazin asked what is lost when they are not allowed to. Metz asked what rules govern the order of segments, and Deleuze asked what kind of image results when the order stops obeying the movement of bodies. Each of these questions has the same grammatical shape: a collection of local pieces, a way of putting them together, and a claim about what the whole carries that the pieces do not.

Mathematics developed a precise language for exactly this shape during the twentieth century. Category theory describes composition and what is preserved by it. Sheaf theory describes how local data assemble into global data, and measures the obstruction when they cannot. Type theory describes dependency, the way later events may be licensed or forbidden by earlier ones. At the same time, generative models have turned film production into a problem of sampling from constrained spaces of histories, which makes these three languages practical rather than decorative.

This book is a first-year graduate text that develops the three languages side by side, with cinema as the running subject. It assumes mathematical maturity at the level of a first course in topology and linear algebra, and basic exposure to probability. Category theory, sheaf theory, and type theory are introduced from the ground up, in the amount needed, in the main chapters and in the appendices. No background in film studies is assumed beyond having watched films with attention, although the reader will profit from the works cited.

The book has seven parts. Part~I sets up scenario spaces and the algebra of editing. Part~II develops the categorical treatment of shots, cuts, and meaning. Part~III introduces sheaves and cohomology as the mathematics of continuity. Part~IV treats narrative dependency with dependent and linear types. Part~V turns to generation: histories, state, appearance, admissible continuation, and the multiplexing of one screen among several viewers. Part~VI addresses what a viewer can distinguish, what a transformation of a film must preserve, and how tropes and hierarchical commands can be given a typed semantics. Part~VII applies the apparatus to a set of films and a synthetic scene, and lists open problems.

Each chapter closes with exercises. Some are routine verifications, some extend a construction, and a few are open-ended. Claims are marked as theorems only when a proof is given or a standard reference is cited. Where a statement is a modeling hypothesis, it is labeled a principle and its status is stated plainly.

\tableofcontents

\mainmatter

\part{Film as Structured Motion}
\chapter{Film Theory as a Theory of Composition}
\label{ch:intro}

\section{Four questions about the cut}

The cut is the founding operation of cinema as an art distinct from recorded theater. Four influential bodies of theory can be read as four answers to what the cut does.

For Eisenstein, a cut is a collision. Two shots placed in sequence produce a meaning belonging to neither, and the editor's craft is the engineering of that surplus~\cite{eisenstein1949}. The Soviet montage tradition, including Kuleshov's well-known experiments on the juxtaposition of an actor's face with different following images, treats the whole as strictly richer than the sum of its parts.

For Bazin, the cut is a loss. The long take with deep focus preserves the spatial and temporal unity of the event and leaves interpretive freedom to the spectator, while montage substitutes the director's analysis for the world's ambiguity~\cite{bazin1967}. Here continuity is a value, and a cut is a cost that must be justified.

For Metz and the semiological tradition, a cut belongs to a syntax. Segments of a film are classified by how they articulate time and space, and some sequences of segments are well formed while others are not~\cite{metz1974}. The film is a string over an alphabet of segment types, and the question is which strings are admissible.

For Bordwell and the cognitive tradition, a cut is a cue within a perceptual and inferential process. The spectator builds a representation of a story world, and editing succeeds when it supports that construction~\cite{bordwell1985}. Deleuze, from a different direction, distinguishes the movement-image, in which cuts serve a sensory-motor chain of action, from the time-image, in which that chain has broken and the cut exposes time itself~\cite{deleuze1986,deleuze1989}.

These answers are not rivals in any simple sense. They address different aspects of one structure, and each can be restated in terms of three questions. What are the local pieces? How do they compose? What does the composite carry? The book is organized around supplying mathematically exact answers to all three, in a form that also tells a generative system what it must respect.

\section{Why mathematics, and why now}

Two developments make a formal treatment more than an exercise in analogy.

The first is the maturation of the relevant mathematics. Category theory gives a general account of composition: a collection of objects, a collection of composable arrows between them, and the requirement that composition be associative and unital~\cite{maclane1998,awodey2010}. Its value for the present subject is that it separates what is true of every composition from what is true of a particular one, and provides a vocabulary for structure-preserving translation between domains, for instance between a space of world states and a space of images. Sheaf theory, developed for geometry, formalizes the passage from local to global and attaches to every failure of that passage an algebraic invariant, a cohomology class~\cite{maclane1992,bott1982}. Type theory in the tradition of Martin-L\"of formalizes dependency: a type may depend on a value, so that what can happen next is a function of what has happened~\cite{martinlof1984,hottbook2013}.

The second is practical. Systems that generate video, interactive environments, and game-like worlds from learned models are now real objects~\cite{ha2018worldmodels,bruce2024genie,valevski2024gamengen}. A generative system differs from a camera in that it must be told, or must learn, what a legal continuation is. The more a system is asked to sustain a persistent world over long durations, the more it needs an explicit account of what is committed, what is derived, and what is merely displayed. Those distinctions are the content of Chapters~\ref{ch:layers} and~\ref{ch:admissibility}.

\section{The thesis in three sentences}

A film is a path through a space of world states, rendered to images, and cut at chosen points. The conditions under which pieces may be joined are describable as the existence of morphisms in a category, as the gluing of local sections of a sheaf, and as the well-typedness of a sequence of events in a dependent type theory. A generative film system is a sampler over the histories that satisfy those conditions.

The first sentence is developed in Chapters~\ref{ch:scenario} and~\ref{ch:editing}. The second spreads over Parts~II to IV. The third occupies Part~V.

\section{A running example}\label{sec:scene}

Abstract constructions are easiest to assess against a scene that can be varied at will. The book therefore returns, in bridge sections of most chapters, to a single synthetic scene. It is invented for teaching, it is not drawn from any film, and its times are stipulated by its specification and not measured from footage. Findings about it are findings about its constructed variations and do not transfer automatically to existing films.

The scene lasts ninety-six seconds in a kitchen. Two characters, $A$ and $B$, are present, and a letter lies on the table. $A$ knows from the start that the letter is forged and $B$ does not. $A$ enters at $t=0$ and $B$ greets $A$ at $t=6$. $A$ glances at the letter at $t=18$, and street noise passes outside between $t=20$ and $t=22$. A pause of seven seconds begins at $t=31$. At $t=38$ $B$ asks where the letter came from, and at $t=46$ $A$ deflects. $B$'s phone buzzes at $t=52$. At $t=58$ $A$ tells $B$ that the letter is forged. A reaction shot of $B$ runs from $t=62$ to $t=66$. $B$ leaves, and a door slams at $t=78$. Music enters at $t=80$ and continues to the end of the scene at $t=96$. \Cref{fig:timeline} shows the timeline.

\begin{figure}[ht]
\centering
\begin{tikzpicture}[x=0.135cm, font=\scriptsize]
\draw[thick] (0,0) -- (96,0);
\fill[gray!25] (31,-0.18) rectangle (38,0.18);
\fill[gray!45] (62,-0.18) rectangle (66,0.18);
\fill[gray!15] (80,-0.18) rectangle (96,0.18);
\foreach \t/\l in {0/enter,18/glance,38/question,52/buzz,62/reaction,80/music}{
  \draw (\t,0) -- (\t,0.35);
  \node[anchor=south west, rotate=45, inner sep=1pt] at (\t,0.35) {\l\,(\t)};
}
\foreach \t/\l in {6/greeting,20/street noise,31/pause,46/deflection,58/reveal,78/door slam,96/end}{
  \draw (\t,0) -- (\t,-0.35);
  \node[anchor=north east, rotate=45, inner sep=1pt] at (\t,-0.35) {\l\,(\t)};
}
\end{tikzpicture}
\caption{Timeline of the forged-letter scene, in seconds. Shaded bands mark the pause (31 to 38), the reaction shot (62 to 66), and the music (80 to 96).}
\label{fig:timeline}
\end{figure}

Four questions about this scene recur. Which of its changes are physical movement, which are changes in what a character knows, which are changes in what the audience has been told, and which are changes in sensory intensity? Which edits of the same footage differ in pacing and which differ only in appearance? What may be altered, for instance by removing the pause or quieting the sound, without changing what matters to a viewer? And how would a person direct such a scene without specifying every movement? The bridge sections answer these in turn, each time with something more exact to observe, change, or investigate.

\section{A map of the argument}

The dependency among the parts is shown in \Cref{fig:map}. Part~I is needed throughout. Parts~II, III, and IV can be read in any order after Part~I, though the examples in Part~III use categorical language and the examples in Part~V draw on all three.

\begin{figure}[ht]
\centering
\begin{tikzpicture}[
  node distance=1.2cm and 1.6cm,
  box/.style={draw, rounded corners=2pt, minimum width=3.0cm, minimum height=0.9cm, align=center, font=\small},
  arr/.style={-{Stealth[length=2.2mm]}, thick}
]
\node[box] (I) {I. Scenario spaces\\and editing};
\node[box, below left=of I] (II) {II. Categories};
\node[box, below=of I] (III) {III. Sheaves};
\node[box, below right=of I] (IV) {IV. Types};
\node[box, below=of III] (V) {V. Generation};
\node[box, below=of V] (VI) {VI. Synthesis};
\draw[arr] (I) -- (II);
\draw[arr] (I) -- (III);
\draw[arr] (I) -- (IV);
\draw[arr] (II) -- (V);
\draw[arr] (III) -- (V);
\draw[arr] (IV) -- (V);
\draw[arr] (V) -- (VI);
\end{tikzpicture}
\caption{Dependencies among the parts of the book.}
\label{fig:map}
\end{figure}

\section{Conventions}

Mathematical objects are introduced in definitions and their basic properties are stated as propositions. Statements asserting something about cinema or perception, rather than about the mathematical model, are called principles and are not proved. A principle is a hypothesis that the model makes testable, and the text indicates in each case what would count against it.

Films are cited by title and year. Where a film is used as an example, the claim made about it concerns the model's reading of a feature the film is widely known for, such as the use of long takes or looped time, and no claim is made about the intentions of its makers.

Notation is collected in Appendix~\ref{app:notation}.

\section*{Exercises}

\begin{exercise}
Choose a sequence of three shots from a film you know well. For each of the four theoretical traditions above, state in one sentence what that tradition would take the cuts in the sequence to be doing. Identify which sentences conflict and which merely address different aspects.
\end{exercise}

\begin{exercise}
The three questions of this chapter were: what are the pieces, how do they compose, what does the composite carry. Apply them to a non-film domain, such as a proof in a textbook or a piece of chamber music, and note which of the three has the least obvious answer.
\end{exercise}

\chapter{Scenario Spaces and Films as Paths}
\label{ch:scenario}

A camera pointed at a staged event records one of many possible views of one of many possible world states. To reason about film mathematically, the first task is to name the space in which those states live and to say what a film is in relation to it. This chapter does so, and then adds just enough geometry to state a hypothesis about camera movement.

\section{Scenario spaces}

\begin{definition}[Scenario space]
A \emph{scenario space} is a pair $(\Sc,\Ren)$ where $\Sc$ is a smooth manifold, possibly with corners, and $\Ren:\Sc\to\Img$ is a continuous map into a metric space $\Img$ of images. The elements of $\Sc$ are \emph{scenarios} or \emph{world states}, and $\Ren$ is the \emph{rendering}.
\end{definition}

In the standard example $\Img=L^2(\Omega;\R^3)$ for a rectangle $\Omega\subset\R^2$, so that an image is a color field on the screen. The choice of $L^2$ is a convenience that makes $\Img$ a Hilbert space and is not essential to anything below. What matters is that $\Ren$ is generally far from injective: many world states produce the same image, and this is the source of most of the interesting phenomena in later chapters.

A concrete scenario space has a product decomposition
\[
\Sc=\mathcal{C}\times\mathcal{A}\times\mathcal{L}\times\mathcal{E},
\]
where $\mathcal{C}=SE(3)\times(0,\infty)$ records camera pose and focal length, $\mathcal{A}$ records the configurations of the actors, $\mathcal{L}$ the lighting, and $\mathcal{E}$ the remaining environment. The camera factor is the only one that is directly under the cinematographer's control during a take, and the other factors are under the control of the production designer, the actors, and the physics of the set.

\begin{example}
A two-person dialogue in a room has $\mathcal{A}=Q_1\times Q_2$, where $Q_i$ is the configuration space of the $i$th actor's body, and $\mathcal{L}$ a space of light-source parameters. Fixing everything but the camera gives the subspace $\mathcal{C}\times\{(a,\ell,e)\}$, which is the space of all possible framings of one frozen moment.
\end{example}

\section{Shots and films}

Time enters through paths. To avoid the awkwardness that concatenation of paths on fixed intervals is not strictly associative, we use Moore paths, which carry their own duration.

\begin{definition}[Shot]
A \emph{shot} is a pair $(\tau,\gamma)$ with $\tau>0$ and $\gamma:[0,\tau]\to\Sc$ continuous and piecewise $C^1$. We write $\dur(\gamma)=\tau$.
\end{definition}

\begin{definition}[Film]
A \emph{film} is a finite sequence $f=(\gamma_1,\dots,\gamma_n)$ of shots. Its \emph{rendered form} is the sequence of image-valued paths $\Ren\circ\gamma_k:[0,\tau_k]\to\Img$. Its total duration is $\sum_k\tau_k$. The \emph{seams} of $f$ are the pairs $(\gamma_k(\tau_k),\gamma_{k+1}(0))$ for $1\le k<n$.
\end{definition}

A seam is \emph{continuous} when the two endpoints coincide in $\Sc$. A seam is a \emph{match} when they are distinct in $\Sc$ but close in $\Img$, which is the situation of a match cut and of the hidden cuts in films assembled to appear as one take. A seam is a \emph{hard cut} when both the scenario endpoints and their renderings are far apart. The three cases are distinguished by the numbers
\[
\delta_\Sc=d_\Sc\bigl(\gamma_k(\tau_k),\gamma_{k+1}(0)\bigr),\qquad
\delta_\Img=d_\Img\bigl(\Ren\gamma_k(\tau_k),\Ren\gamma_{k+1}(0)\bigr),
\]
for any chosen metric $d_\Sc$ on $\Sc$. Continuity is $\delta_\Sc=0$, which implies $\delta_\Img=0$. Matching is $\delta_\Img$ small with $\delta_\Sc$ not small. Because $\Ren$ is continuous, small $\delta_\Sc$ forces small $\delta_\Img$ locally, but the converse fails, and its failure is exactly what a match cut exploits.

\begin{proposition}\label{prop:seam-lipschitz}
If $\Ren$ is $L$-Lipschitz with respect to $d_\Sc$ and $d_\Img$, then $\delta_\Img\le L\,\delta_\Sc$ at every seam.
\end{proposition}

\begin{proof}
Immediate from the definition of the Lipschitz condition applied to the two endpoints.
\end{proof}

The proposition says that a viewer cannot see a cut whose scenario jump is small, which is the principle behind continuity editing. The interesting direction, and the one that makes cinema possible, is the converse failure: a cut with large $\delta_\Sc$ and small $\delta_\Img$ is invisible to the eye and nonetheless transports the story to a different world state.

\section{A perceptual geometry}

To speak of natural camera motion, a metric on $\Sc$ is needed that measures how different two nearby scenarios look to a viewer, not how different they are as physical configurations. Information geometry supplies one~\cite{amari2016}.

Let $X$ be a finite set of screen regions, and suppose each scenario $s$ induces a probability distribution $p(\cdot\mid s)$ on $X$ describing where a viewer's attention falls, with $s\mapsto p(\cdot\mid s)$ smooth and strictly positive. The \emph{Fisher--Rao pullback metric} is
\[
g_s(u,v)=\sum_{x\in X}p(x\mid s)\,\bigl(\partial_u\log p(x\mid s)\bigr)\bigl(\partial_v\log p(x\mid s)\bigr),\qquad u,v\in T_s\Sc.
\]
It is positive semidefinite everywhere, and positive definite exactly where the map $s\mapsto p(\cdot\mid s)$ is an immersion. Directions in which attention does not change are null directions of $g$, and these are the scenario changes that a viewer does not perceive as change.

Given $g$, and a smooth function $U:\Sc\times[0,T]\to\R$ representing a narrative cost, one defines the action of a shot by
\[
\mathcal{E}(\gamma)=\int_0^{\tau}\Bigl(\tfrac12\,g_{\gamma(t)}(\gamma'(t),\gamma'(t))+U(\gamma(t),t)\Bigr)\,dt.
\]
The first term charges for perceptual speed, and the second for passing through states that are narratively unwarranted at time $t$.

\begin{proposition}[Equation of natural motion]\label{prop:EL}
Suppose $g$ is positive definite, so that $(\Sc,g)$ is a Riemannian manifold. A shot $\gamma$ of class $C^2$ (we restrict to this class, although shots are only piecewise $C^1$) is a critical point of $\mathcal{E}$ among shots with the same endpoints and duration if and only if
\[
\nabla_{\gamma'}\gamma'=\operatorname{grad}_g U(\cdot,t)\big|_{\gamma(t)},
\]
where $\nabla$ is the Levi-Civita connection of $g$. In particular, if $U\equiv0$, critical points are exactly the geodesics of $g$ parametrized with constant speed.
\end{proposition}

\begin{proof}
Let $\gamma_\epsilon$ be a smooth variation with variation field $\xi$ vanishing at $t=0$ and $t=\tau$. Differentiating $\mathcal{E}(\gamma_\epsilon)$ at $\epsilon=0$ and using metric compatibility of $\nabla$ gives
\[
\frac{d}{d\epsilon}\Big|_0\mathcal{E}(\gamma_\epsilon)=\int_0^\tau\Bigl(g(\gamma',\nabla_t\xi)+dU(\xi)\Bigr)dt.
\]
Integrating the first term by parts and using $\xi(0)=\xi(\tau)=0$ yields
\[
\int_0^\tau g\bigl(-\nabla_{\gamma'}\gamma'+\operatorname{grad}U,\ \xi\bigr)\,dt.
\]
This vanishes for every $\xi$ if and only if the integrand's first slot vanishes identically, by the fundamental lemma of the calculus of variations. For $U\equiv0$ the equation is the geodesic equation, and geodesics have constant speed.
\end{proof}

The sign in the equation deserves a remark. A critical shot accelerates toward regions of higher narrative cost, which is to say that it hurries through them, since time spent there is charged at the higher rate. Lingering is reserved for narratively warranted regions.

\begin{principle}[Geodesic principle of cinematic naturalness]\label{pr:geodesic}
Camera movements that viewers experience as natural are, to good approximation, critical points of an action of the form $\mathcal{E}$ for a perceptual metric $g$ of Fisher--Rao type and a narrative cost $U$.
\end{principle}

This is a hypothesis about perception and craft, not a theorem. It would be undermined by evidence that the camera paths experienced as natural systematically fail to satisfy the equation of \Cref{prop:EL} for any metric in the family, or that viewers' judgments of naturalness are insensitive to deviations from critical paths. The proposition is nonetheless useful as a specification: it tells a generative system what family of paths to prefer, and it makes the hypothesis operational.

\section{Bridge: four coordinates of one scene}

Return to the forged-letter scene of Section~\ref{sec:scene}. A watcher of the scene tracks at least four different things that a single point of the scenario space does not separate. \emph{Physical activity} $m(t)$ is the speed of the scenario path in a chosen metric: it is large when $A$ enters and when $B$ leaves, and near zero through the pause. \emph{Character knowledge} $k(t)$ is the set of facts each character holds: it changes once, at $t=58$, when $B$ learns that the letter is forged. \emph{Audience discovery} $d(t)$ is the set of facts the audience has been given, which in the baseline version also changes at $t=58$, since nothing earlier settles that $A$ knows. \emph{Sensory intensity} $i(t)$ is the non-dialogue sound level, which rises at the street noise, the door slam, and the music.

These four processes are independent in the following sense: any one of them can be made to change while the others are held fixed. The door slam raises $i$ while $k$ and $d$ stay constant, and the reveal changes $k$ and $d$ while $m$ is small. Consider two edits of the same footage that share dialogue and timing. The first has fourteen cuts in ninety-six seconds, a rate of about $0.146$ cuts per second. The second has four cuts, a rate of about $0.042$ cuts per second, which is lower by a factor of $3.5$. Both place the question at $t=38$ and the reveal at $t=58$, so each contains the same two events that change the semantic state of the story. The cut rate is a descriptor of the visual channel and says nothing about narrative progress, which is the same in both. Pacing is accordingly better described as the vector of processes $(m,k,d,i)$ together with the cut process than as a single speed. Chapter~\ref{ch:equivalence} asks how far each coordinate can be altered before a viewer notices.

\section*{Exercises}

\begin{exercise}
Let $\Sc=\R^2$ with $g$ the Euclidean metric and $U(x,y,t)=\kappa\,e^{-(x^2+y^2)}$ for $\kappa>0$. Write the equation of \Cref{prop:EL} explicitly and describe qualitatively the solutions with endpoints on opposite sides of the origin.
\end{exercise}

\begin{exercise}
Show that if $\Ren$ factors through a quotient $\Sc\to\Sc/G$ by a group action, then every $G$-orbit consists of scenarios that are indistinguishable in $\Img$. Give a filmic example with $G\cong\Z/2$.
\end{exercise}

\begin{exercise}
In \Cref{prop:EL}, replace the narrative cost by one that depends on velocity, $U(s,v,t)$. State the resulting Euler--Lagrange equation. Why might a velocity-dependent cost be natural for modeling the preference for smooth pans over jerky ones?
\end{exercise}

\chapter{Editing as Algebra}
\label{ch:editing}

An editor works with a stock of shots and a small repertoire of operations on them. This chapter isolates the operations and the algebraic structures they generate, and states precisely where the theory of films leaves the territory of ordinary algebra, which is where the first genuinely cinematic phenomena begin.

\section{The free monoid of cuts}

Let $\Sigma$ be a set of shots. A film assembled without further processing is a finite word over $\Sigma$, and the set of all such films is the free monoid $\Sigma^*$ under concatenation, with the empty film as unit. Duration is a monoid homomorphism
\[
\dur:\Sigma^*\to(\R_{\ge0},+),
\]
since the duration of a concatenation is the sum of durations. Every statement about timing that is independent of content, such as total running time or the position of a given shot on the timeline, factors through this homomorphism.

Concatenation of arbitrary shots is always defined, and this reflects a basic fact about the medium: any two shots can be cut together. The theory of continuity concerns the special cases in which the cut is not seen, and those require the more refined structure of the next section.

\section{Seamless films and the Moore path category}

Let $\Sc$ be a scenario space. A \emph{Moore path} in $\Sc$ is a pair $(\tau,\gamma)$ with $\tau\ge0$ and $\gamma:[0,\tau]\to\Sc$ continuous. Two Moore paths $(\tau_1,\gamma_1)$ and $(\tau_2,\gamma_2)$ with $\gamma_1(\tau_1)=\gamma_2(0)$ have a concatenation $(\tau_1+\tau_2,\gamma_1\cdot\gamma_2)$ defined by running $\gamma_1$ and then $\gamma_2$.

\begin{proposition}\label{prop:moore}
There is a category $\Path(\Sc)$ whose objects are the points of $\Sc$ and whose morphisms $s\to s'$ are Moore paths from $s$ to $s'$, with composition given by concatenation and identities given by the constant paths of duration $0$. Composition is strictly associative and strictly unital.
\end{proposition}

\begin{proof}
For composable $(\tau_1,\gamma_1),(\tau_2,\gamma_2),(\tau_3,\gamma_3)$, both bracketings of the triple concatenation have duration $\tau_1+\tau_2+\tau_3$ and agree pointwise: for $t\in[0,\tau_1]$ the value is $\gamma_1(t)$, for $t\in[\tau_1,\tau_1+\tau_2]$ it is $\gamma_2(t-\tau_1)$, and so on. This is where durations are essential: with paths on the fixed interval $[0,1]$, the two bracketings reparametrize differently and agree only up to reparametrization. The unit laws hold because a duration-$0$ path contributes nothing.
\end{proof}

A \emph{seamless film}, which is to say a single continuous take possibly composed of several recorded segments, is a morphism in $\Path(\Sc)$. The duration functor $\dur:\Path(\Sc)\to(\R_{\ge0},+)$, viewing the monoid as a one-object category, is a functor. The inclusion of seamless films into all films is not a functor into $\Sigma^*$ viewed as a category in the obvious way, because concatenation in $\Sigma^*$ is total while in $\Path(\Sc)$ it is partial. The relationship is that $\Path(\Sc)$ is the subcategory of continuity-respecting composition inside the totalizing structure of the monoid.

Bazin's preference for the long take is thus a preference for staying inside $\Path(\Sc)$, in which every composite is again a path. A cut leaves that category and enters the larger monoid, and the cost of leaving is what later chapters will measure.

\section{Parallel editing and the shuffle}

Cross-cutting between two lines of action, as in a chase intercut with a rescue, produces a film that is an interleaving of two shot sequences. If $u=u_1\cdots u_m$ and $v=v_1\cdots v_n$ are words, an \emph{interleaving} is a word $w$ of length $m+n$ whose letters can be partitioned into two subsequences reading $u$ and $v$ respectively. The set of interleavings of $u$ and $v$ is the \emph{shuffle} $u\shuffle v$.

\begin{proposition}\label{prop:shuffle-count}
Counting interleavings by position of origin, $\lvert u\shuffle v\rvert\le\binom{m+n}{m}$, with equality when all letters of $u$ and $v$ are pairwise distinct. The shuffle is commutative and associative as an operation on finite sets of words: $u\shuffle v=v\shuffle u$ and $(u\shuffle v)\shuffle w=u\shuffle(v\shuffle w)$ in the sense that both equal the set of words that can be partitioned into three subsequences reading $u$, $v$, $w$.
\end{proposition}

\begin{proof}
An interleaving is determined by choosing which $m$ of the $m+n$ positions are occupied by letters of $u$; the remaining positions carry $v$ in order. This gives at most $\binom{m+n}{m}$ words, with distinct choices giving distinct words when all letters are distinct. Commutativity is the symmetry of the definition in $u$ and $v$. For associativity, a word in either bracketing admits a partition of its positions into three subsequences reading $u$, $v$, $w$, and conversely such a partition determines a word in both bracketings.
\end{proof}

The count matters to the generative setting. For two storylines of ten shots each there are $\binom{20}{10}=184756$ distinct pure interleavings, so the editor of a cross-cut sequence is selecting from a large space. Almost all of it is ruled out by narrative constraints such as causal order and rhythm, which is the work of Part~IV.

\section{Meaning is not a homomorphism}

The central difficulty in film theory, from the Kuleshov experiments onward, is that the meaning of a concatenation is not determined by the meanings of its parts. Make this precise by letting $M$ be a set of meanings equipped with a way of combining them.

Take $M$ to be a \emph{preordered monoid} $(M,\le,\otimes,1)$, so $\le$ is a preorder, $\otimes$ is associative with unit $1$, and $\otimes$ is monotone in each argument. A natural reading of $x\le y$ is that $y$ contains at least the information of $x$.

\begin{definition}[Compositional and superadditive meaning]
An \emph{interpretation} is a function $\mu:\Sigma^*\to M$. It is \emph{compositional} if $\mu(uv)=\mu(u)\otimes\mu(v)$ and $\mu(\varepsilon)=1$, so that $\mu$ is a monoid homomorphism. It is \emph{superadditive} if
\[
\mu(u)\otimes\mu(v)\le\mu(uv)\quad\text{and}\quad1\le\mu(\varepsilon).
\]
\end{definition}

A superadditive interpretation is the same thing as a lax monoidal functor from the discrete monoidal category $\Sigma^*$ to the monoidal preorder $M$. The Kuleshov effect and Eisenstein's thesis that montage generates meaning not present in the parts are the claim that real viewer interpretations are superadditive and strictly so for suitable pairs. The \emph{montage surplus} of a pair is the failure of equality, which cannot be expressed as a number in general but is meaningful as the statement $\mu(u)\otimes\mu(v)<\mu(uv)$.

If instead $M$ is an abelian group $(M,+)$, one can measure the surplus directly:
\[
\varepsilon(u,v)=\mu(uv)-\mu(u)-\mu(v).
\]
This function satisfies the cocycle identity $\varepsilon(u,v)+\varepsilon(uv,w)=\varepsilon(v,w)+\varepsilon(u,vw)$, because both sides equal $\mu(uvw)-\mu(u)-\mu(v)-\mu(w)$. It is thus a 2-cocycle on the monoid $\Sigma^*$ with values in $M$, and the identity exhibits it as the coboundary of $\mu$, viewed as a 1-cochain, up to the sign convention $\varepsilon=-\delta\mu$ (with $\mu(\varepsilon)=0$ if normalized cochains are used).

\begin{proposition}\label{prop:no-higher}
Every 2-cocycle on a free monoid with values in an abelian group $M$ (with trivial action) is a coboundary. Consequently, for abelian-group-valued meaning, the montage surplus carries no cohomological information beyond $\mu$ itself.
\end{proposition}

The proposition follows from the fact that the free monoid $\Sigma^*$ has cohomological dimension at most one. One way to see this is that the augmentation ideal of the monoid ring $\Z[\Sigma^*]$ is a free module on the generators, so that $0\to\Z[\Sigma^*]^{(\Sigma)}\to\Z[\Sigma^*]\to\Z\to0$ is a free resolution of length one, and cohomology in degrees above one therefore vanishes; the analogous statement for free groups is standard~\cite{brown1982}. A direct proof is also short: given a 2-cocycle $\varepsilon$, define $\mu$ on generators arbitrarily, put $\mu(\varepsilon_{\text{empty}})=0$, and extend recursively by $\mu(ua)=\mu(u)+\mu(a)+\varepsilon(u,a)$ for generators $a$; the cocycle identity then gives $\mu(uv)=\mu(u)+\mu(v)+\varepsilon(u,v)$ for all words, by induction on the length of $v$. We omit the details. Its message is methodological. Over the free monoid of films, montage effects are not an obstruction of the global kind that sheaf theory will detect in Part~III. They are a failure of $\mu$ to be a homomorphism, measured by a coboundary, and they live entirely in the choice of $\mu$. The nontrivial global phenomena in cinema, such as impossible spaces and time loops, arise when the base is not a free monoid of cuts but a space with topology, and that is why the next parts introduce sheaves.

\section{Bridge: how many edits does one scene allow?}

Suppose the forged-letter scene was covered by three cameras, and consider edits that cut at whole-second times. An edit with three cuts selects three of the $95$ interior grid points, which can be done in $\binom{95}{3}=138{,}415$ ways, and assigns one of three cameras to each of the four resulting segments with no camera repeated on adjacent segments, which can be done in $3\cdot2^3=24$ ways. The space of three-cut edits therefore has $138{,}415\times24=3{,}321{,}960$ elements, each a distinct word over the alphabet of segment-camera pairs. Almost all of them are ruled out by constraints that no editor states explicitly: a cut should not interrupt an utterance, the reveal line should be on a shot that shows $A$ or both characters, and the reaction shot should follow the reveal and not precede it. A practical editor navigates this space by such constraints, which are the subject of Parts~IV and~V.

The two edits compared in the previous chapter, with fourteen and four cuts, are words of different lengths over the same alphabet, and the length of a word records only the density of cuts. Whether the reaction shot at $t=62$ follows the reveal at $t=58$ or precedes it is a difference in the order of letters, and changes the interpretation $\mu$ of the whole. If $\mu$ were compositional, the meaning of the pair would be the product $\mu(\text{reveal})\otimes\mu(\text{reaction})$ of the meanings of the parts, with nothing contributed by their adjacency. The reaction shot, however, is read as a response to what has just been said, a meaning present in neither part alone. This is superadditive behavior in the sense of the definition above, in concrete form.

\section*{Exercises}

\begin{exercise}
Verify that $\Path(\Sc)$ fails to be a category if paths are taken on the fixed interval $[0,1]$ with concatenation by reparametrization, in the sense that associativity holds only up to a reparametrization. Which weaker structure results?
\end{exercise}

\begin{exercise}
Compute $\lvert u\shuffle v\rvert$ when $u=ab$ and $v=ac$, and explain why the bound of \Cref{prop:shuffle-count} fails to be attained.
\end{exercise}

\begin{exercise}
Let $M=([0,\infty),\le,+,0)$ with $\mu(w)$ the number of distinct emotional associations the audience forms from film $w$. Describe a hypothetical experiment, in the style of Kuleshov's, that would test superadditivity for a given pair of shots.
\end{exercise}

\begin{exercise}
Prove that the set of interleavings of $u$ and $v$ is closed under taking prefixes of the interleaved words in the sense that every prefix of an interleaving of $u$ and $v$ is an interleaving of a prefix of $u$ and a prefix of $v$.
\end{exercise}


\part{Categories of Shots and Meanings}
\chapter{Categories of Shots}
\label{ch:categories}

Chapter~\ref{ch:editing} met the category $\Path(\Sc)$ in passing. This chapter develops the categorical vocabulary needed to speak about film, namely categories, functors, and monoidal structure, and uses it to give a precise account of the relation between a world and its images and of the simultaneous composition of picture and sound. Appendix~\ref{app:cat} collects the definitions; the reader meeting the subject for the first time should read it alongside this chapter.

\section{Categories, functors, and the rendering functor}

A \emph{category} $\mathbf{C}$ consists of objects, morphisms between them, an associative composition, and identities. A \emph{functor} $F:\mathbf{C}\to\mathbf{D}$ assigns objects to objects and morphisms to morphisms so as to preserve composition and identities~\cite{maclane1998}.

\begin{proposition}[Rendering as a functor]\label{prop:rendering-functor}
Let $(\Sc,\Ren)$ be a scenario space. Then $\Ren$ induces a functor $\Ren_*:\Path(\Sc)\to\Path(\Img)$ with $\Ren_*(s)=\Ren(s)$ on objects and $\Ren_*(\tau,\gamma)=(\tau,\Ren\circ\gamma)$ on morphisms.
\end{proposition}

\begin{proof}
Since $\Ren$ is continuous, $\Ren\circ\gamma$ is a Moore path in $\Img$ from $\Ren(\gamma(0))$ to $\Ren(\gamma(\tau))$. Composition and identities are preserved because post-composition with $\Ren$ commutes with concatenation of paths and takes constant paths to constant paths.
\end{proof}

The proposition says that what the audience sees of a seamless take is determined by what happens in the world, in a way that respects the joining of takes. When $\Ren$ has positive-dimensional fibers, $\Ren_*$ fails to be faithful, which is the technical statement that different takes can look identical: two distinct morphisms of $\Path(\Sc)$ may have the same image. (Non-injectivity alone does not force this: a covering map is non-injective, yet paths lift uniquely.) It typically fails to be full as well, because most image paths do not arise from continuous paths of world states, which is the formal content of the idea that cinema can show what no physical set could produce.

\begin{example}[Faithful but not essentially surjective]
Suppose $\Ren$ is an embedding. Then $\Ren_*$ is faithful, because $\Ren\circ\gamma=\Ren\circ\gamma'$ implies $\gamma=\gamma'$, and the image of $\Ren_*$ consists of paths in $\Img$ that lie in $\Ren(\Sc)$. If $\Ren(\Sc)$ is a proper subset of $\Img$, which is generic, then $\Ren_*$ is not essentially surjective. A dissolve between two shots produces an image path leaving $\Ren(\Sc)$, so it is not in the image of any functor of this form; it must be modeled by a separate operation on $\Path(\Img)$ and not by a movement of the world.
\end{example}

\section{Products and the synchronization of tracks}

A film is more than a picture. Sound, picture, and subtitles run in parallel and are required to be synchronized. Categorically this is a limit construction.

Consider two scenario spaces $\Sc_1$ and $\Sc_2$, for instance the visual world and the acoustic world. Moore paths in the product $\Sc_1\times\Sc_2$ are pairs of Moore paths of the same duration.

\begin{proposition}[Synchronization as a fiber product]\label{prop:fiber}
The category $\Path(\Sc_1\times\Sc_2)$ is isomorphic to the fiber product of categories
\[
\Path(\Sc_1)\times_{(\R_{\ge0},+)}\Path(\Sc_2),
\]
taken over the two duration functors. Explicitly, its objects are pairs $(s_1,s_2)$ and its morphisms are pairs of morphisms with equal duration.
\end{proposition}

\begin{proof}
A Moore path $(\tau,(\gamma_1,\gamma_2))$ in the product has components that are Moore paths of the same duration $\tau$, and conversely any pair of Moore paths of equal duration assembles into a path in the product. Composition and identities are computed componentwise, with durations adding in both components simultaneously. The map is bijective on objects and morphisms and preserves composition, hence is an isomorphism of categories.
\end{proof}

This is the mathematical form of the editor's rule that picture and sound may be cut at different places but must sum to the same running time. A film in which the sound track is allowed to lead or lag the picture, as in the L-cut and J-cut, is not a path in the product but a pair of paths whose seams are offset, and this offset is exactly the failure of the pair to be composed of morphisms of the fiber product at each seam. The relaxed structure, in which durations are required to agree only in total, is the subject of Exercise~\ref{ex:lcut}.

\section{Monoidal structure and string diagrams}

The fiber product is a special case of monoidal composition. A \emph{monoidal category} has, in addition to composition of morphisms in sequence, a tensor product that composes them in parallel, subject to coherence conditions. Pictures called \emph{string diagrams} represent such compositions: sequential composition runs downward, parallel composition runs across~\cite{selinger2011}.

\begin{figure}[ht]
\centering
\begin{tikzpicture}[
  box/.style={draw, rounded corners=2pt, minimum width=1.5cm, minimum height=0.7cm, font=\small},
  wire/.style={thick}
]
\node[box] (p1) at (0,0)   {$\gamma_1$};
\node[box] (s1) at (3,0)   {$\sigma_1$};
\node[box] (p2) at (0,-1.6){$\gamma_2$};
\node[box] (s2) at (3,-1.6){$\sigma_2$};
\draw[wire] (0,1) node[above]{picture} -- (p1);
\draw[wire] (p1) -- (p2);
\draw[wire] (p2) -- (0,-2.6);
\draw[wire] (3,1) node[above]{sound} -- (s1);
\draw[wire] (s1) -- (s2);
\draw[wire] (s2) -- (3,-2.6);
\draw[dashed] (-1.2,-0.8) -- (4.2,-0.8);
\node[right] at (4.3,-0.8) {\small cut};
\end{tikzpicture}
\caption{A two-track film as a string diagram. Vertical composition is the sequence of shots on each track, and horizontal juxtaposition is parallel composition. A synchronized cut crosses both strands at the same level.}
\label{fig:string}
\end{figure}

What the monoidal viewpoint adds is a way of stating when tracks are independent. Two edits on different tracks are independent when the film's change is a tensor $g\otimes g'$ of a change $g$ on one track and a change $g'$ on the other. Composition then satisfies the interchange law
\[
(g\circ f)\otimes(g'\circ f')=(g\otimes g')\circ(f\otimes f').
\]
When this holds, it does not matter whether one thinks of the film as first doing both first-shots and then both second-shots, or as doing each track's two shots in turn. The interchange law itself is an axiom of every monoidal category, so it cannot fail. What can fail is that a given edit is of the form $g\otimes g'$. A sound effect whose meaning depends on the exact image it overlays is a change of the pair that is not a tensor of changes on the tracks, and this is a coupling between them. (In a premonoidal category, where only one track can advance at a time, interchange itself may fail; we do not need that refinement.)

\section{Kuleshov, revisited}

The superadditive interpretation $\mu$ of Chapter~\ref{ch:editing} keeps its form: it is a lax monoidal functor from the word monoid $\Sigma^*$, with concatenation as tensor, to a monoidal preorder of meanings, with comparison $\mu(u)\otimes\mu(v)\le\mu(uv)$ going from parts to whole. A lax monoidal functor need not be strong, and the cases in which it is not are the cases that film theory calls montage. Whether a comparable structure exists on $\Path(\Sc)$ depends on choosing a monoidal structure there, which we do not do; the claim of this section concerns concatenation of shots.

\section{Bridge: three kinds of composition}

Three different operations are all called composition in editing, and the forged-letter scene exhibits each. \emph{Composing tracks} places dialogue over a picture. This is the fiber product of \Cref{prop:fiber}: the two tracks have equal duration and run in parallel, and nothing is added to the running time. \emph{Composing events} places one event after another in the history. Consider a reaction shot of $B$ inserted at $t=62$. If the shot is inserted into the timeline as a new segment of four seconds, the monoid homomorphism $\dur$ adds four seconds and every later event is delayed. If the same four seconds of picture replace the picture under existing dialogue, the sound track continues unchanged and the running time is the same. Insertion and overlay are different operations, with different effects on $\dur$, and a film description that does not say which was used does not determine the running time. \emph{Composing an interpretation} concerns what the reaction shot does to the meaning of the reveal that precedes it. This is the lax structure of the previous sections, and it is not an operation on picture or sound at all.

The rendering functor $\Ren_*$ preserves the first two by definition: it preserves the joining of takes and the durations of paths. It does not preserve the third. Music that continues across a cut, which is widely understood to smooth the perceived transition, is a relation between the sound track and the seam in the picture, and neither the scenario path nor the rendering contains it as a morphism. Likewise the inference that $B$ is shaken by the reveal is not a morphism in $\Path(\Img)$. Whatever relates a seam to its perceived continuity or a shot to its expressive meaning must be supplied by an interpretation functor into a category of meanings, and, as argued above, such a functor is in general only lax.

\section*{Exercises}

\begin{exercise}
Show that the duration map $\dur:\Path(\Sc)\to(\R_{\ge0},+)$ is a functor and that it is faithful if and only if for each pair of points $s,s'$ there is at most one Moore path of each duration from $s$ to $s'$. Under what conditions on $\Sc$ does this fail?
\end{exercise}

\begin{exercise}\label{ex:lcut}
Define a category of two-track films in which the durations of the two tracks need agree only in total, and describe its relation to the fiber product of \Cref{prop:fiber}. Which edits are morphisms of the fiber product and which require the larger category?
\end{exercise}

\begin{exercise}
In the setting of \Cref{prop:rendering-functor}, let $\Ren':\Sc\to\Img$ be a second rendering with $\Ren'=\Ren$ on a subspace $\Sc_0\subset\Sc$. Show that $\Ren_*$ and $\Ren'_*$ agree on the full subcategory $\Path(\Sc_0)$.
\end{exercise}

\chapter{Scripts, Models, and Generation}
\label{ch:scripts}

A script specifies a film without being one. Many films satisfy a given script, and a given film satisfies many scripts. This chapter formalizes that two-sided relation as a Galois connection, which gives the right setting for asking what it means to generate a film from a description.

\section{Satisfaction and the Galois connection}

Let $F$ be a set of films and $\Phi$ a set of constraints, such as ``the protagonist is in the kitchen at the end of the second scene'' or ``no shot exceeds forty seconds''. Let $\models\ \subseteq F\times\Phi$ be a satisfaction relation. For subsets $Y\subseteq F$ and $\Psi\subseteq\Phi$ define
\[
\Th(Y)=\{\psi\in\Phi:\ y\models\psi\ \text{for all }y\in Y\},\qquad
\Mod(\Psi)=\{x\in F:\ x\models\psi\ \text{for all }\psi\in\Psi\}.
\]
A script is a set $\Psi$ of constraints, and $\Mod(\Psi)$ is the set of films that realize it. A set of films has a theory $\Th(Y)$, the set of constraints they all satisfy.

\begin{proposition}[Galois connection]\label{prop:galois}
For all $Y\subseteq F$ and $\Psi\subseteq\Phi$,
\[
Y\subseteq\Mod(\Psi)\iff\Psi\subseteq\Th(Y).
\]
Consequently $\Th$ and $\Mod$ are order-reversing, $Y\subseteq\Mod\Th(Y)$, $\Psi\subseteq\Th\Mod(\Psi)$, and $\Th\Mod\Th=\Th$, $\Mod\Th\Mod=\Mod$.
\end{proposition}

\begin{proof}
Both sides of the equivalence say that $y\models\psi$ for every $y\in Y$ and every $\psi\in\Psi$. Taking $\Psi=\Th(Y)$ in the equivalence gives $Y\subseteq\Mod\Th(Y)$, and symmetrically for the other inclusion. Order reversal is immediate from the definitions. For the idempotence, applying the order-reversing $\Th$ to $Y\subseteq\Mod\Th(Y)$ gives $\Th\Mod\Th(Y)\subseteq\Th(Y)$, while the inclusion $\Th(Y)\subseteq\Th\Mod\Th(Y)$ is the second inclusion applied to $\Psi=\Th(Y)$. The case of $\Mod\Th\Mod$ is the same argument with the roles reversed.
\end{proof}

The closed sets of the connection are the \emph{definable} families of films: those that are exactly the set of realizations of some script. The operator $\Cl=\Th\Mod$ takes a script to its \emph{closure}, the set of all constraints implied by it relative to $F$. A script is \emph{complete} if $\Mod(\Psi)$ is a single film, and \emph{underdetermined} if $\Mod(\Psi)$ contains many films that differ in ways that matter to the director.

\section{Generation as selection from a model class}

In this language, the shooting of a film is the selection of an element of $\Mod(\Psi)$, and a generative system is a mechanism that makes this selection.

\begin{definition}[Generative model for a script]
A \emph{generative model} for a script $\Psi$ is a probability measure $P$ on $F$ with $P(\Mod(\Psi))=1$. It is \emph{faithful} if its support is exactly $\Mod(\Psi)$.
\end{definition}

A system trained on a corpus of films without being told $\Psi$ gives a measure on $F$ that typically assigns positive probability to films outside $\Mod(\Psi)$. This is the formal form of an everyday difficulty: such a system may produce a film in which a character who died in the second scene reappears in the fourth, because nothing in its training objective forbids that. The remedy, developed in Chapter~\ref{ch:admissibility}, is to condition on the script, which replaces $P$ by its restriction to $\Mod(\Psi)$, renormalized, assuming $P(\Mod(\Psi))>0$ (and, for continuous $F$, that $\Mod(\Psi)$ is measurable).

\begin{remark}
The connection runs in the opposite direction as well. Given a corpus $Y$ of films, $\Th(Y)$ is the set of constraints all of them satisfy, which is an inferred house style or genre. Learning a genre from examples is the computation of $\Th(Y)$, or of a statistical relaxation of it, and the result is a script against which new films can be checked.
\end{remark}

\section{Compression as adjunction}

The Galois connection above is the poset-level shadow of an adjunction between categories, and it is worth stating the more general principle. Let $\mathbf{Film}$ be a category whose objects are films and whose morphisms record that one film is an edit or elaboration of another, and let $\mathbf{Script}$ be a category of descriptions ordered by refinement. A \emph{description functor} $D:\mathbf{Film}\to\mathbf{Script}$ and a \emph{realization functor} $G:\mathbf{Script}\to\mathbf{Film}$ form an adjunction $D\dashv G$ if there is a natural bijection
\[
\Hom_{\mathbf{Script}}(D(f),s)\cong\Hom_{\mathbf{Film}}(f,G(s)).
\]
Reading this: a film $f$ is a faithful elaboration of the realization of $s$ exactly when the description of $f$ refines to $s$ in the script category. When the categories are posets, this specializes to a monotone Galois connection, in which the composite $GD$ is a closure operator and $DG$ an interior (kernel) operator. The antitone connection of \Cref{prop:galois} is the same structure with one order reversed.

The principle behind the ``text as substrate'' approach to generative media is that compression, the passage from a film to a description, and reconstruction, the passage back, can be arranged as an adjoint pair, so that the reconstruction is the best approximation consistent with the description. Nothing the description does record is lost exactly when the counit $DG\to\mathrm{id}$ is an isomorphism, that is, when $G$ is fully faithful, which is an additional hypothesis. Whether any practical system realizes such a pair exactly is an empirical question, and the framework serves to state what exactness would mean.

\section*{Exercises}

\begin{exercise}
Let $F=\{0,1\}^*$ be finite binary words regarded as toy films and let $\Phi$ consist of the constraints ``has length at most $n$'' for $n\in\N$ together with ``begins with $0$''. Compute $\Th(\{01,011\})$ and $\Mod(\{\text{begins with }0\})$.
\end{exercise}

\begin{exercise}
Show that $\Mod(\Psi_1\cup\Psi_2)=\Mod(\Psi_1)\cap\Mod(\Psi_2)$, so that adding constraints to a script can only shrink its set of realizations. Show that dually $\Th(Y_1\cup Y_2)=\Th(Y_1)\cap\Th(Y_2)$, and give an example in which $\Mod(\Psi_1)\cup\Mod(\Psi_2)$ is strictly smaller than $\Mod(\Psi_1\cap\Psi_2)$.
\end{exercise}

\begin{exercise}
A script is \emph{unrealizable} if $\Mod(\Psi)=\emptyset$. Give a script for a film in which two constraints are individually realizable but jointly not. How does this relate to the sheaf-theoretic notion of a compatible family with no global section, to be introduced in Chapter~\ref{ch:sheaves}?
\end{exercise}


\part{Sheaves and the Mathematics of Continuity}
\chapter{Presheaves, Sheaves, and Continuity}
\label{ch:sheaves}

The first five chapters treated a film as a global object, a path or a word, and asked how to compose. The mathematics of this chapter starts from the opposite end: local pieces of footage, each defined on a limited stretch of time or region of space, and the question of when they assemble into one consistent whole.

\section{Presheaves and sheaves}

Let $X$ be a topological space and let $\Open(X)$ be the poset of its open sets ordered by inclusion.

\begin{definition}[Presheaf, sheaf]
A \emph{presheaf of sets} on $X$ is a functor $F:\Open(X)^{\op}\to\Set$. For $V\subseteq U$ the image of the inclusion is the \emph{restriction} map $F(U)\to F(V)$, written $s\mapsto\restr{s}{V}$. A presheaf is a \emph{sheaf} if for every open set $U$ and every open cover $U=\bigcup_iU_i$ the following hold.
\begin{description}
\item[Locality.] If $s,t\in F(U)$ satisfy $\restr{s}{U_i}=\restr{t}{U_i}$ for all $i$, then $s=t$.
\item[Gluing.] If $s_i\in F(U_i)$ satisfy $\restr{s_i}{U_i\cap U_j}=\restr{s_j}{U_i\cap U_j}$ for all $i,j$, then there exists $s\in F(U)$ with $\restr{s}{U_i}=s_i$ for all $i$.
\end{description}
\end{definition}

Elements of $F(U)$ are called \emph{sections over $U$}. The two axioms say that a section is determined by its local behavior and that locally compatible data can be assembled. Both are the mathematical form of the idea that continuity is a local property that can be checked piece by piece~\cite{maclane1992}.

\begin{proposition}[The sheaf of takes]\label{prop:takes}
Let $T=[0,\tau]$ with its usual topology and let $\Sc$ be a scenario space. The assignment $F(U)=C(U,\Sc)$, continuous maps from $U$ to $\Sc$, with restriction given by restriction of maps, is a sheaf on $T$.
\end{proposition}

\begin{proof}
Restriction of continuous maps is continuous, and restriction is functorial, so $F$ is a presheaf. For locality, two maps that agree on each $U_i$ agree at every point of $U=\bigcup U_i$. For gluing, a compatible family of continuous maps $s_i:U_i\to\Sc$ defines a unique function $s:U\to\Sc$ with $\restr{s}{U_i}=s_i$, and $s$ is continuous because its restriction to each member of an open cover is continuous.
\end{proof}

Sections of this sheaf over an open interval $U\subseteq T$ are the \emph{takes} defined on $U$. The next section explains how an editor's practice maps onto the sheaf axioms.

\section{Handles, overlaps, and the failure of gluing}

An editor typically has more footage than the cut requires. Each shot comes with \emph{handles}: extra frames before the in-point and after the out-point. Handles are what permit a dissolve, a trim, and a re-cut, and they are the overlaps in the sheaf-theoretic picture.

Consider a film cut at time $t_0$. Choose $\epsilon>0$ and let $U_1=[0,t_0+\epsilon)$ and $U_2=(t_0-\epsilon,\tau]$. Let $s_1\in F(U_1)$ be the first shot extended by its out-handle and $s_2\in F(U_2)$ the second shot extended by its in-handle. The overlap is $U_1\cap U_2=(t_0-\epsilon,t_0+\epsilon)$.

\begin{proposition}\label{prop:cut-fails}
The family $(s_1,s_2)$ glues to a continuous global take if and only if $s_1=s_2$ on the overlap. In particular, if the shots depict different world states at some time in the overlap, no continuous global take extends both, and the edit is a cut in the sense of Chapter~\ref{ch:scenario}.
\end{proposition}

\begin{proof}
This is the gluing axiom for the sheaf of \Cref{prop:takes} applied to the cover $\{U_1,U_2\}$. If the family is compatible on $U_1\cap U_2$ it glues to a unique global section, and conversely a global section restricts to a compatible family.
\end{proof}

This reframes the central quantities of editing. A \emph{hard cut} is a compatible-failure of the pair at the overlap. A \emph{match cut} is a pair that is compatible after rendering but not before: the sections of the sheaf of images $G(U)=C(U,\Img)$, defined in the same way, glue, while the sections of $F$ do not. The comparison is natural: rendering induces a morphism of sheaves $F\to G$, $s\mapsto\Ren\circ s$, and a match cut is a compatible family in $G$ that does not lift to a compatible family in $F$.

\begin{remark}
A morphism of sheaves that is surjective on stalks need not be surjective on sections over a given open set. For sheaves of abelian groups the discrepancy is measured by cohomology of the kernel; for sheaves of sets, as here, no such measure is available in general. The phenomenon is related to, but distinct from, the observation of Chapter~\ref{ch:categories} that $\Ren_*$ is not full, which arises mainly because image paths can leave $\Ren(\Sc)$.
\end{remark}

\section{Sheaves of discrete state}

Not every aspect of a story world varies continuously. Whether a door is open, who holds a gun, and how many drinks have been consumed are discrete. The relevant sheaves are \emph{locally constant}.

\begin{definition}
Let $A$ be a set. The \emph{constant sheaf} $\underline{A}$ on $X$ assigns to each open $U$ the set of locally constant functions $U\to A$.
\end{definition}

A global section of $\underline{A}$ on a connected space $X$ is a constant function. On a disconnected space, it is a choice of value on each connected component. If the time of a film is $T=[0,\tau]$ minus a finite set of cut points, then a global section of $\underline{A}$ over that space is a choice of discrete state per scene, which is the form of the discrete continuity bookkeeping done by a script supervisor. The more interesting case is the one in which the underlying space has loops, taken up next.

\section*{Exercises}

\begin{exercise}
Show that the presheaf of bounded continuous functions on $\R$ is not a sheaf. Which axiom fails?
\end{exercise}

\begin{exercise}
Let $F(U)$ be the set of constant functions on $U$ (constant on all of $U$, not locally constant). Show that $F$ is a presheaf but not a sheaf on a disconnected space.
\end{exercise}

\begin{exercise}
Describe the sheaf $F$ of \Cref{prop:takes} when the time domain is a circle $T=\R/\tau\Z$ representing a looped film. What is a global section? Is it automatic that the take at time $0$ equals the take at time $\tau$?
\end{exercise}

\chapter{Continuity as Cohomology}
\label{ch:cohomology}

When local data fail to assemble, one wants not merely to record the failure but to measure it. Cohomology does this: it assigns to each cover and each sheaf of abelian groups a sequence of groups whose nonzero elements are obstructions to global consistency. This chapter defines Cech cohomology in the lowest degrees, computes two examples that bear directly on film, and states the general principle.

\section{Cech cochains in low degree}

Let $A$ be a sheaf of abelian groups on $X$ and let $\mathcal{U}=(U_i)_{i\in I}$ be an open cover with $I$ totally ordered. Write $U_{ij}=U_i\cap U_j$ and $U_{ijk}=U_i\cap U_j\cap U_k$.

\begin{definition}[Cech complex]
Define
\begin{align*}
C^0(\mathcal{U};A)&=\prod_iA(U_i),\\
C^1(\mathcal{U};A)&=\prod_{i<j}A(U_{ij}),\\
C^2(\mathcal{U};A)&=\prod_{i<j<k}A(U_{ijk}),
\end{align*}
with differentials $d^0:C^0\to C^1$ and $d^1:C^1\to C^2$ given by
\[
(d^0a)_{ij}=\restr{a_j}{U_{ij}}-\restr{a_i}{U_{ij}},\qquad
(d^1c)_{ijk}=\restr{c_{jk}}{U_{ijk}}-\restr{c_{ik}}{U_{ijk}}+\restr{c_{ij}}{U_{ijk}}.
\]
Then $d^1d^0=0$, and the \emph{first Cech cohomology} of the cover is $H^1(\mathcal{U};A)=\ker d^1/\operatorname{im}d^0$.
\end{definition}

The meaning is as follows. A 0-cochain is a choice of local datum on each piece. A 1-cochain is a choice of datum on each pairwise overlap, which one thinks of as a proposed correction relating the local pieces. A 1-cochain is a \emph{cocycle} if its corrections are consistent on triple overlaps, and it is a \emph{coboundary} if it arises as the difference of some choice of local data. The group $H^1$ therefore measures \emph{consistent systems of corrections that are not explained by any choice of local data}.

When the cover is sufficiently fine, so that the sheaf has no higher cohomology on any finite intersection, the groups $H^p(\mathcal{U};A)$ agree with the sheaf cohomology $H^p(X;A)$ of the space, independently of the cover; this is Leray's theorem~\cite{bott1982}. We use it only as a guide to interpretation. All computations below are made directly in the Cech complex of a stated cover.

\section{The time loop}

A looping film, in which the end of the film returns to its beginning, has as time domain a circle $X=\R/T\Z$. Consider a viewer maintaining a \emph{world clock}: a quantity $w$ that increases at unit rate with film time, locally well defined but not obviously globally so. Locally one may take $w(t)=t$ with respect to a chart, and the question is whether a single-valued global $w$ with $w'=1$ exists.

Cover $X$ by the two arcs $U_1=X\setminus\{0\}$ and $U_2=X\setminus\{T/2\}$, where $0$ and $T/2$ denote the classes of those numbers modulo $T$. The intersection $U_{12}=X\setminus\{0,T/2\}$ has two connected components $B=(0,T/2)$ and $B'=(T/2,T)$. Let $A=\underline{\R}$ be the constant sheaf, whose sections over an open set are the locally constant real functions.

Define local world clocks $w_1:U_1\to\R$ by $w_1(t)=t$ for the representative $t\in(0,T)$, and $w_2:U_2\to\R$ by $w_2(t)=t$ for the representative $t\in(-T/2,T/2)$. On $B$ both representatives agree and $w_2-w_1=0$. On $B'$ the class of $t$ has representative $t\in(T/2,T)$ for $w_1$ and $t-T\in(-T/2,0)$ for $w_2$, so $w_2-w_1=-T$. The difference $c=w_2-w_1\in C^1(\mathcal{U};A)=A(U_{12})$ is the locally constant function that equals $0$ on $B$ and $-T$ on $B'$.

\begin{proposition}\label{prop:loop}
The cochain $c$ is a cocycle, and its class in $H^1(\mathcal{U};\underline{\R})\cong\R$ is $-T$. In particular no global single-valued world clock with unit rate exists on a loop of period $T>0$.
\end{proposition}

\begin{proof}
Triple intersections are covered by only two sets, so $C^2=0$ and every 1-cochain is a cocycle. A coboundary has the form $(d^0a)_{12}=\restr{a_2}{U_{12}}-\restr{a_1}{U_{12}}$ with $a_i\in A(U_i)=\R$, since $U_1$ and $U_2$ are connected and locally constant functions on them are constants. Hence coboundaries are exactly the cochains that take the same value on $B$ and on $B'$, that is, the diagonal $\{(r,r)\}\subset\R^2=A(U_{12})$. The quotient $\R^2/\text{diagonal}$ is isomorphic to $\R$ via $(r_B,r_{B'})\mapsto r_{B'}-r_B$, and the class of $c$ maps to $-T-0=-T\neq0$. A global unit-rate world clock would give $w$ on $X$ whose restrictions $w|_{U_1}$ and $w|_{U_2}$ differ from $w_1,w_2$ by constants $a_1,a_2$, forcing $c=d^0(a)$, which contradicts the class being nonzero.
\end{proof}

Cinematically, the loop gives the following statement. In a film whose structure is cyclic, the viewer's local sense of progress cannot be integrated into a single consistent time. Something must be given up: either the clock is allowed to jump at the seam, or the world state accumulates a record, so that the state at the end differs from the state at the start although the images coincide. The second option is the one taken by films in which a day repeats but a character retains memory of earlier repetitions, and it is analyzed in Chapter~\ref{ch:layers}, where the record is placed in the \emph{history} while the image is placed in the \emph{appearance}.

\section{Impossible figures}

Penrose observed that the local geometry of certain impossible figures is consistent while the global geometry is not, and he recorded the obstruction as a cohomology class~\cite{penrose1992}. The same computation applies to any film that shows an impossible space by moving the camera around a closed route in which each local view is geometrically sound.

\begin{proposition}[Tribar obstruction]\label{prop:tribar}
Let $X$ be a circle covered by three arcs $U_1,U_2,U_3$ whose pairwise intersections are nonempty and connected and whose triple intersection is empty. Let $A=\underline{\R}$. Then $H^1(\mathcal{U};A)\cong\R$, with the isomorphism given by $c\mapsto c_{12}+c_{23}-c_{13}$.
\end{proposition}

\begin{proof}
Since $U_{123}=\emptyset$ and $A(\emptyset)=0$, $C^2=0$ and every 1-cochain is a cocycle. Each $U_{ij}$ is connected so $A(U_{ij})=\R$ and $C^1=\R^3$ with coordinates $(c_{12},c_{13},c_{23})$. Each $U_i$ is connected, $C^0=\R^3$, and $d^0(a_1,a_2,a_3)=(a_2-a_1,\ a_3-a_1,\ a_3-a_2)$. The image of $d^0$ is the plane $\{c_{12}+c_{23}-c_{13}=0\}$, since the three coordinates satisfy that identity for every choice of $a$, and the image has dimension $2$ because the kernel of $d^0$ is the diagonal. The functional $c\mapsto c_{12}+c_{23}-c_{13}$ therefore vanishes exactly on the coboundaries and is surjective, giving $H^1\cong\R^3/\R^2\cong\R$.
\end{proof}

Interpret $c_{ij}$ as the depth offset required to reconcile the local depth assignments on the $i$th and $j$th arcs. Each pairwise offset can be absorbed locally, but their alternating sum around the loop is an invariant, and the figure is impossible exactly when it is nonzero. A film that moves the camera along the loop of a Penrose staircase and cuts so that each local view is realizable has the same structure.

\section{The general principle}

The two examples share a pattern, which can be stated as a principle of method.

\begin{principle}[Obstruction principle]\label{pr:obstruction}
Let a film be modeled by local sections of a sheaf over a time or space domain, and let the sections be compatible up to corrections drawn from a sheaf of abelian groups. The film is globally consistent if and only if the associated cohomology class vanishes. The class is invariant under any change of the local data that alters the corrections by coboundaries, which is the sheaf-theoretic form of the statement that re-editing local details cannot repair a global contradiction.
\end{principle}

This is a methodological hypothesis. Its mathematical content, the equivalence of global consistency with vanishing class, is a theorem once the corrections are placed in an abelian sheaf as above and the cover is fixed (it is the classification of torsors by \v{C}ech $H^1$), but its application to a particular film depends on whether the film's continuity constraints can be so modeled. Where they can, the principle yields a precise statement of what no amount of local fixing can achieve.

\section*{Exercises}

\begin{exercise}
Repeat the loop computation with the cover $U_1=X\setminus\{0\}$, $U_2=X\setminus\{T/3\}$, and verify that the class is again $\pm T$.
\end{exercise}

\begin{exercise}
Compute $H^1(\mathcal{U};\underline{\Z})$ for the two-arc cover of the circle and show that it is isomorphic to $\Z$. What does the generator represent for a looped film with a discrete state variable such as a counter?
\end{exercise}

\begin{exercise}
Let $X$ be the figure-eight, two circles joined at a point. Choose a cover and compute $H^1(\mathcal{U};\underline{\R})$. Interpret the answer for a film whose camera path traverses two independent loops in a house.
\end{exercise}

\begin{exercise}
Show that $H^1(\mathcal{U};A)=0$ for any sheaf $A$ whenever $\mathcal{U}$ consists of a single open set equal to $X$. What does this say about a film given as one uninterrupted take?
\end{exercise}

\chapter{Coverage, Descent, and Multiple Views}
\label{ch:descent}

Chapters~\ref{ch:sheaves} and~\ref{ch:cohomology} took time as the base space. Production also works with a cover of the \emph{scene} by many cameras. This chapter treats the second kind of cover, relates it to the problem of contextuality studied in quantum foundations, and states a result on when a family of local reconstructions determines a global one.

\section{Multi-camera coverage}

A scene is staged in a region $X\subseteq\R^3$ over a time interval. Cameras $i\in I$ each see an open subregion $U_i\subseteq X$, their \emph{field of view}, and each camera's footage supports a local reconstruction of the scene state on $U_i$. Let $F$ be the sheaf of scene states on $X$, with $F(U)$ the set of consistent descriptions of the state of the region $U$, and suppose each camera produces a section $s_i\in F(U_i)$.

\begin{definition}[Coverage]
The family of cameras provides \emph{coverage} if $\bigcup_iU_i=X$. It provides \emph{consistent coverage} if it provides coverage and $\restr{s_i}{U_{ij}}=\restr{s_j}{U_{ij}}$ for all $i,j$.
\end{definition}

By the sheaf axioms, consistent coverage determines a unique global state. Production practice contains a method for checking the hypotheses: continuity supervisors compare overlapping views for matching props, costume position, and the like, which is a check of the restriction equalities on $U_{ij}$.

\begin{proposition}\label{prop:cover}
If coverage is consistent, the global scene state $s\in F(X)$ is uniquely determined by $(s_i)$. If coverage is not consistent, there exist indices $i,j$ with $\restr{s_i}{U_{ij}}\ne\restr{s_j}{U_{ij}}$, and no global state restricts to both $s_i$ and $s_j$.
\end{proposition}

\begin{proof}
The first statement is the gluing and locality axioms. For the second, suppose $\restr{s_i}{U_{ij}}\ne\restr{s_j}{U_{ij}}$ in $F(U_{ij})$ for some $i,j$. A global section $s\in F(X)$ restricting to both $s_i$ and $s_j$ would satisfy $\restr{s_i}{U_{ij}}=\restr{(\restr{s}{U_i})}{U_{ij}}=\restr{s}{U_{ij}}=\restr{s_j}{U_{ij}}$ by functoriality of restriction, a contradiction.
\end{proof}

The proposition has a purely practical counterpart: an inconsistency visible in the overlap of two cameras is a certificate that the footage cannot all depict the same event.

\section{Contextuality and empirical models}

The pattern of local consistency without global consistency is not specific to film. It is the structure of contextuality in quantum mechanics, and the sheaf-theoretic account of Abramsky and Brandenburger~\cite{abramsky2011} gives a precise analogue. In that setting, one has a set of measurements, a cover by the maximal sets of jointly performable measurements, and local probability distributions on each such context that agree on overlaps. The empirical model is \emph{non-contextual} if it extends to a global distribution and \emph{contextual} otherwise. Local agreement is a section of a compatible family, and the failure to extend is the failure of gluing.

(The analogy uses a presheaf of local distributions or accounts, which need not be a sheaf; this differs from the sheaf $F$ of scene states above, for which gluing holds by assumption.) The filmic analogue is a collection of viewpoints, each internally coherent and each agreeing with the others where they overlap, with no single world state accounting for all of them. This is a technical rendering of the claim that some narratives are told from perspectives that cannot be reconciled in one objective account. Whether a specific film, such as \emph{Last Year at Marienbad} (1961), has this structure is a matter of reading, and the model allows one to state the claim exactly: that the family of recalled scenes is locally compatible and admits no global section.

\begin{remark}
The analogy with quantum contextuality should be taken only as far as the mathematics goes. The cover and sheaf condition are the common structure. Nothing here implies that films exhibit quantum behavior, and nothing in the physics of contextuality is claimed for narrative.
\end{remark}

\section{Descent}

A more refined statement relates the sheaf condition to a universal property. Given a cover $\mathcal{U}$ of $X$, the sheaf condition says that $F(X)$ is the equalizer of two maps
\[
\prod_iF(U_i)\rightrightarrows\prod_{i,j}F(U_{ij}),
\]
which send a family $(s_i)$ to $(\restr{s_i}{U_{ij}})$ and $(\restr{s_j}{U_{ij}})$ respectively. In the language of descent, a global section is the same as a family of local sections together with the equality of their two restrictions on all overlaps~\cite{maclane1992}.

When the local data are not sets but categories, as when each camera's footage supports a category of possible reconstructions rather than a single one, equality on overlaps is replaced by a specified isomorphism, and these isomorphisms must satisfy a cocycle condition on triple overlaps. The resulting structure is a \emph{descent datum}, and its effectiveness (the existence of a global object inducing it) is a stronger condition than gluing of sets. The cocycle condition is the non-abelian form of the \v{C}ech 1-cocycle condition of the previous chapter. Obstructions to the existence of descent data at all lie one level higher, in second cohomology (gerbes), which we do not pursue. (Strictly, the local data then form a stack, with restriction defined up to canonical isomorphism.)

\begin{definition}[Descent datum for coverage]
A descent datum for cameras $(U_i)$ with local reconstruction categories $F(U_i)$ consists of objects $s_i\in F(U_i)$ and isomorphisms $\phi_{ij}:\restr{s_i}{U_{ij}}\to\restr{s_j}{U_{ij}}$ such that $\restr{\phi_{jk}}{U_{ijk}}\circ\restr{\phi_{ij}}{U_{ijk}}=\restr{\phi_{ik}}{U_{ijk}}$ on triple overlaps. The datum is \emph{effective} if there exists a global $s\in F(X)$ with isomorphisms $\restr{s}{U_i}\cong s_i$ compatible with the $\phi_{ij}$.
\end{definition}

An editor assembling a scene from several cameras supplies the isomorphisms $\phi_{ij}$ implicitly by identifying props and actors across views. The cocycle condition is the requirement that this identification be transitive: if prop $a$ in camera $1$ is prop $b$ in camera $2$ and prop $b$ is prop $c$ in camera $3$, then $a$ is $c$. A violation means the identifications do not form a descent datum, a continuity error of the \emph{identity} kind. Effectiveness is a further property of a datum that does satisfy the cocycle condition, namely that some single global scene accounts for it.

\section{Bridge: when is gluing the right demand?}

Suppose the forged-letter scene is covered by three cameras. Each camera yields local sections of several separately recorded sheaves: the position of objects, the clothing of the characters, the lighting, what each character knows, and the timing of events. Comparison on overlaps proceeds sheaf by sheaf. Camera one shows the letter on the left of the table at $t=30$, and camera two, whose view overlaps camera one's, shows it on the right at $t=38$. Camera three shows $B$'s sleeve rolled up at $t=44$ and rolled down at $t=46$. These readings are at different instants, so they lie on no common overlap in time; they become a disagreement only under the assumption that the letter's position and the sleeve are constant over the interval, and that assumption is what a gluing demand makes.

A mismatch of this kind has several possible explanations, and the mathematics does not choose among them. It may be a production error, of the sort a continuity supervisor exists to catch. It may be a deliberate ambiguity, where the film wishes the audience to be unsure of the scene's physical state. It may be a subjective section, such as a recollection, for which the demand for agreement with another view was never appropriate. Or it may show that the model is inadequate: the letter may have been moved during the pause between $t=31$ and $t=38$, in which case the two position readings, taken before and after the pause, do not conflict once the history contains the event of its being moved, and the apparent contradiction came from asking for a single state where the story has two.

This suggests when gluing is an appropriate demand. The requirement that local sections agree on overlaps is justified for a sheaf that models objective world state, over a stretch of time that the film presents as continuous and objective. It is not justified across a cut to another time, or for a section marked as someone's recollection, which belongs in a separate sheaf indexed by the person who recollects. Only after the sheaf and the stretch have been chosen is a failure to glue evidence of anything. The last explanation above also shows how the apparatus of Chapter~\ref{ch:layers} enters: a section that fails to glue as a state may glue once the history is enlarged by the missing event.

\section*{Exercises}

\begin{exercise}
Let $X=[0,3]$ with $U_1=[0,2)$, $U_2=(1,3]$. Give an example of a sheaf $F$ and local sections $s_1,s_2$ that agree on $U_{12}$ but whose glued section fails to belong to a given subpresheaf $F'\subset F$ of ``allowed'' states. What does this say about checking constraints locally?
\end{exercise}

\begin{exercise}
Verify that the cocycle condition of a descent datum is automatic when the isomorphisms $\phi_{ij}$ arise from a global object, and that it is the only extra condition needed at the triple-overlap level. (Show also that $\phi_{ii}=\mathrm{id}$ follows from the cocycle condition.)
\end{exercise}

\begin{exercise}
Describe a three-camera scene in which each pair of cameras is consistent but the three together are not, in the sense that the transitivity of prop identification fails. Is such a configuration ruled out by the sheaf axiom for sets?
\end{exercise}


\part{Types for Narrative}
\chapter{Dependent Types for Narrative}
\label{ch:types}

Sheaves answer the question of whether local pieces fit together. A different question is whether a piece is \emph{allowed to occur} given what came before. A gun may fire only if it has been introduced; a character may recognize another only if they have met. Such conditions make the legality of an event depend on the history, and the mathematics of dependency is type theory. This chapter introduces the fragment of dependent type theory needed for narrative; Appendix~\ref{app:types} gives the formal rules.

\section{Types, terms, and contexts}

In type theory, a judgment $a:A$ says that $a$ is a term of type $A$. A \emph{context} $\Gamma$ is a list of declarations $x_1:A_1,\,x_2:A_2(x_1),\dots$ in which later types may depend on earlier terms~\cite{martinlof1984}. A judgment $\Gamma\vdash a:A$ says that $a$ has type $A$ in context $\Gamma$.

For narrative, a context is a \emph{story state}: a record of which characters, objects, and facts have been established. A type is a \emph{kind of fact} and a term of that type is a \emph{particular witness} to it. The word witness is used here in its type-theoretic sense, a value that certifies a proposition.

\begin{example}[A minimal story language]\label{ex:chekhov}
Let $\mathsf{Obj}$ be a type of objects, and for each $o:\mathsf{Obj}$ let $\mathsf{Introduced}(o)$ and $\mathsf{Loaded}(o)$ be types of facts about $o$. A story in which a gun is introduced and then fired has the context
\[
g:\mathsf{Obj},\quad i:\mathsf{Introduced}(g),\quad l:\mathsf{Loaded}(g),
\]
and the event $\mathsf{fire}(g)$ is licensed by a rule
\[
\frac{\Gamma\vdash g:\mathsf{Obj}\qquad\Gamma\vdash i:\mathsf{Introduced}(g)\qquad\Gamma\vdash l:\mathsf{Loaded}(g)}{\Gamma\vdash\mathsf{fire}(g,i,l):\mathsf{Fired}(g)}.
\]
The rule demands witnesses of introduction and loading as arguments, so that an attempt to fire a gun that has not been introduced is not an ill-formed story in the sense of poor taste, but an ill-typed one in the sense that the term cannot be constructed.
\end{example}

This is the formal content of what writers call Chekhov's gun, in its permissive direction: what may happen is limited by what has been set up. The converse direction, that what has been set up must be used, is addressed by linear types in the next chapter.

\section{Dependent products and sums}

Two constructors give dependent types their expressive power.

The \emph{dependent sum} $\Sigma_{x:A}B(x)$ is the type of pairs $(a,b)$ with $a:A$ and $b:B(a)$. A character together with a motive is a term of $\Sigma_{c:\mathsf{Character}}\mathsf{Motive}(c)$, where the type of the motive depends on which character it is. The projection to the first component forgets the motive and is the formal version of the observation that a viewer sees the character before understanding the reason for the act.

The \emph{dependent product} $\Pi_{x:A}B(x)$ is the type of functions assigning to each $a:A$ a term of $B(a)$. An obligation is a term of such a type: a function that, for every setup of a given kind, produces a payoff. A promise made to the audience early in a film has the type $\Pi_{s:\mathsf{Setup}}\mathsf{Payoff}(s)$ and is discharged by exhibiting a term.

\begin{definition}[Story]
A \emph{story} over a signature of event rules is a finite sequence $e_1,\dots,e_n$ of events together with contexts $\Gamma_0,\dots,\Gamma_n$ such that $\Gamma_0$ is the initial context and, for each $k$, the rule for $e_k$ is applicable in $\Gamma_{k-1}$ and $\Gamma_k$ is the context obtained by extending $\Gamma_{k-1}$ with the declarations introduced by $e_k$. A story is \emph{well typed} if all such rule applications are derivable.
\end{definition}

\begin{proposition}[Prefix closure]\label{prop:prefix}
Every prefix of a well-typed story is well typed. Hence the set $\Adm$ of well-typed stories is a prefix-closed subset of the free monoid on events.
\end{proposition}

\begin{proof}
The conditions for $e_k$ refer only to $\Gamma_{k-1}$, which depends only on $e_1,\dots,e_{k-1}$. A prefix $e_1,\dots,e_m$ with $m\le n$ therefore uses exactly the first $m$ of the derivations.
\end{proof}

Prefix closure is a small observation with large consequences. It says that admissibility is a property of the past: once a continuation is legal, no later event can retroactively make it illegal. The set $\Adm$ is a tree, and Chapter~\ref{ch:admissibility} studies paths through it.

\section{Refinement and genre}

A \emph{refinement type} $\{x:A\mid\varphi(x)\}$ is the subtype of $A$ cut out by a predicate $\varphi$. Genre conventions are refinements. A whodunit restricts stories to those in which the culprit is introduced before the reveal. A farce restricts to those in which each of a given list of misunderstandings is eventually resolved. Writing a genre as a refinement makes it checkable and, more importantly for generation, makes it a constraint that a sampler can enforce.

The relationship of this structure to the Galois connection of Chapter~\ref{ch:scripts} is direct. The constraints of a script $\Psi$ define a refinement of the type of all films, and $\Mod(\Psi)$ is the type so defined. A well-typed story is thus a story in $\Mod(\Psi)$ for the script whose constraints are the typing rules.

\section{Types as spaces}

The homotopy interpretation of type theory regards a type as a space, a term as a point, and an equality $a=_Ab$ as a path from $a$ to $b$~\cite{hottbook2013}. Two paths between the same endpoints may fail to be equal, in which case the type has nontrivial higher structure. We adopt this reading as a modeling vocabulary, with the same caution as for the quantum analogy of the preceding chapter.

\begin{example}[Rashomon]\label{ex:rashomon}
Let $W$ be a type of world states, $a$ the state in which a crime has occurred and $b$ the state in which it has been adjudicated. Each witness account of the events between $a$ and $b$ is a path $p_k:a=_Wb$. The film \emph{Rashomon} (1950) presents several such paths, and the viewer's question is whether $p_1=p_2$ as paths. If $W$ is a set, in the sense that any two parallel paths are equal, the accounts agree and the question has a trivial answer. If $W$ has nontrivial identity structure, the accounts are distinct and the film is about exactly that distinction. The model thus separates two statements often conflated: that the accounts reach the same end state, which is the statement that $p_1,p_2$ have the same endpoints, and that they are the same account, which is the statement $p_1=p_2$.
\end{example}

\section*{Exercises}

\begin{exercise}
Extend \Cref{ex:chekhov} with a type $\mathsf{Dead}(c)$ for characters and a rule forbidding any event $\mathsf{speak}(c)$ in a context containing a term of $\mathsf{Dead}(c)$. Show that the story in which a character speaks after dying is ill typed. What changes if the character is permitted to speak in flashback?
\end{exercise}

\begin{exercise}
Give a signature for a fair-play detective story: a rule $\mathsf{reveal}$ whose premises include witnesses that every clue used in the solution has been displayed to the audience. Show that prefix closure still holds.
\end{exercise}

\begin{exercise}
Express the Kuleshov effect as a dependency: the type of the viewer's inference from shot $v$ depends on the preceding shot $u$. Write the type of the inference as a dependent product over the type of preceding shots.
\end{exercise}

\chapter{Linear Types and Causal Resources}
\label{ch:linear}

Dependent types say what a past permits. They do not say what a past demands. In a well-made film, a setup is not merely licensed but obligatory to resolve, and the number of times a given setup may be used is limited. These are properties of \emph{resources}, and the type-theoretic apparatus for resources is linear logic~\cite{girard1987}.

\section{Resources in a story}

In ordinary logic a hypothesis may be used any number of times, or not at all. In linear logic, each hypothesis in the linear part of a context must be used \emph{exactly once}. A linear implication $A\multimap B$ is a procedure that consumes an $A$ and produces a $B$. A story fragment ``a loaded gun is fired'' consumes a loaded gun, and what it produces is a fired gun, a wound, an alarm, or whatever the dramatic situation supplies.

The central object is a judgment of the form $\Gamma\mid\Delta\vdash e:B$, where $\Gamma$ is the unrestricted context of facts that persist and $\Delta$ the linear context of resources that must be consumed. A story starts with a linear context $\Delta_0$ of initial resources and ends with a linear context $\Delta_n$ of unconsumed resources. We call a story \emph{balanced} if $\Delta_n=\emptyset$.

\begin{definition}[Resource signature]
A \emph{resource signature} consists of a finite set $\mathsf{R}$ of resource types and a finite set of event rules, each of the form
\[
e:\ \rho_1,\dots,\rho_p\ \multimap\ \sigma_1,\dots,\sigma_q,
\]
with $\rho_i,\sigma_j\in\mathsf{R}$. The event $e$ requires a multiset of $p$ resources and replaces them with a multiset of $q$ resources.
\end{definition}

The state of a story is a finite multiset of resources, that is, a function $\Delta:\mathsf{R}\to\N$. An event rule applies to $\Delta$ if $\Delta\ge$ the multiset of its premises componentwise, and its application replaces premises by conclusions. This is exactly a Petri net: $\mathsf{R}$ is the set of places, events are transitions, and $\Delta$ is the marking.

\section{Balance is decidable and efficient}

\begin{proposition}[Balance check]\label{prop:balance}
Given a resource signature and a finite sequence of events $e_1,\dots,e_n$, one can decide in time $O\bigl(n\cdot m+\lvert\mathsf{R}\rvert\bigr)$, where $m$ is the maximum number of resources in any rule, whether the sequence is a valid run from an initial multiset $\Delta_0$ and whether it ends balanced.
\end{proposition}

\begin{proof}
Maintain the marking $\Delta$ as an array indexed by $\mathsf{R}$. At step $k$, check for each premise of $e_k$ that the count in $\Delta$ is at least its multiplicity in the rule, which takes $O(m)$ time. If the check fails the sequence is not a valid run. Otherwise subtract the premises and add the conclusions, again $O(m)$. After $n$ steps the sequence is a valid run, and it is balanced exactly if every entry of the final array is zero, which is checked by one pass over $\mathsf{R}$ in time independent of $n$.
\end{proof}

The proposition separates a decidable verification problem from the much harder \emph{existence} problem: whether a balanced run exists from a given $\Delta_0$ is the reachability problem for Petri nets, which is decidable but of high complexity, a result due to Mayr and Kosaraju, with later work establishing its precise difficulty. We do not use that result here; the point is that checking a proposed story is cheap, and that finding one is a separate matter to be solved by search.

\begin{example}[Chekhov's gun, linear form]
Take resources $\mathsf{Gun}$, $\mathsf{LoadedGun}$, $\mathsf{Spent}$ and rules
\[
\mathsf{intro}:\ \cdot\multimap\mathsf{Gun},\qquad
\mathsf{load}:\ \mathsf{Gun}\multimap\mathsf{LoadedGun},\qquad
\mathsf{fire}:\ \mathsf{LoadedGun}\multimap\mathsf{Spent}.
\]
The run $\mathsf{intro},\mathsf{load},\mathsf{fire}$ ends with the marking $\{\mathsf{Spent}:1\}$, which is non-empty. The balance requirement is therefore best phrased as ending in a marking contained in a set of \emph{acceptable terminal states}; here $\mathsf{Spent}$ is acceptable and $\mathsf{LoadedGun}$ is not. A story that introduces and loads a gun and never fires it ends with $\{\mathsf{LoadedGun}:1\}$, which is the type-theoretic statement of a gun left hanging on the wall.
\end{example}

The refinement to acceptable terminal markings is typical: the useful requirement is not that every resource be consumed but that the unconsumed ones be of a kind a story may leave behind.

\section{Linearity and the shuffle}

Chapter~\ref{ch:editing} counted interleavings of two storylines. With resources, the count refines to a count of \emph{valid} interleavings. Two storylines that draw on disjoint resources interleave freely, since no event in one affects the applicability of any event in the other. If they share a resource, interleavings that would consume it twice are invalid. This gives a precise form to the editor's rule that intercutting is free only when the lines of action do not interfere.

\begin{proposition}\label{prop:commute}
Let events $e$ and $e'$ have disjoint sets of premise and conclusion resources. Then they commute: for any marking $\Delta$ at which both are applicable in sequence in either order, the two orders reach the same marking.
\end{proposition}

\begin{proof}
Applying $e$ subtracts its premises and adds its conclusions, and likewise for $e'$. Because the two changes act on disjoint coordinates of $\Delta$, the final marking is $\Delta$ plus the sum of the two changes in both orders. Applicability of the second event is unaffected by the first because it depends only on coordinates the first did not touch.
\end{proof}

Disjointness of resource usage is the type-theoretic form of independence, matching the interchange law of Chapter~\ref{ch:categories}. A story in which two events commute is one in which the editor may reorder them, and the conservation of resources is the quantity that determines which reorderings change the story.

\section{Bridge: reusable, consumable, and transforming resources}

A concrete narrative obligation illustrates the three behaviors of resources. A character receives a key, learns a fact, makes a promise, and spends a coin. A \emph{reusable} resource, such as the fact that the letter is forged, can be used any number of times and is copied freely: it belongs to the unrestricted context $\Gamma$ of Chapter~\ref{ch:types}. A \emph{consumable} resource, such as a single-use key or a coin, belongs to the linear context $\Delta$ and disappears when used. A \emph{transforming} resource changes kind when used, as when a key turns a locked door into an open door: the rule consumes $\mathsf{Key}$ and $\mathsf{LockedDoor}$ and produces $\mathsf{OpenDoor}$. The type-theoretic bookkeeping makes explicit which of these a story has in mind, and a promise is a consumable obligation: it is created by one event and must be discharged by another.

Reverse parables are fiction whose developments remain constrained by a principle. Take as principle the one-shot prisoner's dilemma with payoffs $(3,3)$ for mutual cooperation, $(1,1)$ for mutual defection, and $(5,0)$ and $(0,5)$ when one defects against a cooperator. Defection is strictly dominant: against a cooperator it yields $5>3$ and against a defector $1>0$.

\begin{proposition}\label{prop:parable}
In a story whose characters $A$ and $B$ each make one choice, if each choice is the strictly dominant action for the stipulated payoffs, the story ends with mutual defection and payoffs $(1,1)$. Both characters would have received strictly more under mutual cooperation, $(3,3)$.
\end{proposition}

\begin{proof}
Each character chooses defection by dominance, so the outcome is mutual defection with payoff $1$ each, and $3>1$ for each character.
\end{proof}

An unearned exception is a cooperative choice by a character without any earlier event that changes the payoffs. In the typed setting, a cooperative move is admitted only when the context contains a witness of a payoff-changing event, such as a binding promise, a repeated interaction, or a side payment. The check for evasion is thus a check on the typing of the story, and a story that lets its characters cooperate without such a witness is ill typed relative to its own principle. The distinction between satisfying the constraint and communicating it is separate. A story can satisfy the proposition and still fail to let an audience recover the payoff structure from what it shows, which is a question about what the audience has been given, and belongs to the layer of Chapter~\ref{ch:layers} that tracks the audience's knowledge.

\section*{Exercises}

\begin{exercise}
Model a heist story as a resource signature with at least five rules and an acceptable terminal set. Find one balanced run and one run that is valid but ends in an unacceptable marking.
\end{exercise}

\begin{exercise}
Show that if every rule has exactly as many conclusions as premises, then the total number of resources is conserved along every run. Give a filmic example of a story world with a conservation law of this kind.
\end{exercise}

\begin{exercise}
Prove the converse partial result to \Cref{prop:commute}: if two events share a premise resource that has multiplicity one in the marking and neither event produces that resource, at most one of the two orders is applicable.
\end{exercise}


\part{Generation}
\chapter{History, State, and Appearance}
\label{ch:layers}

A system that generates a persistent world must keep track of three different things, which are easy to conflate: what has \emph{happened}, what is the case \emph{now}, and what is \emph{shown}. This chapter separates them, proves that they cannot in general be recovered from one another, and states the design consequence for generative film.

\section{Three layers}

Let $E$ be a set of \emph{events}, the atoms of what can happen in the world: an actor enters, a door opens, a line is spoken. A \emph{history} is a finite word $h\in E^*$, and the set of histories that are admissible, that is, that respect the constraints of the story world, is a prefix-closed set $\Adm\subseteq E^*$ in the sense of \Cref{prop:prefix}.

Three maps organize the layers.

\begin{definition}[Layers]
A \emph{layered world} consists of the following data.
\begin{description}
\item[History.] A prefix-closed set $\Adm\subseteq E^*$ of admissible histories.
\item[State.] A set $Q$ and a function $M:\Adm\to Q$, the \emph{fold}, assigning to each admissible history the current state of the world.
\item[Appearance.] A function $\Ren:Q\to\Img$ assigning to each state the image shown.
\end{description}
\end{definition}

The three layers are thus related by $\Adm\xrightarrow{M}Q\xrightarrow{\Ren}\Img$. The composite $\Ren\circ M$ is what a spectator receives from a history, and the intermediate $M$ is what the world remembers of it. We call $M$ a fold because in practice $M(he)$ is computed from $M(h)$ and $e$ by a transition function, as in a computation that reads events left to right.

The history is the layer that is \emph{authoritative at commit}: when an event is recorded, it is a fact about what happened and is not later reinterpreted. The state is authoritative \emph{between} events and is derived from the history. The appearance is derived from the state and is disposable: it can be recomputed, re-rendered, or replaced by an alternative rendering without changing anything that happened.

\section{State does not determine the future}

The point at issue is whether the state suffices to determine what may happen next. The following result, a version of the Myhill--Nerode theorem~\cite{hopcroft2006}, gives a precise answer.

For $h,h'\in\Adm$ define $h\sim h'$ if for all words $w\in E^*$ we have $hw\in\Adm\iff h'w\in\Adm$. This is an equivalence relation, the \emph{future equivalence}: two histories are equivalent if they permit exactly the same admissible continuations.

\begin{proposition}[Fold sufficiency]\label{prop:nerode}
Let $M:\Adm\to Q$ be a fold. The admissible continuations of a history $h$ are determined by $M(h)$, in the sense that $M(h)=M(h')$ implies $h\sim h'$, if and only if $\ker M\subseteq{\sim}$, where $\ker M=\{(h,h'):M(h)=M(h')\}$. The quotient map $q:\Adm\to\Adm/{\sim}$ is such a fold, and every fold with this property determines $q$: there is a function $g:M(\Adm)\to\Adm/{\sim}$ with $q=g\circ M$. In particular $\Adm/{\sim}$ is the coarsest state space that determines admissible futures.
\end{proposition}

\begin{proof}
The first equivalence is a restatement of the definition of $\ker M\subseteq{\sim}$. The map $q$ has kernel exactly $\sim$, so it determines futures. Now let $M$ determine futures. Define $g(M(h))=q(h)$. This is well defined: if $M(h)=M(h')$ then $h\sim h'$, so $q(h)=q(h')$. By construction $q=g\circ M$, and since $g$ is defined on the image of $M$, the image of $M$ has at least as many elements as $\Adm/{\sim}$. Hence no fold that determines futures can identify more histories than $q$ does.
\end{proof}

The proposition holds for any prefix-closed language. In computational mechanics the classes of $\sim$ are the \emph{causal states}, and the minimal sufficient statistic of the past for the future is the quotient by $\sim$~\cite{shalizi2001}.

Now compare with the appearance. Because $\Ren$ is generally not injective, the composite $\Ren\circ M$ is coarser than $M$, and may be coarser than $\sim$.

\begin{proposition}[Appearance does not determine the future]\label{prop:appearance}
Suppose there exist admissible histories $h,h'$ with $\Ren M(h)=\Ren M(h')$ and $h\not\sim h'$. Then no function of the image of $h$ alone determines the admissible continuations of $h$. Equivalently, if $\Ren\circ M$ determines futures then no such pair $h,h'$ exists; we say the world has a \emph{hidden state} when such a pair does.
\end{proposition}

\begin{proof}
If a function $\Phi$ of the image existed with $\Phi(\Ren M(h))=\{w:hw\in\Adm\}$, then $h$ and $h'$ would have the same image and hence the same set of continuations, which says $h\sim h'$ and contradicts the hypothesis. The equivalent form is the contrapositive of the same argument.
\end{proof}

\begin{example}[Identical images, different pasts]\label{ex:groundhog}
A film in which a day repeats while one character retains memory of earlier repetitions, as in \emph{Groundhog Day} (1993), has appearance that returns to a given image at the start of each day. The histories ending at the start of day one and at the start of day ten share an image but are not future-equivalent, since after day ten the character possesses knowledge that makes different continuations admissible. Here the loop of Chapter~\ref{ch:cohomology} is repaired by the history: the world clock obstruction disappears once the state space is enlarged to $Q=\R/T\Z\times K$ with $K$ a record of accumulated knowledge, because the map $\Adm\to Q$ no longer has to be periodic.
\end{example}

\section{One history, several folds}

The preceding results suggest a design in which \emph{one} history supports \emph{several} folds, each computing a different projection for a different purpose. A \emph{physical fold} $M^{\mathrm{phys}}$ tracks where objects are and how they move. A \emph{semantic fold} $M^{\mathrm{sem}}$ tracks what has been established for the audience: which facts are known, which promises are open. A \emph{provenance fold} tracks which events were attested by observation and which merely inferred. All are computed from $h\in\Adm$, and each is sufficient for its own purposes without being sufficient for the others.

The appearance is then a function of the folds, $x=\Ren(M^{\mathrm{phys}}(h),M^{\mathrm{sem}}(h),\omega)$, where $\omega$ is a source of randomness for rendering. Replaying the same history with the same $\omega$ gives the same frames. Replaying with different $\omega$ gives different frames that depict the same events. This is the property that makes the layered design useful: appearance may be regenerated, but the history is what the world is committed to.

\begin{principle}[Authority principle]\label{pr:authority}
In a persistent generative world, the history is authoritative at the moment of commitment, the state is authoritative between commitments, and the appearance is never authoritative. A system violating this order, for instance by deriving state from rendered frames, will accumulate inconsistencies that cannot be repaired by improving the renderer.
\end{principle}

The principle is a design hypothesis consistent with the failure modes reported for systems that predict frames from frames without an explicit state: such systems tend to forget off-screen objects and to lose consistency when the camera returns to a previously seen place~\cite{bruce2024genie,valevski2024gamengen}. Whether those failures are fully explained by the absence of a history layer is an empirical question. The model states precisely what a remedy would have to supply.

\section{Distributed worlds}

When several processes, such as many players or many rendering nodes, contribute events, there is no single global order. The result of \Cref{prop:nerode} extends if the history is taken to be a partially ordered set of events whose linear extensions are all admissible and give the same state, which is the condition of \emph{commutation} for independent events established in \Cref{prop:commute}. The relevant machinery is the theory of consistent cuts in distributed systems~\cite{lamport1978,chandy1985}: a snapshot of a world is meaningful only relative to a cut of the event structure, and two nodes can agree on a state precisely when their cuts have a common refinement.

\section*{Exercises}

\begin{exercise}
Let $E=\{a,b\}$ and $\Adm=\{\text{words with no two consecutive }b\}$. Compute the future-equivalence classes and the minimal state space. Draw the transition graph.
\end{exercise}

\begin{exercise}
Show that the future equivalence $\sim$ is a right congruence: if $h\sim h'$ then $he\sim h'e$ for every event $e$ with $he\in\Adm$. Conclude that the minimal state space carries a well-defined transition function.
\end{exercise}

\begin{exercise}
In \Cref{ex:groundhog}, give an explicit set $K$ and transition function on $Q=\R/T\Z\times K$ for which the fold is not periodic although the appearance is.
\end{exercise}

\chapter{Admissibility, Policies, and Sampling}
\label{ch:admissibility}

A generative system must produce histories that are admissible and, beyond that, that are good. This chapter studies the two conditions in turn. The first is existence: when is it possible to continue a story indefinitely without becoming inadmissible? The second is distribution: how should a probabilistic model be corrected so that it never leaves the admissible set, and in what sense is the correction optimal?

\section{Infinite admissible histories}

Let $\Adm\subseteq E^*$ be a prefix-closed set of admissible histories. A history $h\in\Adm$ is \emph{extendable} if $he\in\Adm$ for some event $e$. The set $\Adm$ is \emph{non-blocking} if every $h\in\Adm$ is extendable, so that a story is never forced to stop.

\begin{proposition}[Existence of infinite admissible chains]\label{prop:chains}
Suppose $\Adm$ is nonempty, prefix-closed, and non-blocking. Then there exists an infinite sequence of events $e_1,e_2,\dots$ all of whose finite prefixes lie in $\Adm$.
\end{proposition}

\begin{proof}
Choose $h_0=\varepsilon$, which lies in $\Adm$ because $\Adm$ is nonempty and prefix-closed. Given $h_k\in\Adm$, non-blockingness provides an event $e_{k+1}$ with $h_ke_{k+1}\in\Adm$; set $h_{k+1}=h_ke_{k+1}$. The sequence $(e_k)$ so obtained has every prefix $h_k$ in $\Adm$. The construction uses a choice at each step, that is, dependent choice, and no further assumption.
\end{proof}

The hypotheses are not vacuous. Non-blockingness can fail in a way that is easy to miss: a story with a linear resource that must be consumed, as in Chapter~\ref{ch:linear}, may reach a marking from which the acceptable terminal set cannot be reached. Such a history is a \emph{dead end}. The check for dead ends is a reachability problem for the underlying resource net, and a generative system that samples without lookahead will, with positive probability, enter them.

When the set of events is finite, a stronger conclusion holds under weaker hypotheses.

\begin{proposition}[K\"onig]\label{prop:konig}
Suppose $E$ is finite and $\Adm$ is prefix-closed and contains histories of arbitrarily large length. Then there exists an infinite sequence all of whose prefixes lie in $\Adm$.
\end{proposition}

\begin{proof}
Call $h\in\Adm$ \emph{good} if $\Adm$ contains arbitrarily long extensions of $h$. The empty history is good by hypothesis. If $h$ is good, then, since $E$ is finite and extensions of length greater than $\lvert h\rvert$ begin with $he$ for some $e\in E$, at least one $he\in\Adm$ is good, for otherwise the length of admissible extensions of $h$ would be bounded by the maximum of the bounds for the finitely many $he$. Choosing a good $e_{k+1}$ at each step yields the sequence.
\end{proof}

\section{Policies}

A \emph{policy} is a rule that selects the next event given the history. Formally it is a function $\pi:\mathcal{A}'\to E$ on a prefix-closed subset $\mathcal{A}'\subseteq\Adm$ with $h\cdot\pi(h)\in\mathcal{A}'$ for all $h\in\mathcal{A}'$, or more generally a function to probability distributions on $E$ supported on admissible events.

\begin{definition}[Admissible policy]
A policy $\pi$ is \emph{admissible} if the history generated by iterating $\pi$ from $\varepsilon$ lies in $\Adm$ at every finite stage, for every realization of the randomness.
\end{definition}

\begin{corollary}\label{cor:policy}
An admissible policy exists if and only if $\Adm$ contains a nonempty non-blocking prefix-closed subset. In particular, if $\Adm$ is non-blocking, the policy defined in the proof of \Cref{prop:chains} is admissible.
\end{corollary}

The practical content is that admissibility alone does not make a good generator. A policy that always chooses the first admissible event is admissible and typically dull. What distinguishes a good policy is a distribution over admissible continuations that also reflects style, surprise, and rhythm, and for that the language of probability is needed.

\section{Conditioning as optimal projection}

Let $P$ be a probability measure on a countable set $W$ of candidate continuations, for instance the output distribution of a trained model, and let $A\subseteq W$ be the admissible ones, with $P(A)>0$. The \emph{admissibility projection} of $P$ is the conditional measure
\[
P_A(w)=\frac{P(w)\,\mathbf{1}_A(w)}{P(A)}.
\]
It is realized in practice by rejection sampling: draw from $P$ and discard samples outside $A$.

\begin{theorem}[Optimality of conditioning]\label{thm:projection}
Among all probability measures $Q$ on $W$ with $\supp Q\subseteq A$, the measure $P_A$ is the unique minimizer of $\KL(Q\,\|\,P)$, and
\[
\KL(P_A\,\|\,P)=-\log P(A).
\]
\end{theorem}

\begin{proof}
For $Q$ supported on $A$ with $\KL(Q\|P)<\infty$, write
\[
\KL(Q\|P)=\sum_{w\in A}Q(w)\log\frac{Q(w)}{P(w)}
=\sum_{w\in A}Q(w)\log\frac{Q(w)}{P_A(w)}+\sum_{w\in A}Q(w)\log\frac{P_A(w)}{P(w)}.
\]
The first sum is $\KL(Q\|P_A)\ge0$, with equality if and only if $Q=P_A$, by Gibbs' inequality. In the second sum, $P_A(w)/P(w)=1/P(A)$ for $w\in A$, so it equals $-\log P(A)$ independently of $Q$. Therefore $\KL(Q\|P)=\KL(Q\|P_A)-\log P(A)\ge-\log P(A)$, with equality exactly when $Q=P_A$. Evaluating at $Q=P_A$ gives the stated value.
\end{proof}

The theorem identifies the cost of enforcing admissibility. If the model puts probability $P(A)$ on admissible continuations, enforcement through conditioning incurs a KL divergence of $-\log P(A)$ from the model. Rejection sampling needs on average $1/P(A)$ attempts per accepted sample. A model whose admissible mass is $2^{-20}$ will need on the order of a million attempts per accepted sample, and the cost of enforcing admissibility by rejection scales exactly as the reciprocal of the admissible mass. This is the quantitative argument for building admissibility into training or decoding, as opposed to filtering outputs afterwards.

\begin{remark}[Masking is not conditioning]\label{rem:masking}
Constrained decoding, which masks inadmissible next events at each step and renormalizes, differs in general from whole-sequence conditioning, even when $\Adm$ is prefix-closed and non-blocking. At a prefix $h$, conditioning on eventual admissibility selects the next event $e$ with probability proportional to $P(e\mid h)\,P(\Adm\mid he)$, where $P(\Adm\mid he)$ is the probability that a continuation of $he$ is admissible. Masking omits the factor $P(\Adm\mid he)$ and weights each locally admissible event by $P(e\mid h)$ alone. The two agree exactly when (assuming $P(e\mid h)>0$ for the locally admissible events) at every admissible prefix, this factor takes the same value for all locally admissible next events. A small example shows the difference: let $E=\{0,1\}$, let $P$ be uniform on words of length two, and let the admissible length-two words be $00$, $01$, $10$. This set is prefix-closed and non-blocking, since the prefix $1$ extends by $0$. Conditioning gives each admissible word probability $1/3$, so the first event is $0$ with probability $2/3$. Masking leaves both first events locally admissible and chooses each with probability $1/2$. The difference is the weight $P(\Adm\mid he)$ that lookahead supplies. When $\Adm$ can block, masking can also enter dead ends, which is a further reason lookahead is needed.
\end{remark}

\section{Where learned systems stand}

Current learned systems divide, broadly, into those that maintain an explicit latent state used for planning and control~\cite{ha2018worldmodels,hafner2020dreamer} and those that generate frames conditioned on recent frames and actions with no persistent state beyond a compressed context~\cite{valevski2024gamengen,bruce2024genie}. In the vocabulary of Chapter~\ref{ch:layers}, the first kind has a state layer and the second kind has an appearance layer with a short window of history. Neither, as described in the cited work, maintains the committed history as a separate authoritative object. The mathematical analysis suggests what such a layer would give: a set $\Adm$ against which every generated event can be checked and a fold from which state can be recomputed.

\section{Bridge: from valid continuation to cinematic decision}

Admissibility determines which events may follow, not how they are shown. In the forged-letter scene the reveal at $t=58$ is admissible under many realizations. Suppose the pause before the question can last any whole number of seconds from $0$ to $8$, the reaction shot any whole number of seconds from $0$ to $6$, and the music can begin at the door slam, two seconds after it, or not at all. There are $9\times7\times3=189$ realizations on this grid. All are admissible, they share a single history, and they differ in anticipation, emphasis, and the opportunities they give a viewer to settle. Choosing among them is a decision about timing, framing, performance, and sound that occurs after admissibility has been settled and is not made by it.

Where such a decision belongs in the layered world is a design question. A choice that changes only appearance, for instance the exact length of a reaction shot, can be discarded and regenerated without affecting what the world has committed to. A choice that alters what the audience has been told, or what a character can reasonably do next, changes the fold and must be recorded as an event in the history. The decision rule is thus the one stated in the authority principle: a realization choice enters the history exactly when it affects the semantic fold.

\section*{Exercises}

\begin{exercise}
Give an example of a prefix-closed admissible set that contains arbitrarily long histories but is not non-blocking, and show that \Cref{prop:konig} still produces an infinite chain when $E$ is finite. Does the example contradict \Cref{prop:chains}?
\end{exercise}

\begin{exercise}
Show that conditioning is idempotent: $(P_A)_A=P_A$. Show also that for $A'\subseteq A$ with $P(A')>0$ one has $(P_A)_{A'}=P_{A'}$.
\end{exercise}

\begin{exercise}
Suppose the model assigns admissible mass $P(A)=0.01$. How many samples are needed on average to obtain one admissible sample by rejection? Compute $\KL(P_A\|P)$ in bits.
\end{exercise}

\chapter{Selective Cinema and Multiplexing}
\label{ch:selective}

A film is ordinarily one object shown to many viewers. Several technologies break this arrangement: shutter glasses that pass alternate frames to different eyes, polarized and spectrally filtered displays, and interactive branching. In each, a single screen carries several films at once, and each viewer receives a projection of the whole. This chapter gives a model for such systems and states what limits them.

\section{Frame multiplexing}

Let a display refresh at $r$ frames per second, and index frames by $n\in\N$. Let $V$ be a finite set of \emph{viewer classes}. A \emph{schedule} is a partition of the frame indices
\[
\N=\bigsqcup_{v\in V}N_v,
\]
and a viewer of class $v$ wears a selector, for instance a synchronized shutter, that passes only the frames in $N_v$. If $f:\N\to\Img$ is the displayed sequence, viewer $v$ sees $f|_{N_v}$, which we may reindex as a film $f_v:\N\to\Img$ with $f_v(k)=f(n_k^v)$ for the $k$th element $n^v_k$ of $N_v$.

\begin{proposition}[Rate budget]\label{prop:rate}
If the schedule is periodic with the classes sharing frames equally, each viewer perceives $r/\lvert V\rvert$ frames per second. If viewers require at least $r_{\min}$ frames per second for acceptable motion, then $\lvert V\rvert\le\lfloor r/r_{\min}\rfloor$.
\end{proposition}

\begin{proof}
Over a window of $r$ frames, the $\lvert V\rvert$ classes partition the frames equally, so each receives $r/\lvert V\rvert$. The constraint $r/\lvert V\rvert\ge r_{\min}$ rearranges to the stated bound. More generally, for a non-equal schedule the rates $r_v=\lvert N_v\cap\{1,\dots,r\}\rvert$ sum to $r$, so at most $\lfloor r/r_{\min}\rfloor$ classes can have $r_v\ge r_{\min}$.
\end{proof}

At $r=120$ frames per second and $r_{\min}=24$, up to five classes are admissible in this rough accounting; at $r=240$, up to ten. This is a crude bound, since viewers vary in their tolerance and since temporal coherence of low-rate streams can be improved by the content, but it shows that the multiplexing capacity of a display is a resource that is allocated, in the sense of the previous chapters, and not merely a technical feature.

\section{Joint constraints}

In the simplest case the $f_v$ are independent films. More interesting are schedules in which the films are \emph{complementary}: a single audio track serves two viewers who see different pictures, or the two films must be compatible in the sense that scenes that overlap in time satisfy a joint condition.

Let $\mathcal{F}_v$ be the set of admissible films for class $v$, and let $J\subseteq\prod_v\mathcal{F}_v$ be a set of compatible tuples. A \emph{multiplexed film} is a tuple $(f_v)_{v\in V}\in J$ together with a schedule.

\begin{proposition}\label{prop:sheafV}
Consider the discrete space $V$, a nonempty joint constraint $J$, and the presheaf $\mathcal{G}$ on $V$ with $\mathcal{G}(S)=\prod_{v\in S}\mathcal{F}_v$ for $S\subseteq V$, restriction being projection. Then $\mathcal{G}$ is a sheaf. For a joint constraint $J\subseteq\prod_{v\in V}\mathcal{F}_v$, the assignment $\mathcal{J}(S)=\pi_S(J)$, the set of restrictions to $S$ of jointly admissible tuples, is always a subpresheaf of $\mathcal{G}$. It is a subsheaf if and only if $J=\prod_v\pi_v(J)$, that is, if and only if the constraint decomposes into independent constraints on each viewer.
\end{proposition}

\begin{proof}
Every subset of a discrete space is open, and a cover of $S$ is a family of subsets with union $S$. A compatible family of local elements of $\mathcal{G}$ determines a value in $\mathcal{F}_v$ for each $v\in S$, and compatibility on overlaps says these values do not conflict, so the family glues uniquely to an element of $\prod_{v\in S}\mathcal{F}_v$. Hence $\mathcal{G}$ is a sheaf. Projection of a tuple in $J$ to $S$ and then to $S'\subseteq S$ equals projection directly to $S'$, so $\mathcal{J}$ is closed under restriction and is a subpresheaf, whatever $J$ is.

Suppose $\mathcal{J}$ is a subsheaf. Cover $V$ by the singletons $\{v\}$. The family of elements $f_v\in\pi_v(J)=\mathcal{J}(\{v\})$ is compatible, because distinct singletons do not overlap. Gluing in $\mathcal{J}$ gives an element of $\mathcal{J}(V)=J$ whose $v$th coordinate is $f_v$. Thus every tuple of individually admissible films lies in $J$, which says $J=\prod_v\pi_v(J)$. Conversely, if $J=\prod_v\pi_v(J)$ then $\pi_S(J)=\prod_{v\in S}\pi_v(J)$, and a compatible family of local elements of $\mathcal{J}$ glues in $\mathcal{G}$ to a tuple whose coordinates all lie in the sets $\pi_v(J)$, hence to an element of $\pi_S(J)$.
\end{proof}

The proposition says that non-product constraints are exactly the obstruction to the gluing property for the multiplexed film: local admissibility of each viewer's film does not imply joint admissibility. The restrictions of jointly admissible tuples still behave as a presheaf, but a family of individually admissible films need not come from any jointly admissible tuple. The genuinely multiplexed case, in which two films shown together tell a story only when taken together, is the case with a non-product $J$, and the cost of this coupling is that the films cannot be generated independently. A generative system for such a film must sample from $J$, not from each $\mathcal{F}_v$.

\section{Branching and selection}

A branching film is a tree $\mathcal{T}$ of segments, rooted at the opening, in which each internal node has several children. A viewer's experience is a maximal path from the root, chosen by the viewer or by an external condition. The common prefix of two viewers' paths is a stretch they share. Shared segments are a cover in the sense of Chapter~\ref{ch:sheaves}: the viewers agree on the shared stretch and diverge after it.

\begin{proposition}\label{prop:branching}
Let $\mathcal{T}$ be a finite rooted tree with $L$ leaves. The number of distinct viewer experiences is $L$, and the total running time of the content that must be produced is the sum of the durations of all segments, which is at most $L$ times the longest root-to-leaf duration, with equality only when no segments are shared.
\end{proposition}

\begin{proof}
Each maximal path from the root ends at a leaf, and distinct leaves give distinct paths in a tree, so there are $L$ experiences. The content that must be produced is the set of all segments, and each segment lies on at least one root-to-leaf path. A segment lying on $k$ paths is counted $k$ times in the sum of the path durations and once in the content total, so the content total is at most the sum of the path durations, which is at most $L$ times the longest. Equality throughout requires $k=1$ for every segment.
\end{proof}

The sharing of segments is where branching films save work, and it is also where they are constrained: any shared segment must be admissible for every experience that traverses it, and so must be consistent with every history that leads to it. A segment shared by two branches that diverge earlier must not depend on the facts on which they diverged. This is the type-theoretic requirement of Chapter~\ref{ch:types} read on a tree: a term in the shared segment cannot require a witness that exists only on one branch.

\begin{principle}[Shared-segment principle]\label{pr:shared}
A segment can be shared among experiences if it is well typed in the intersection of the contexts of those experiences; and if it is well typed in each experience's context by the same term, the intersection is the weakest context that suffices. The more branches share a segment, the smaller the context in which it must be well typed, and therefore the weaker the dependencies it can carry.
\end{principle}

\section*{Exercises}

\begin{exercise}
A display at 144 frames per second serves three viewer classes with rates $r_1,r_2,r_3$ and $r_{\min}=36$. Find all integer schedules of rates compatible with the budget of \Cref{prop:rate}.
\end{exercise}

\begin{exercise}
Give an example of a non-product joint constraint $J$ on two viewer classes in which each viewer's film is individually admissible but the pair is not, and explain what a generative system must do to sample from $J$.
\end{exercise}

\begin{exercise}
In a branching film with the binary tree of depth $d$, compute the number of leaves and the total content duration if each segment lasts one minute, and compare with the duration of one experience.
\end{exercise}


\part{Equivalence, Transformation, and Direction}
\chapter{Perceptual Equivalence and Effective Dimension}
\label{ch:equivalence}

A film has many adjustable features: durations, levels, color, timing of disclosures, positions of reaction shots. A viewer distinguishes far fewer variations than a production can make. Two questions follow. When may two versions be treated as the same experience? And how many independently meaningful variations does a scene sustain? Both are easily asked loosely and easily answered wrongly. This chapter states them precisely, proves what the formalism does and does not permit, and marks where evidence about viewers must replace mathematics.

\section{Versions, tasks, and dissimilarity}

A \emph{version} is a film together with the presentation conditions relevant to a comparison: the soundtrack as mixed, the playback environment, the subtitles shown. Let $\mathcal{V}$ be a set of versions. Perceptual comparison is conditioned on a \emph{situation} $\theta$ that bundles a viewer or population, a task such as following the plot or judging tension, and a context. For each $\theta$, suppose an operational procedure yields a \emph{dissimilarity} $d_\theta:\mathcal{V}\times\mathcal{V}\to[0,\infty)$, for instance a calibrated detection rate in a discrimination experiment. Everything in this chapter is relative to a fixed $\theta$, and the subscript is dropped.

Three response variables must not be conflated: whether a change is \emph{detected}, whether it changes \emph{comprehension}, and whether it changes \emph{aesthetic or affective response}. Each defines its own $d$. A version pair can be indistinguishable for detection and yet differ in comprehension, as when a deleted line is never noticed but removes the explanation of a later event.

\begin{definition}[Pseudometric dissimilarity]
The function $d$ is a \emph{pseudometric} if $d(x,x)=0$, $d(x,y)=d(y,x)$, and $d(x,z)\le d(x,y)+d(y,z)$ for all versions. It is a \emph{metric} if additionally $d(x,y)=0$ implies $x=y$.
\end{definition}

Whether empirical dissimilarities satisfy the triangle inequality is itself an empirical matter. The results below that rely on it are therefore conditional, and the text says so where it matters.

\section{Exact equivalence and the quotient}

\begin{proposition}[Quotient by zero dissimilarity]\label{prop:quotient}
If $d$ is a pseudometric on $\mathcal{V}$, then the relation $x\approx_0y\iff d(x,y)=0$ is an equivalence relation, and $d$ descends to a well-defined metric on the quotient $\mathcal{V}/{\approx_0}$.
\end{proposition}

\begin{proof}
Reflexivity and symmetry follow from $d(x,x)=0$ and symmetry of $d$. For transitivity, if $d(x,y)=d(y,z)=0$ then $0\le d(x,z)\le d(x,y)+d(y,z)=0$. To see that $d$ descends, suppose $x\approx_0x'$ and $y\approx_0y'$. Then $d(x,y)\le d(x,x')+d(x',y')+d(y',y)=d(x',y')$, and symmetrically $d(x',y')\le d(x,y)$, so the value depends only on classes. The descended function satisfies the metric axioms, and distinct classes have positive distance by construction.
\end{proof}

This is the only sense in which one may speak of a \emph{perceptual equivalence class}: versions at dissimilarity exactly zero, under a pseudometric. Real viewers never produce exact zero. What they produce are thresholded judgments, treated next.

\section{Thresholded similarity is not transitive}

For $\varepsilon>0$ define $x\approx_\varepsilon y$ iff $d(x,y)\le\varepsilon$, read as \emph{indistinguishable at tolerance $\varepsilon$}.

\begin{proposition}[Failure of transitivity]\label{prop:nontransitive}
For every $\varepsilon>0$ the relation $\approx_\varepsilon$ is reflexive and symmetric. It need not be transitive, even when $d$ is a metric. If $x_0\approx_\varepsilon x_1\approx_\varepsilon\cdots\approx_\varepsilon x_n$ and $d$ is a pseudometric, then
\[
d(x_0,x_n)\le n\varepsilon,
\]
and the bound is attained for some pseudometric spaces.
\end{proposition}

\begin{proof}
Reflexivity and symmetry are inherited from $d$. For non-transitivity take $\mathcal{V}=\R$ with $d(x,y)=\lvert x-y\rvert$ and the points $0$, $\varepsilon$, $2\varepsilon$: then $0\approx_\varepsilon\varepsilon$ and $\varepsilon\approx_\varepsilon2\varepsilon$, but $d(0,2\varepsilon)=2\varepsilon>\varepsilon$. The chain bound follows by induction on $n$ from the triangle inequality: $d(x_0,x_n)\le d(x_0,x_{n-1})+d(x_{n-1},x_n)\le(n-1)\varepsilon+\varepsilon$. In $\R$ the points $x_k=k\varepsilon$ show that equality $d(x_0,x_n)=n\varepsilon$ occurs.
\end{proof}

The phenomenon is old in the psychology of discrimination: sequences of stimuli each indistinguishable from the next may connect visibly different endpoints, and the appropriate order structure for such judgments is that of a semiorder rather than an equivalence relation~\cite{luce1956}. The consequence for editing is practical.

\begin{corollary}[Accumulated drift]\label{cor:drift}
Let $T_1,\dots,T_n$ be transformations each satisfying $d(x,T_k(x))\le\varepsilon$ for all $x$ in its domain. Then $d(x,T_n\circ\cdots\circ T_1(x))\le n\varepsilon$ for a pseudometric $d$. To guarantee a final drift of at most $\delta$ after $n$ steps by this argument alone, each step must satisfy $\varepsilon\le\delta/n$.
\end{corollary}

\begin{proof}
Apply \Cref{prop:nontransitive} to the chain $x,\,T_1x,\,T_2T_1x,\dots$, each consecutive pair being within $\varepsilon$ by hypothesis.
\end{proof}

The corollary explains a recurring failure of iterated processing: a pipeline of steps each judged imperceptible can produce a result that is not. It also states what validation must do. Imperceptibility of each step does not certify the composite, and the composite must be compared with the original directly. If the triangle inequality fails for the empirical $d$, the bound itself may fail, and direct comparison is then the only safe course.

\section{Effective degrees of freedom}

Suppose a scene is controlled by $p$ real parameters $\theta_1,\dots,\theta_p$ forming an open set $\Theta\subseteq\R^p$: cut times, gains, durations, gesture timings. Suppose a differentiable \emph{response model} $r:\Theta\to\R^q$ maps parameters to a vector of response measures, such as predicted comprehension scores or ratings of tension, and that perceptual dissimilarity between parameter settings is $d(\theta,\theta')=\lVert r(\theta)-r(\theta')\rVert_W$, where $\lVert z\rVert_W=(z^{\top}Wz)^{1/2}$ for a fixed symmetric positive definite weight matrix $W$ (the Euclidean norm when $W=I$). Perturbation sizes in parameter space are measured in the Euclidean norm. The weights encode the relative importance of the response measures, and replacing $J$ by $W^{1/2}J$ reduces the weighted case to the Euclidean one, so below $J$ denotes the weighted sensitivity matrix $W^{1/2}Dr(\theta)$ when $W\neq I$. Such a model is a hypothesis to be fit and tested, not something the mathematics supplies.

The \emph{sensitivity matrix} at $\theta$ is the Jacobian $J=Dr(\theta)\in\R^{q\times p}$, weighted as just described. Its rank counts the locally independent directions in which parameters change the response.

\begin{proposition}[Redundant controls]\label{prop:redundant}
If $J=Dr(\theta)$ has rank $k$ and the rank of $Dr$ is constant, equal to $k$, on a neighborhood of $\theta$, then there is a neighborhood of $\theta$ in which the level set of $r$ through each point is an embedded submanifold of dimension $p-k$. In particular, there are $p-k$ independent directions of parameter change that leave the modeled response unchanged.
\end{proposition}

\begin{proof}
This is the constant rank theorem~\cite{lee2013}: if $Dr$ has constant rank $k$ near $\theta$, there are local coordinates in which $r$ takes the form $(x_1,\dots,x_p)\mapsto(x_1,\dots,x_k,0,\dots,0)$. The fibers are then the coordinate planes of dimension $p-k$.
\end{proof}

Exact invariance is an idealization. A viewer who detects response changes only above a tolerance $\varepsilon$ treats directions of small sensitivity as redundant as well. Let $\sigma_1\ge\cdots\ge\sigma_{\min(p,q)}$ be the singular values of $J$ and let $v_i$ be unit right singular vectors.

\begin{proposition}[First-order effective dimension]\label{prop:effdim}
With the Euclidean norm on parameters and the (weighted) Euclidean norm on responses, a parameter perturbation $\rho v_i$ of size $\rho>0$ changes the modeled response by $\sigma_i\rho$ to first order, that is, $r(\theta+\rho v_i)-r(\theta)=\rho Jv_i+o(\rho)$ with $\lVert Jv_i\rVert=\sigma_i$. Hence the number of singular directions in which a perturbation budget $\rho$ produces a first-order change exceeding tolerance $\varepsilon$ is
\[
d_{\mathrm{eff}}(\theta;\varepsilon,\rho)=\#\{i:\ \sigma_i>\varepsilon/\rho\}.
\]
\end{proposition}

\begin{proof}
Differentiability gives the first-order expansion. By the definition of the singular value decomposition $Jv_i=\sigma_iu_i$ with $u_i$ a unit vector in the Euclidean norm, so $\lVert Jv_i\rVert_2=\sigma_i$. The argument uses that the norms are induced by inner products, which is why an arbitrary norm is not permitted. The first-order change in direction $v_i$ exceeds $\varepsilon$ exactly when $\sigma_i\rho>\varepsilon$.
\end{proof}

The quantity $d_{\mathrm{eff}}$ is not a property of the film. It depends on the model $r$, the base point $\theta$, the tolerance $\varepsilon$ that encodes a viewer and task, and the budget $\rho$ that encodes how much one is willing to vary. It never exceeds $\min(p,q)$, and it generally decreases as $\varepsilon$ grows or $\rho$ shrinks. The proposition therefore gives a precise form to the question of how many parameters can be varied before a difference is noticed, without asserting that cinema has a single intrinsic dimension. The first-order character matters: when several parameters change at once by large amounts, the linearization may fail, and the interaction of changes must be examined by experiment.

\begin{principle}[Conditional dimension]\label{pr:conditional}
The number of perceptually independent variations a scene sustains is a conditional quantity, indexed by scene, viewer population, task, tolerance, and perturbation budget. Claims of a universal number are unsupported, and estimates should be reported with all five conditions.
\end{principle}

Its status is methodological: it constrains what may be claimed, and it would be contradicted by a demonstration that the conditional estimates were stable across scenes and populations to a degree that justified a single figure.

\section{Worked perturbations in the forged-letter scene}\label{sec:perturb}

The apparatus of the previous section is abstract until a model is written down. This section writes down a deliberately simple one for the forged-letter scene of Section~\ref{sec:scene}, computes every quantity of the preceding propositions, and then shows where the first-order analysis fails. The response model is \emph{stipulated}. Its functional forms are chosen for transparency, not fitted to viewers, and none of the numbers below is a measurement. What the section demonstrates is the method: which units appear, what a singular value means in those units, how the tolerance enters, and what the local analysis can and cannot say.

\paragraph{Parameters and units.} Four controls are varied, with base values taken from the baseline edit.
\begin{center}
\begin{tabular}{llll}
\toprule
symbol & control & unit & base value\\
\midrule
$p_1$ & length of the pause before $B$'s question (cut timing) & s & $7$\\
$p_2$ & delay of music onset after the door slam & s & $2$\\
$p_3$ & music level relative to the reference & dB & $-2$\\
$p_4$ & duration of $B$'s reaction shot (gesture duration) & s & $4$\\
\bottomrule
\end{tabular}
\end{center}
The disclosure order, namely whether the audience learns that $A$ knows before or after $B$'s question, is a fifth control and is discrete. It has no derivative and is treated separately below.

\paragraph{Response model.} Three response measures are modeled, each scaled to $[0,1]$: anticipation $q_1$ (how strongly the pause builds expectation of the question), fit of the music $q_2$, and salience of the reaction $q_3$. We stipulate
\[
q_1=\tanh\!\Bigl(\frac{p_1+\tfrac12p_4}{6}\Bigr),\qquad
q_2=\frac{1}{1+e^{-(p_3+5)/4}}\;e^{-(p_2-2)^2/18},\qquad
q_3=1-e^{-p_4/2}.
\]
All three measures are dimensionless and lie in $[0,1]$, so we take $W=I$, and tolerances $\varepsilon$ are in the same units. At the base point $\theta_0=(7,2,-2,4)$ the responses are $r(\theta_0)=(0.9051,\,0.6792,\,0.8647)$. The Jacobian $J=Dr(\theta_0)$, with columns ordered $p_1,\dots,p_4$, is
\[
J=\begin{pmatrix}
0.0301 & 0 & 0 & 0.0151\\
0 & 0 & 0.0545 & 0\\
0 & 0 & 0 & 0.0677
\end{pmatrix},
\]
with entries in response units per second for $p_1,p_2,p_4$ and per decibel for $p_3$. The entries are not comparable as they stand, because a second and a decibel are not commensurable. The remedy is to state a perturbation budget in the units of each control.

\paragraph{A budget with units.} Let $D=\operatorname{diag}(1\,\mathrm{s},\,1\,\mathrm{s},\,2\,\mathrm{dB},\,1\,\mathrm{s})$. The budget is the set of perturbations $\theta_0+\rho Du$ with $\lVert u\rVert\le1$ and $\rho=1$: an editor is willing to move the timing controls by about a second and the level by about two decibels. In the variable $u$ the sensitivity matrix is $JD$, and \Cref{prop:effdim} applies to the matrix $JD$ (in place of $J$) with the Euclidean norm on $u$, and the budget is $\rho=1$. The singular values of $JD$ are
\[
\sigma_1=0.1089,\qquad\sigma_2=0.0697,\qquad\sigma_3=0.0292 .
\]
\Cref{prop:effdim} then gives $d_{\mathrm{eff}}=\#\{i:\sigma_i>\varepsilon/\rho\}$. For a viewer tolerance $\varepsilon=0.05$ it gives $d_{\mathrm{eff}}=2$, and for the stricter $\varepsilon=0.02$ it gives $d_{\mathrm{eff}}=3$. Four controls therefore support two perceptually independent variations for the coarser viewer and three for the finer one, in this model, for this budget. This is the content of \Cref{pr:conditional} in a single example: the number depends on the tolerance, and it would change with a different budget $D$.

Two structural features are visible. The matrix $J$ has rank $3$ with a nontrivial kernel, and $p_1$ and $p_4$ act through a common response $q_1$, so a change in the pause can be partly offset by an opposite change in the reaction duration. Indeed $\partial q_1/\partial p_1=0.0301$ and $\partial q_1/\partial p_4=0.0151$, so lengthening $p_4$ by $1$ s and shortening $p_1$ by $0.5$ s leaves $q_1$ unchanged to first order, although $q_3$ changes. The column for $p_2$ is zero because the music onset delay is at its modeled optimum, where the response is quadratic and its derivative vanishes.

\paragraph{Where the first-order analysis fails.} The zero column suggests that the onset delay is irrelevant. It is not. Shifting $p_2$ by $1.5$ s changes $q_2$ by $-0.0798$, and by $3$ s by $-0.2672$, both beyond $\varepsilon=0.05$, although the first-order prediction is exactly zero. The general statement is a bound on the remainder.

\begin{proposition}[Validity of the first-order prediction]\label{prop:taylor}
Let $r$ be $C^2$ on an open set containing the segment from $\theta$ to $\theta+h$. Let $M_h$ be an upper bound for $\lVert D^2r(\xi)[h,h]\rVert/\lVert h\rVert^2$ over $\xi$ on the segment. Then
\[
\bigl\lVert r(\theta+h)-r(\theta)-Dr(\theta)h\bigr\rVert\le\tfrac12M_h\lVert h\rVert^2 .
\]
Consequently, if $\lVert Dr(\theta)h\rVert>\varepsilon+\tfrac12M_h\lVert h\rVert^2$ the change exceeds $\varepsilon$, and if $\lVert Dr(\theta)h\rVert<\varepsilon-\tfrac12M_h\lVert h\rVert^2$ it does not. Between these values the first-order estimate does not decide.
\end{proposition}

\begin{proof}
Put $\phi(s)=r(\theta+sh)$. By Taylor's theorem with integral remainder, $\phi(1)-\phi(0)-\phi'(0)=\int_0^1(1-s)\,\phi''(s)\,ds$, and $\phi''(s)=D^2r(\theta+sh)[h,h]$. Taking norms, the integrand is at most $(1-s)M_h\lVert h\rVert^2$, and $\int_0^1(1-s)\,ds=\tfrac12$.
\end{proof}

The bound depends on the direction $h$ and on the segment, not on a global bound for the whole map, which is what the worked perturbations need. For a shift of $p_2$ the segment lies in the line $\theta_0+te_2$, along which $q_2$ is the only response that varies, with $\phi''=\partial^2q_2/\partial p_2^2$. Its absolute value is at most $M_h=0.0755$ on $p_2\in[0,3]$, the maximum being attained at the base point (verified numerically in the accompanying script).

With $M_h=0.0755$ along this line, the first-order prediction is decisive only when $\tfrac12M_h\lVert h\rVert^2$ is small compared with $\varepsilon$, which means $\lVert h\rVert\ll\sqrt{2\varepsilon/M_h}=1.15$ s. At $h=1.5$ s the remainder bound is $0.085$, which contains the observed change $0.080$; at $h=3$ s the bound $0.340$ contains the observed $0.267$. For the pause the same bound along the $p_1$ line is $0.162$ at $h=-4$ s against an actual remainder of $0.102$, so the bound is valid but not sharp. The practical rule is that a sensitivity analysis is reported together with the radius within which it was checked.

\paragraph{Asymmetry of cut timing.} The linearization also misses asymmetry. For the pause, $\partial q_1/\partial p_1=0.0301$ per second predicts a symmetric tolerance of $\pm0.05/0.0301=\pm1.66$ s. The model itself gives $q_1\ge0.9051-0.05$ only if $p_1\ge5.65$, so the pause may be shortened by $1.35$ s, and $q_1\le0.9051+0.05$ up to $p_1=9.32$, so it may be lengthened by $2.32$ s. A reduction of $4$ s is predicted by the linearization to cost $0.120$ and costs $0.223$ in the model, because $\tanh$ is concave. Saturating responses are typical of perceptual measures, and for such responses the tolerance region is not symmetric about the base point.

\paragraph{Weighted measures.} If $q_1$ is judged to matter twice as much as the others, $W=\operatorname{diag}(4,1,1)$ and $W^{1/2}=\operatorname{diag}(2,1,1)$. The sensitivity matrix becomes $W^{1/2}JD$, whose first row is doubled, and the singular values must be recomputed from that matrix; the unweighted values do not carry over. The weights are a statement about the task and are part of the conditions that \Cref{pr:conditional} requires one to report.

\paragraph{Disclosure order.} Let $t_d$ be the time at which the audience is first given the fact that $A$ knows the letter is forged. If $t_d<38$ the question at $t=38$ is received with the audience already in possession of the fact, and the scene reads as dramatic irony. If $t_d>38$ the audience shares $B$'s ignorance when the question is asked and the scene reads as a mystery resolved at $t_d$. Although $t_d$ is a number, the change in how the question is read is not expected to vary smoothly across $t_d=38$, and a response model fitted on either side will not extrapolate across it. The honest treatment is to partition the parameter space at the discontinuity, compute sensitivity within each piece, and compare the two pieces directly by experiment. No Jacobian is defined at the boundary, and \Cref{prop:effdim} says nothing about it. This is a concrete instance of the earlier remark that the local analysis describes smooth motion within a class of versions and not a change of class.

\paragraph{What an experiment would supply.} To replace the stipulated model one needs, for each control, a set of stimuli at several levels, paired judgments by a viewer population on a stated task, a fitted $r$, and a check of \Cref{prop:taylor} at the radii used. The singular spectrum so obtained is the quantity the conjecture of the next section is about.

\section{Painting and minimal cues}

A traditional observation in the study of pictures is that a small repertoire of cues, such as balance, directional guidance, repeated forms, and selective emphasis, suffices for a work to be seen as composed. The film analogue would be a repertoire of temporal cues: an anticipatory pause, a reaction shot, a delayed disclosure. In the present vocabulary the analogue is the claim that $d_{\mathrm{eff}}$ is small relative to $p$ near typical production settings, because most parameters are either redundant or masked by a few dominant cues. That is an empirical conjecture. It can be tested by estimating the singular spectrum of $J$ from controlled perturbation experiments, and a rapidly decaying spectrum would support it.

\section*{Exercises}

\begin{exercise}
Let $\mathcal{V}=\{0,1,\dots,N\}$ with $d(x,y)=\lvert x-y\rvert$ and $\varepsilon=1$. For which pairs does $x\approx_1y$ hold? Show that the transitive closure of $\approx_1$ is the full relation, and explain why this illustrates the difficulty of defining equivalence classes from thresholded judgments.
\end{exercise}

\begin{exercise}
Suppose the empirical dissimilarity satisfies only a relaxed triangle inequality $d(x,z)\le K\bigl(d(x,y)+d(y,z)\bigr)$ for some $K>1$. Show that $\approx_0$ is still an equivalence relation, and that for $n=2$ the drift bound becomes $d(x_0,x_2)\le2K\varepsilon$.
\end{exercise}

\begin{exercise}
Let $r(\theta_1,\theta_2,\theta_3)=(\theta_1+\theta_2,\ \theta_3)$. Compute the rank of $Dr$ and describe the level sets. Interpret $\theta_1$ and $\theta_2$ as two timing parameters that jointly determine one perceived delay.
\end{exercise}

\begin{exercise}
Design a controlled perturbation experiment to estimate the singular values of the sensitivity matrix for a two-parameter scene, specifying the response measures, the perturbation sizes, and a rule for deciding when the linearization is unreliable.
\end{exercise}

\begin{exercise}
In Section~\ref{sec:perturb}, change the weight matrix to $W=\operatorname{diag}(4,1,1)$ and recompute the singular values of $W^{1/2}JD$. For which tolerances $\varepsilon$ does $d_{\mathrm{eff}}$ change compared with the unweighted case?
\end{exercise}

\begin{exercise}
For the response $q_1=\tanh((p_1+\tfrac12p_4)/6)$, find the set of $(p_1,p_4)$ with $q_1\ge0.855$ and compare it with the half-plane predicted by linearization at $(7,4)$. Where do the two differ most?
\end{exercise}

\chapter{Transformations and Preservation Contracts}
\label{ch:contracts}

A system that changes a film must say what it will leave alone. A quieter soundtrack, a compressed scene, a genre substitution, and a redirected character do not preserve the same things, and a transformation described only by what it changes cannot be evaluated. This chapter formalizes the promise a transformation makes, shows when promises compose, and works through acoustic occupation as a case in which the formalism yields exact results.

\section{Contracts}

Let $\mathcal{V}$ be a set of versions as in Chapter~\ref{ch:equivalence}. A \emph{preservation relation} is a binary relation $\rho\subseteq\mathcal{V}\times\mathcal{V}$ read as ``$y$ preserves what $\rho$ cares about in $x$''. Examples are preservation of event order, of causal relations, of dialogue intelligibility, of dramatic function, and tolerance relations $\approx_\varepsilon$ from the previous chapter.

\begin{definition}[Transformation contract]
A \emph{transformation contract} consists of a function $T:\mathcal{D}\to\mathcal{V}$ on a domain $\mathcal{D}\subseteq\mathcal{V}$ of versions it applies to, together with preservation relations $\rho_1,\dots,\rho_m$, written $(T,\mathcal{D},\rho_1,\dots,\rho_m)$, such that $x\,\rho_j\,T(x)$ for all $x\in\mathcal{D}$ and all $j$. The \emph{target change} is whatever $T$ alters that is not in the scope of any $\rho_j$.
\end{definition}

A contract is silent about everything outside $\rho_1,\dots,\rho_m$, which is the point: it makes explicit which aspects are guaranteed and, by omission, which are not. Preservation of one aspect does not entail preservation of another, and Section~\ref{sec:conflict} gives a precise example.

\begin{proposition}[Composition of contracts]\label{prop:compose-contracts}
Let $(T,\mathcal{D},\rho)$ and $(S,\mathcal{D}',\rho)$ be contracts with $T(\mathcal{D})\subseteq\mathcal{D}'$.
\begin{enumerate}
\item If $\rho$ is transitive, then $(S\circ T,\mathcal{D},\rho)$ is a contract.
\item If $\rho=\approx_\varepsilon$ for a pseudometric $d$, then $S\circ T$ satisfies $\approx_{2\varepsilon}$.
\end{enumerate}
\end{proposition}

\begin{proof}
For the first claim, $x\,\rho\,T(x)$ and $T(x)\,\rho\,S(T(x))$ give $x\,\rho\,S(T(x))$ by transitivity. For the second, $d(x,ST(x))\le d(x,T(x))+d(T(x),ST(x))\le2\varepsilon$.
\end{proof}

The lesson is the same as \Cref{cor:drift}. Exact preservation relations such as equality of event order are transitive and compose without loss. Thresholded relations do not, and a chain of transformations must be validated end to end.

\section{Conflicting preservation}\label{sec:conflict}

\begin{example}[Compression preserves order but not recovery]
Model a version as a finite sequence of intervals, each tagged \emph{active} or \emph{quiet} and carrying a duration. Let $\rho_{\mathrm{ord}}$ say that two versions have the same sequence of tags, ignoring durations, and let $\rho_{\mathrm{rec}}$ say that every maximal run of quiet intervals in $x$ has a counterpart in $y$ of at least a stated minimum total duration $\tau_0$. A transformation $T$ that shortens all intervals by a common factor $c<1$ preserves $\rho_{\mathrm{ord}}$ exactly, but violates $\rho_{\mathrm{rec}}$ for any quiet run whose scaled duration falls below $\tau_0$. The two relations are independent: a transformation preserving one can destroy the other, and a contract that lists only $\rho_{\mathrm{ord}}$ licenses a compression that removes every opportunity for recovery.
\end{example}

The example is deliberately simple. Its point is the logical form: preservation relations are not ordered by implication in general, and a contract must name each aspect it guarantees. The chapters of this book have introduced several candidates, namely event order and causal relations (Chapters~\ref{ch:editing} and~\ref{ch:types}), continuity (Chapter~\ref{ch:cohomology}), admissibility (Chapter~\ref{ch:admissibility}), and the perceptual tolerance of Chapter~\ref{ch:equivalence}. A transformation of film should state its contract in terms of these.

\section{Acoustic occupation}\label{sec:occupation}

Sound offers a setting where the temporal organization of a feature, as distinct from its distribution, can be treated exactly. Let a soundtrack be reduced to a discrete-time sequence of levels $L_1,\dots,L_n$ for non-dialogue sound and $D_1,\dots,D_n$ for dialogue, in decibels over a fixed window. Fix a margin $\mu\ge0$ and call time step $t$ \emph{exceeding} if $L_t>D_t+\mu$. The \emph{exceedance set} is $E=\{t:L_t>D_t+\mu\}$.

An \emph{occupation block} is a maximal run of consecutive exceeding steps, and a \emph{recovery interval} is a maximal run of consecutive non-exceeding steps lying between two blocks. These definitions are dialogue-referenced: a loud effect under silent dialogue and the same effect under loud dialogue are different with respect to $E$.

\begin{proposition}[Distribution does not determine arrangement]\label{prop:arrangement}
Let $n$ time steps contain $m$ exceeding steps, with $0<m<n$. Then the number $B$ of occupation blocks can be any integer with $1\le B\le\min(m,\,n-m+1)$, and each value occurs for some arrangement. Consequently the histogram of levels, which determines $m$, does not determine $B$ or the lengths of recovery intervals.
\end{proposition}

\begin{proof}
Encode an arrangement as a binary sequence of length $n$ with $m$ ones. The number of blocks is the number of maximal runs of ones. Placing all ones consecutively gives $B=1$. For a given $B$ with $1\le B\le m$, split the $m$ ones into $B$ nonempty groups, which is possible since $B\le m$, and separate consecutive groups by at least one zero, which requires $B-1\le n-m$ zeros, that is, $B\le n-m+1$. Place the remaining zeros anywhere compatible with these separations. No arrangement can have more than $m$ runs, nor more than one run more than the number of zeros, since each pair of adjacent runs is separated by at least one zero. Thus $B$ ranges exactly over $1\le B\le\min(m,n-m+1)$.
\end{proof}

A soundtrack with an isolated peak every ten seconds and a soundtrack with a single thirty-second sustained passage may thus have identical level histograms and identical total exceedance time, and differ completely in the opportunities they give a listener to recover. This is the reason for treating acoustic occupation as a temporal structure and not a statistic.

\section{A selective attenuator and its guarantee}

Consider a proposed class of tools that attenuate selected non-dialogue stems, for instance explosions or music under speech, by a gain factor $\alpha\in(0,1]$ applied to the level $L_t$, in decibels as an offset $-a$ with $a\ge0$, leaving dialogue unchanged.

\begin{proposition}[Monotone reduction of occupation]\label{prop:mufflerizer}
Let $T_a$ reduce the non-dialogue level by $a\ge0$ decibels at all times, so that $L_t^{(a)}=L_t-a$ and $D_t^{(a)}=D_t$. Then the exceedance sets are nested, $E^{(a')}\subseteq E^{(a)}$ for $a\le a'$, the number of exceeding steps is non-increasing in $a$, and dialogue level is unchanged.
\end{proposition}

\begin{proof}
If $t\in E^{(a')}$ then $L_t-a'>D_t+\mu$, and since $a\le a'$ this gives $L_t-a>D_t+\mu$, so $t\in E^{(a)}$. The inclusion implies the count is non-increasing. Dialogue is untouched by construction.
\end{proof}

The proposition is a guarantee about a defined quantity under a defined intervention, and it is exact. It says nothing about whether listeners prefer the result, whether intelligibility improves beyond what the margin captures, or whether the reduced version carries the same dramatic effect. Each of those is a separate preservation relation requiring its own operational definition and evidence. Selective attenuation by source category adds a further risk that the proposition does not address, namely misclassification: if screams or music are labeled as another category, the wrong material is attenuated, and the contract must state the cost of such errors.

\section{Alternative versions of one scene}\label{sec:altversions}

The contracts above become useful when several versions of the same scene are laid beside one another and each is checked against a stated list of relations. Take the forged-letter scene of Section~\ref{sec:scene} as the baseline $V_0$ and fix the following relations, all stipulated for this example.
\begin{itemize}
\item $\rho_{\mathrm{ord}}$: the sequence of history events is unchanged.
\item $\rho_{\mathrm{ant}}$: a quiet interval of at least $\tau_0=5$ s precedes $B$'s question. In $V_0$ the pause lasts $7$ s.
\item $\rho_{\mathrm{dia}}$: dialogue is unchanged in content and level.
\item $\rho_{\mathrm{cue}}$: each causal cue sound, namely the phone buzz at $t=52$ and the door slam at $t=78$, remains at or above an audibility floor of $50$ dB.
\item $\rho_{\mathrm{know}}$: the schedule of what the audience has been told is unchanged.
\end{itemize}
Four alternatives are considered.
\begin{description}
\item[$V_1$, compressed.] The pause is shortened from $7$ s to $3$ s and everything after it moves earlier by $4$ s; the running time becomes $92$ s, a reduction of about $4.2\%$.
\item[$V_2$, uniformly quiet.] All non-dialogue sound is lowered by $4$ dB.
\item[$V_3$, stem-wise quiet.] The street noise is lowered by $1$ dB and the music by $7$ dB; the door slam and phone buzz are left alone (policy $S$ of the next section, which also discusses whether the slam should stay loud).
\item[$V_4$, early disclosure.] An insert shot at $t=18$ shows $A$ at the desk where the letter was forged, so that the audience learns that $A$ knows before $B$ asks the question.
\end{description}
Whether each version satisfies each relation can be settled by inspection of the specification, with the computations of the next section for the sound.

\begin{center}
\begin{tabular}{lccccc}
\toprule
 & $\rho_{\mathrm{ord}}$ & $\rho_{\mathrm{ant}}$ & $\rho_{\mathrm{dia}}$ & $\rho_{\mathrm{cue}}$ & $\rho_{\mathrm{know}}$\\
\midrule
$V_1$ compressed & holds & fails & holds & holds & holds\\
$V_2$ uniform $-4$ dB & holds & holds & holds & holds & holds\\
$V_3$ stem-wise & holds & holds & holds & holds & holds\\
$V_4$ early disclosure & holds & holds & holds & holds & fails\\
\bottomrule
\end{tabular}
\end{center}

The table lists guarantees. It does not list gains and losses, which are the quantities the transformation was undertaken for, and these differ.

\begin{itemize}
\item $V_1$ gains $4$ s of running time. In the stipulated response model of Section~\ref{sec:perturb}, shortening $p_1$ from $7$ to $3$ lowers anticipation $q_1$ from $0.905$ to $0.682$, a loss of $0.223$ that exceeds the tolerance $\varepsilon=0.05$. The contract that listed only $\rho_{\mathrm{ord}}$ would have licensed it.
\item $V_2$ reduces the number of exceeding steps from $14$ to $6$ (computed below). In the same model, lowering the music by $4$ dB lowers the fit measure $q_2$ from $0.679$ to $0.438$, a loss of $0.241$.
\item $V_3$ reduces exceeding steps from $14$ to $2$, leaves the causal cues at their original levels, and lowers the fit measure $q_2$ from $0.679$ to $0.269$. The occupation is better than in $V_2$ and the modeled music fit is worse, which is the conflict of Section~\ref{sec:conflict} in numerical form (the music fit and occupation pull in opposite directions in this model).
\item $V_4$ preserves every guarantee except $\rho_{\mathrm{know}}$, which it breaks on purpose. It changes the genre reading from mystery to dramatic irony, and no response model with parameters in $\R^p$ compares it with $V_0$, since the two differ by a discrete choice.
\end{itemize}

Two lessons follow. A version can satisfy every listed relation and still be one the author does not want, because the list omitted something that mattered, as with $V_3$ and the music. And the choice among $V_1,\dots,V_4$ is not a calculation: it is a decision about which relations the transformation should promise, and the contract records that decision where the viewer, the editor, and a generative system can read it.

\section{A Whisperizer-style attenuator, in detail}\label{sec:whisper}

Consider a tool, here called a Whisperizer, that lowers selected non-dialogue sound so that the scene's loud passages no longer crowd out speech. Four claims are commonly run together when such a tool is described, and they are logically independent.
\begin{enumerate}
\item \emph{Reduced occupation}: the exceedance set $E$ of Section~\ref{sec:occupation} shrinks or fragments.
\item \emph{Preserved intelligibility}: dialogue is unchanged and non-dialogue sound stays within a stated margin of it. This is a proxy defined by levels. Actual intelligibility depends on spectrum, masking, room, and listener, and is not captured by the margin.
\item \emph{Preserved causal cues}: sounds that are the only evidence of a plot event remain audible.
\item \emph{Preferred experience}: listeners judge the result better. This is empirical, and nothing below speaks to it.
\end{enumerate}

\paragraph{The scene's sound, stipulated.} Let the reference dialogue level be $D_t=60$ dB at every step and the margin $\mu=3$ dB, so a step exceeds when $L_t>63$. Non-dialogue events are: ambient sound at $40$ dB throughout; street noise at $64$ dB for $t=20,21,22$; the phone buzz at $54$ dB at $t=52$; the door slam at $72$ dB for $t=78,79$; and music beginning at $t=80$ with level $58+0.8k$ dB at $t=80+k$ for $k=0,\dots,15$. These events do not overlap in time, so $L_t$ at each step is the level of the single active event. The audibility floor is $50$ dB.

\paragraph{Uniform attenuation.} With no change, the exceeding steps are $t\in\{20,21,22\}$, $\{78,79\}$, and $\{87,\dots,95\}$ (the music exceeds $63$ once $0.8k>5$, that is for $k\ge7$): $14$ steps in three blocks, with recovery intervals of $55$ steps (23 to 77) and $7$ steps (80 to 86) between them. Under uniform attenuation by $a$ dB, \Cref{prop:mufflerizer} says the exceedance set shrinks as $a$ grows. The values are as follows.
\begin{center}
\begin{tabular}{rcccl}
\toprule
$a$ (dB) & exceeding steps & blocks & recovery intervals & buzz level (dB)\\
\midrule
$0$ & $14$ & $3$ & $55,\ 7$ & $54$\\
$4$ & $6$ & $2$ & $12$ & $50$\\
$6$ & $4$ & $2$ & $14$ & $48$\\
\bottomrule
\end{tabular}
\end{center}
At $a=4$ the street noise ($60$ dB) drops out, the door slam ($68$ dB) still exceeds, and the music exceeds only for $k\ge12$, that is $t=92,\dots,95$. At $a=6$ the music exceeds only at $t=94,95$. The last column shows the constraint of the cue relation: at $a=6$ the buzz is at $48$ dB, below the floor.

\begin{proposition}[Cue-safe uniform attenuation]\label{prop:cuesafe}
Let $\mathcal{C}$ be the set of causal cues with levels $L_c$ and floor $f$, with $L_c\ge f$. Uniform attenuation by $a\ge0$ preserves every cue above the floor if and only if $a\le a_{\max}=\min_{c\in\mathcal{C}}(L_c-f)$.
\end{proposition}

\begin{proof}
A cue $c$ stays at or above the floor iff $L_c-a\ge f$, that is $a\le L_c-f$. All cues stay iff $a$ is at most every one of these bounds.
\end{proof}

In the scene, $a_{\max}=\min(72-50,\,54-50)=4$ dB, and the quiet cue, the phone buzz, sets the limit. Combined with \Cref{prop:mufflerizer}, the best any uniform attenuation can do under $\rho_{\mathrm{cue}}$ is the $a=4$ row: six exceeding steps.

\paragraph{Stem-wise attenuation.} Treating each sound event as its own stem and choosing its gain separately removes the coupling between a loud non-cue and a quiet cue.

\begin{proposition}[Cue-safe elimination of a stem]\label{prop:stemwise}
Let a stem occupy a set of steps with levels $L_t\ge f$ and constant dialogue level $D$, and let $\mu$ be the margin. A single gain $a\ge0$ applied to the stem removes all its exceedances and keeps every step at or above the floor $f$ if and only if
\[
\max_tL_t-\min_tL_t\ \le\ D+\mu-f .
\]
\end{proposition}

\begin{proof}
Removing exceedance requires $L_t-a\le D+\mu$ for all $t$, that is $a\ge\max_tL_t-(D+\mu)$. Keeping the floor requires $a\le\min_tL_t-f$. Together with $a\ge0$, which is compatible with the second bound because $L_t\ge f$, such an $a$ exists iff $\max_tL_t-(D+\mu)\le\min_tL_t-f$, which rearranges to the stated inequality.
\end{proof}

Here $D+\mu-f=13$ dB. The music has levels from $58$ to $70$ dB, a range of $12\le13$, so it can be eliminated from the exceedance set with a gain $a\in[7,8]$ keeping it above the floor. The music is not a cue, so the contract does not require the floor of it, but the proposition shows how much room there is. At $a=7$ the music runs from $51$ to $63$ dB and never exceeds. The street noise needs only $a\ge1$. The door slam has range $0$ and needs $a\ge9$. Consider two policies.
\begin{itemize}
\item \emph{Policy $S$}: street $-1$ dB, music $-7$ dB, door slam and buzz untouched. The exceedance set is exactly $\{78,79\}$, a single block, with $78$ non-exceeding steps before it and $16$ after. The door slam retains its full $72$ dB.
\item \emph{Policy $S^*$}: as $S$ with the door slam lowered by $9$ dB to $63$ dB. The exceedance set is empty, and all contract relations $\rho_{\mathrm{dia}}$, $\rho_{\mathrm{cue}}$ hold.
\end{itemize}
\paragraph{Why the slam should or should not stay loud.} The contract, as stated, does not decide between $S$ and $S^*$, and it is worth being explicit about why. Suppose, as a stipulation, that the door is out of frame, so that the slam is the only evidence the audience has of $B$'s exit. Then the \emph{causal information} of the cue is that an exit occurred, and what preserving it requires is that the sound be audible and recognizable. That is $\rho_{\mathrm{cue}}$, with a floor, and $S^*$ meets it at $63$ dB. Nothing in the causal function requires the original $72$ dB. A tool whose purpose is to reduce intrusive effects has no reason internal to its contract to leave the slam alone; if the slam is to stay loud, the reason must come from the author.

The author's reason would be a different relation, \emph{salience} $\rho_{\mathrm{sal}}(\kappa)$: the cue stands above the dialogue by at least $\kappa$ dB, that is $L_c\ge D+\kappa$ at its steps. This expresses a claim about expressive function, that the slam punctuates $B$'s anger and not merely reports an exit, and that claim is a reading, not a consequence of the causal structure. The two relations interact as follows. Zero occupation requires $L_t\le D+\mu$ at every step, and $\rho_{\mathrm{sal}}(\kappa)$ requires $L_c\ge D+\kappa$ at the cue steps, so both can hold exactly when $\kappa\le\mu$; they conflict when $\kappa>\mu$. With $\mu=3$ dB, a slam at $63$ dB satisfies both for $\kappa\le3$ and $S^*$ is a legitimate choice. For $\kappa>3$, for instance $S$ with its slam $12$ dB over dialogue, occupation cannot reach zero and two exceeding steps remain. The threshold $\mu$ is a modeling parameter, and a slam at $63$ dB and one at $64$ dB are far apart for the proxy and presumably close for a listener, so both $\mu$ and $\kappa$ should be treated as decisions of the author, recorded in the contract, and tested with listeners, not as properties of sound.

\paragraph{Separating the four claims.} For policy $S$ the four claims have different statuses. Reduced occupation is a computed fact: $14\to2$ exceeding steps. Preserved intelligibility, in the proxy sense, holds except at the two slam steps, which exceed the margin by design. Preserved causal cues holds because both cues are unchanged. Preferred experience is not established: the stipulated model of Section~\ref{sec:perturb} even predicts a large loss in music fit for $S$, and an experiment with listeners could confirm or overturn that. A report on such a tool should therefore present these four as four rows, each with its own evidence, and should not assert the fourth on the strength of the first three.

\paragraph{Misclassification.} The analysis assumed that each sound event is correctly assigned to a stem. If the phone buzz is labeled as ambient noise and attenuated with it, then $\rho_{\mathrm{cue}}$ fails, and nothing in the guarantees above detects it. The contract of a stem-wise tool must therefore include the classifier's error rate on cues, and a cue-protection list supplied by the author is more reliable than inference.

\section{Contracts for the principal transformations}

Four families of transformation are common enough to describe by the form of their contracts.

\emph{Sensory remixing} changes levels or channel balance and promises preservation of dialogue content and timing, event order, and perhaps spatial cues, while the target change is the occupation structure. \emph{Temporal compression} promises preservation of explicit plot facts and causal order and cannot promise preservation of recovery or anticipation without further conditions, by the example above. \emph{Genre transformation} changes style cues and promises preservation of events, and by itself preserves neither stakes nor expected outcomes; to support the new mode it typically needs changes to what the audience is led to expect, that is, to the semantic fold of Chapter~\ref{ch:layers}, and so cannot in general preserve that fold; this is a design observation, not a theorem. \emph{Generative reconstruction} replaces a version by one produced from a description and promises satisfaction of the description, that is, membership in $\Mod(\Psi)$, and promises nothing about the particular film beyond that.

\section*{Exercises}

\begin{exercise}
For the compression example, suppose every quiet run of $x$ lasts at least $m_0\ge\tau_0$ seconds. Determine the smallest common factor $c<1$ for which $T$ still satisfies $\rho_{\mathrm{rec}}$, and describe a non-uniform compression that preserves $\rho_{\mathrm{rec}}$ while shortening the version.
\end{exercise}

\begin{exercise}
Verify \Cref{prop:arrangement} for $n=6$, $m=3$ by listing all binary sequences with three ones and counting blocks. Which values of $B$ occur and with what multiplicities?
\end{exercise}

\begin{exercise}
Define the dialogue-referenced exceedance differently, using a ratio of linear amplitudes in place of a difference of decibel levels, and show that the nesting result of \Cref{prop:mufflerizer} still holds.
\end{exercise}

\begin{exercise}
Write the contract of a genre transformation from drama to comedy that changes only music and color grading. State which preservation relations it can honestly claim and which it cannot.
\end{exercise}

\begin{exercise}
Using the sound specification of Section~\ref{sec:whisper}, find all real $a\ge0$ for which the number of exceeding steps is $4$ (the answer is an interval), and verify that none of them satisfies $\rho_{\mathrm{cue}}$ with floor $50$ dB.
\end{exercise}

\begin{exercise}
State and prove a version of \Cref{prop:stemwise} in which the dialogue level $D_t$ varies with $t$. What replaces the quantity $D+\mu-f$?
\end{exercise}

\chapter{Tropes and Hierarchical Direction}
\label{ch:direction}

Chapters~\ref{ch:types} and~\ref{ch:linear} gave a typed account of what may happen in a story. Two questions remain for a system that is meant to be directed. First, how can recurring narrative patterns, tropes, be represented so that they can be recognized in an existing film and selected for a new one? Second, how can a person express intentions at several scales, from genre to gesture, without specifying every movement? This chapter answers both with the machinery already built, and it identifies where the machinery reaches its limits.

\section{Tropes as schemas}

A trope label rarely specifies enough to reconstruct a scene. The useful object is a \emph{schema}: roles, enabling conditions, expected developments, and timing.

\begin{definition}[Trope schema]
A \emph{trope schema} $\mathsf{T}$ consists of a context $\Gamma_{\mathsf{T}}$ declaring its roles, a type $\mathsf{Pre}_{\mathsf{T}}(\Gamma_{\mathsf{T}})$ of enabling conditions, a set $\mathsf{Dev}_{\mathsf{T}}$ of admissible developments, each an event word over the story's event alphabet, and a timing window $W_{\mathsf{T}}\subseteq\R_{\ge0}\times\R_{\ge0}$ of permitted onset and duration. An \emph{instance} of $\mathsf{T}$ in a film is a substitution of the roles by characters or objects, a witness of the enabling conditions in the story state at the onset, and a development from $\mathsf{Dev}_{\mathsf{T}}$ whose timing lies in $W_{\mathsf{T}}$.
\end{definition}

A schema is thus a typed fragment of the kind introduced in Chapter~\ref{ch:types}, and its instance requires a witness: a trope cannot occur unless its enabling conditions hold. This yields the compatibility question of trope composition. Two schemas compose only if the developments of the first supply, as postconditions, the enabling conditions of the second. A sequence of recognizable tropes without such links is a list, not a plot.

\section{Reverse inference of tropes}

Given a finished film, which schemas does it instantiate? In the vocabulary of Chapter~\ref{ch:scripts}, let $F$ be a set of films and let the constraints $\Phi$ be statements of the form ``contains an instance of schema $\mathsf{T}$ at interval $I$''. Then the set of schemas a film $f$ satisfies is $\Th(\{f\})$ restricted to this vocabulary.

\begin{proposition}[Underdetermination]\label{prop:underdet}
Let $\Phi_{\mathsf{trope}}$ be a vocabulary of trope constraints and $f$ a film. The set $\Th(\{f\})\cap\Phi_{\mathsf{trope}}$ is closed under the consequences of the satisfaction relation: if $\psi$ is satisfied by every film that satisfies all of a subset $\Psi\subseteq\Th(\{f\})$, then $\psi\in\Th(\{f\})$. Distinct films may have the same trope theory, and a trope theory need not determine a unique film.
\end{proposition}

\begin{proof}
If $\psi$ holds in every model of $\Psi$, and $f\in\Mod(\Psi)$ because $\Psi\subseteq\Th(\{f\})$, then $f\models\psi$, so $\psi\in\Th(\{f\})$. For the second statement, take two films that agree on every constraint in $\Phi_{\mathsf{trope}}$ and differ only in a feature outside that vocabulary, such as the color of a costume; they have the same trope theory. Here $\psi$ is taken in $\Phi_{\mathsf{trope}}$ for the intersection to be defined.
\end{proof}

The proposition formalizes a methodological warning from Chapter~\ref{ch:scripts}. A trope analysis recovers a closed set of constraints, and the original film is only one of the many that realize it. Reverse inference from observed events to schemas is further complicated by the fact that a schema instance is a hypothesis about function, while the observed events are features: the same observable event may instantiate different tropes depending on what the audience knows. Automatic suggestions of tropes should accordingly be treated as hypotheses for review, each with the evidence that supports it.

\section{Commands at several scales}

A director's instruction to an actor and a user's instruction to a simulated character are both commands issued at some level of abstraction. A useful scheme distinguishes levels from world and genre through plot, sequence, scene, relationship, goal, action, and gesture. A higher-level command constrains lower-level choices without determining them.

\begin{definition}[Command language]
A \emph{command alphabet} $K$ is a finite set of keys. A \emph{command} is a word over $K$ selected from a set $\mathcal{K}\subseteq K^*$ of valid command words. Each command $c\in\mathcal{K}$ has a \emph{scope}, a \emph{precondition} that is a type in the context of the current story state, and an \emph{effect}, which is a set $R(c)\subseteq E^*$ of event words called its \emph{realizations}.
\end{definition}

A command at the level of a gesture has a small set of realizations, and a command at the level of a plot has a large one. A realization is \emph{acceptable} if it satisfies the preservation conditions of the command's contract, which for a dramatic command will include the preservation of the intended function. The relation between a high-level command and its realizations is a refinement.

\begin{proposition}[Prefix-free commands decode uniquely]\label{prop:prefixfree}
If $\mathcal{K}$ is a set of nonempty words that is prefix-free, that is, no command is a proper prefix of another, then every word $w\in K^*$ that is a concatenation of commands has a unique factorization $w=c_1c_2\cdots c_r$ with each $c_i\in\mathcal{K}$.
\end{proposition}

\begin{proof}
Suppose $c_1\cdots c_r=c_1'\cdots c_s'$ are two factorizations, and let $i$ be the first index with $c_i\ne c_i'$. After removing the equal common prefix, the two words begin with $c_i$ and $c_i'$ respectively, so one is a prefix of the other, which contradicts prefix-freeness unless they are equal. The same argument applied inductively shows $r=s$ and $c_i=c_i'$ for all $i$.
\end{proof}

Hierarchical, mnemonic key sequences of the kind found in modal editors organize commands as paths in a tree. If every command is a leaf, the key set is prefix-free and every keystroke stream parses uniquely, which is the property that makes menus of this type unambiguous to a machine. The cost is that no command can be a prefix of another, so a command cannot be both complete and extensible. Interfaces resolve this by terminating commands with an explicit key or a timeout, which restores prefix-freeness at the price of an additional symbol. The design question of discoverability and memorization is empirical, and the notation of any specific interface should be designed and tested rather than inferred from the formal property.

\section{Command semantics: order matters}

Commands act on a story state, and their order can change the outcome, because one may establish what another requires.

\begin{example}[Non-commuting commands]\label{ex:noncomm}
Let $c_1$ be the command ``character $A$ learns that the letter is forged'', whose effect is to add the fact $\mathsf{Knows}(A,\mathsf{forged})$ to the state, and let $c_2$ be ``$A$ confronts $B$ about the forgery'', whose precondition is $\mathsf{Knows}(A,\mathsf{forged})$. The sequence $c_1c_2$ is valid. The sequence $c_2c_1$ is invalid, since the precondition of $c_2$ is not satisfied in the initial state. The commands do not commute.
\end{example}

The semantics of \Cref{ex:noncomm} is the semantics of Chapter~\ref{ch:linear} read in the command layer: a command is a rule with premises and conclusions, and a sequence is valid when each command's premises are present in the current marking. Disjointness of resources gives commutation, by \Cref{prop:commute}. Here the shared fact $\mathsf{Knows}(A,\mathsf{forged})$ is exactly what couples the two commands.

\begin{proposition}[Realizations of a sequence]\label{prop:realize}
Let $c_1,\dots,c_r$ be a valid command sequence with realization sets $R(c_i)\subseteq E^*$. Define a \emph{realization of the sequence} to be a concatenation $w_1\cdots w_r$ with $w_i\in R(c_i)$ that is an admissible history. Then the set of realizations is exactly
\[
\{w_1w_2\cdots w_r:\ w_i\in R(c_i)\}\cap\Adm,
\]
where $\Adm$ is the set of admissible histories, and this may be a proper subset of the unconstrained product because a choice of realization for one command may constrain the choices available for later ones.
\end{proposition}

\begin{proof}
Every realization of the sequence is a concatenation of realizations of its commands, which gives the first factor, and every realization must be an admissible history, which gives the intersection with $\Adm$. A realization of $c_1$ that consumes a resource needed by every realization of $c_2$ makes any concatenation using it inadmissible, so the intersection can be a proper subset of the product.
\end{proof}

This is the formal reason why high-level direction of characters cannot be treated as independent instructions: later commands inherit the consequences of earlier realizations. Systems that support such direction need to expose conflicts and not resolve them silently, and to preview alternative realizations at the semantic level where the user is working. They also need undo and revision at that same level, so that a correction to one command does not silently invalidate a setup that a later command depends on, a requirement that follows from the preconditions themselves: a change to an early command can invalidate all later commands whose preconditions it supplied.

\section{A complete trope example: concealed knowledge}\label{sec:tropeexample}

Return to the forged-letter scene. Two schemas in the sense of the definition above are relevant, and the difference between them lies in which layer each constrains.

\begin{definition}[Two schemas on different layers]
\emph{Concealed knowledge}, $\mathsf{T}_{\mathrm{con}}$, has roles $K,I$ (knower and uninformed) and a fact $\varphi$. Its enabling condition is $\mathsf{Knows}(K,\varphi)\times\neg\mathsf{Knows}(I,\varphi)$ in the \emph{state} layer. Its developments are
\[
\mathsf{Dev}=\{\mathsf{ask}\,\mathsf{deflect}\,\mathsf{reveal},\ \mathsf{ask}\,\mathsf{deflect}\,\mathsf{deflect}\,\mathsf{reveal}\},
\]
where $\mathsf{ask}$ is a question by $I$ bearing on $\varphi$, $\mathsf{deflect}$ is an evasive response by $K$, and $\mathsf{reveal}$ adds $\mathsf{Knows}(I,\varphi)$. Its timing window requires the first $\mathsf{deflect}$ to begin within $10$ s of $\mathsf{ask}$ and $\mathsf{reveal}$ to occur within $30$ s of $\mathsf{ask}$.

\emph{Dramatic irony}, $\mathsf{T}_{\mathrm{iro}}$, has the same roles and developments and, in addition, requires the \emph{audience} layer to contain $\mathsf{Knows}_{\mathrm{aud}}(\varphi)$ at the onset of $\mathsf{ask}$.
\end{definition}

\paragraph{Instance in the baseline.} In $V_0$ the history contains $\mathsf{ask}$ at $t=38$ (question), $\mathsf{deflect}$ at $t=46$, and $\mathsf{reveal}$ at $t=58$. The word $\mathsf{ask}\,\mathsf{deflect}\,\mathsf{reveal}$ is in $\mathsf{Dev}$. The deflection begins $8$ s after the question, within the $10$ s bound, and the reveal occurs $20$ s after the question, within $30$ s. The enabling condition holds throughout $[0,58)$ because $A$ knows from the start, so $\mathsf{T}_{\mathrm{con}}$ has an instance with onset at $t=38$. The schema $\mathsf{T}_{\mathrm{iro}}$ does \emph{not} have an instance in $V_0$, because the audience layer has not been told that $A$ knows until $t=58$. In $V_4$, where the insert at $t=18$ tells the audience, $\mathsf{T}_{\mathrm{iro}}$ has an instance and $\mathsf{T}_{\mathrm{con}}$ has the same one. The two schemas therefore separate the versions $V_0$ and $V_4$, even though the state-layer history is identical.

\paragraph{Reverse inference and a competing hypothesis.} A viewer, or a system annotating $V_0$, sees only observable events: $A$ glances at the letter ($t=18$), pauses ($31$ to $38$), is asked a question ($38$), responds evasively ($46$), and says something that reveals the forgery ($58$). From the pause and the question-and-evasion exchange alone, $t=31$ to $46$, a second hypothesis is equally supported. Let $\mathsf{T}_{\mathrm{int}}$ be an \emph{interrogation} schema in which $B$ presses $A$ about a matter $\psi$ unrelated to the letter, for example where $A$ was the previous evening. Under $\mathsf{T}_{\mathrm{int}}$ the same word $\mathsf{ask}\,\mathsf{deflect}$ occurs with $\psi$ in place of $\varphi$. This is a separate point from \Cref{prop:underdet}, which concerns the film: here the observables of this stretch do not determine the theory, as the table below shows.

\begin{center}
\small
\begin{tabular}{p{5.8cm}p{3cm}p{3cm}}
\toprule
observable & $\mathsf{T}_{\mathrm{con}}$ with $\varphi=$ forged & $\mathsf{T}_{\mathrm{int}}$ with $\psi$ unrelated\\
\midrule
$A$ glances at the letter, $t=18$ & expected & unexplained\\
long pause, $t=31$ to $38$ & compatible & compatible\\
question and evasion, $t=38$ to $46$ & compatible & compatible\\
(absent in $V_0$) $A$'s eyes go to the letter during the evasion & expected & unexplained\\
(absent in $V_0$) the evasion names the letter in dialogue & expected & excluded\\
\bottomrule
\end{tabular}
\end{center}

The pause and the question-and-evasion exchange do not discriminate. Observations that discriminate are those that lie in the evidence set of one hypothesis and not the other: a glance toward the letter during the evasion, or a line of dialogue that names the letter. Neither is part of the baseline specification, and adding either to the scene would settle the reading earlier. The annotation of $V_0$ as an instance of $\mathsf{T}_{\mathrm{con}}$ is therefore a hypothesis resting on the glance at $t=18$ and on the reveal at $t=58$, and until $t=58$ it is one of two readings. This is the point of the earlier warning: the reveal at $t=58$ is the first event that collapses the ambiguity, and a film designer who wants the audience to hold both readings has a measurable interval, $t=38$ to $58$, in which to maintain them.

\paragraph{Composition.} The reveal adds $\mathsf{Knows}(B,\varphi)$, which is the enabling condition of two further schemas that depend on it independently. The first, $\mathsf{T}_{\mathrm{rea}}$, is a reaction in which $B$'s face and posture register the knowledge; its timing window requires onset within $5$ s of the reveal, and the reaction shot at $t=62$, $4$ s after $t=58$, is within it (the duration of the shot is not constrained by the window). The second, $\mathsf{T}_{\mathrm{dep}}$, is $B$'s departure, and the door exit at $t=78$ is its development; its enabling condition is again $\mathsf{Knows}(B,\varphi)$, supplied by the reveal and not by the reaction. The three schemas form a plot in the sense of the earlier section because the postcondition of $\mathsf{T}_{\mathrm{con}}$ supplies the enabling conditions of the other two, and removing the reveal removes both.

\section{A complete session of hierarchical direction}\label{sec:session}

A command language for the scene needs a few commands, all of the same length so that the key set is prefix-free. Each command is a path of three keys: a scale, a category, and an action.

\begin{center}
\small
\begin{tabular}{p{1.0cm}p{3.3cm}p{3.7cm}p{3.1cm}}
\toprule
keys & meaning & precondition & effect\\
\midrule
\texttt{s k l} & $A$ learns the letter is forged & none & $\mathsf{Knows}(A,\varphi)$\\
\texttt{s k c} & $A$ conceals & $\mathsf{Knows}(A,\varphi)$, $\neg\mathsf{Knows}(B,\varphi)$ & concealment active\\
\texttt{s k r} & reveal to $B$ & concealment active & $\mathsf{Knows}(B,\varphi)$\\
\texttt{g r b} & reaction shot of $B$ & $\mathsf{Knows}(B,\varphi)$ & reaction recorded\\
\texttt{s d x} & $B$ exits, door slam & reaction recorded & $B$ gone\\
\texttt{s m q} & music with a quiet profile & none & music stem policy set\\
\bottomrule
\end{tabular}
\end{center}

Because every command is a word of exactly three keys, no command is a proper prefix of another, so by \Cref{prop:prefixfree} any concatenation parses uniquely. The user can type a continuous stream of eighteen keys, and the interface need not use a terminator.

\paragraph{The valid orderings.} The preconditions of the table impose the chain
\[
\texttt{skl}<\texttt{skc}<\texttt{skr}<\texttt{grb}<\texttt{sdx}
\]
and leave \texttt{smq} unconstrained. Of the $6!=720$ orderings of the six commands, the valid ones are those that place the five chained commands in order and insert \texttt{smq} at any of the six positions, so there are exactly $6$, a count confirmed by enumerating all $720$ orderings against the preconditions of the table (the script \texttt{verify.py} distributed with the source does this). If the user types \texttt{skc} first, the system refuses at step one with the message that $\mathsf{Knows}(A,\varphi)$ is absent. If the user types \texttt{skl skr}, it refuses at step two because no concealment is active. In both cases the refusal names the precondition, not merely the command, so that the user can repair the plan. The order \texttt{skl skc skr grb sdx smq} is valid, and so are the other five that move \texttt{smq} earlier.

\paragraph{Where choices enter.} A command fixes a semantic effect and leaves discretionary choices open. Three are stipulated here: the form of the deflection under \texttt{skc} (spoken or by gesture), the length of the reaction shot under \texttt{grb} in seconds (any of $2,3,4,5,6$), and the character of the slam under \texttt{sdx} (hard at $72$ dB or soft at $51$ dB). The product of the choice sets has $2\times5\times2=20$ elements. An admissibility rule couples two choices: a soft slam is admissible only if the reaction lasted at least $4$ s, because the audience needs the face to carry the departure. For a hard slam all five durations are admissible, and for a soft slam three are. The admissible realizations number $2\times(5+3)=16$, a proper subset of the $20$ in the product, in agreement with \Cref{prop:realize}.

The choices do not all enter the history. The choice between spoken and gestured deflection leaves the semantic fold unchanged, since both yield the same state, and so it belongs to the appearance layer only. The choice of reaction duration affects the timing window of $\mathsf{T}_{\mathrm{rea}}$ and the audience's access to $B$'s state, and the choice of slam affects whether the cue is salient in the sense of Section~\ref{sec:whisper}. A system that records only the commands and not the realized choices loses exactly this information. A system that records the realized choices in the history, when they affect the fold or an enabling condition, and in the appearance layer otherwise, has what it needs to answer why a later command is or is not admissible.

\paragraph{Revision.} Suppose the user deletes \texttt{skl}. Because preconditions are read in the state the earlier commands produced, every command whose precondition depends on $\mathsf{Knows}(A,\varphi)$ becomes invalid: \texttt{skc}, then \texttt{skr}, \texttt{grb}, and \texttt{sdx}. Only \texttt{smq} survives. A system that edits at the command level should present this cascade before applying the deletion, and offer the alternatives: delete all dependents, or replace \texttt{skl} by another command that establishes the same fact. The cascade is exactly the dependency graph of the five chained commands, which the enumeration in the accompanying script reproduces.

\section*{Exercises}

\begin{exercise}
Write a trope schema for a mistaken-identity plot in the sense of the definition above, specifying roles, enabling conditions, developments, and timing window. Give an instance in a film of your choice and identify the witness of the enabling conditions.
\end{exercise}

\begin{exercise}
Show that if $\mathcal{K}$ is not prefix-free, the factorization of \Cref{prop:prefixfree} can fail to be unique. Give an example over a two-letter alphabet.
\end{exercise}

\begin{exercise}
Using the resource semantics of Chapter~\ref{ch:linear}, model three commands with a shared resource so that exactly two of the six possible orderings are valid. Describe the story-level interpretation.
\end{exercise}

\begin{exercise}
A user issues a high-level command whose set of realizations $R(c)$ is empty in the current state. List three distinct responses a system might give (refuse, replan earlier commands, request revision) and state which formal property of $\Adm$ each relies on.
\end{exercise}

\begin{exercise}
In the command language of Section~\ref{sec:session}, add a command \texttt{s k a} (``$A$ confesses unprompted'') with precondition $\mathsf{Knows}(A,\varphi)\times\neg\mathsf{Knows}(B,\varphi)$ and effect $\mathsf{Knows}(B,\varphi)$. Determine which pairs of commands cannot occur together in a valid sequence, and list the maximal valid sequences.
\end{exercise}

\begin{exercise}
Write a trope schema for \emph{red herring} in the sense of Section~\ref{sec:tropeexample}, giving roles, enabling conditions, developments, and a timing window, and say which observable would discriminate it from a genuine clue.
\end{exercise}


\part{Synthesis}
\chapter{Case Studies}
\label{ch:cases}

The preceding chapters introduced tools for one question at a time. This chapter applies several of them together to a small number of well-known films, and to one synthetic scene. The claims made about each film are claims about what the model can state, given a feature of the film that is widely discussed. They are readings, offered as demonstrations of method, and none depends on attributing intentions to the makers.

\section{Rope: seamless composition with hidden seams}

\emph{Rope} (1948) was designed to appear as a continuous take, and is assembled from a small number of long takes joined by concealed cuts, many of them masked by a dark surface filling the frame. In the vocabulary of Chapter~\ref{ch:scenario}, the film is a word in the free monoid of shots whose seams have $\delta_\Img$ small while $\delta_\Sc$ need not be zero: the camera comes to rest against a dark back, the next take begins from the same configuration, and the two world states are close in appearance. The audience reads the sequence as a morphism in $\Path(\Sc)$, a single path, although it is a composite of path segments joined at matching seams.

The model separates two things that viewers may conflate. The film's effect of unbroken real time is the effect of $\delta_\Img$ being small at every seam. Its actual construction as a word of several shots is the fact that $\delta_\Sc$ need not vanish. By \Cref{prop:seam-lipschitz}, if the rendering is Lipschitz then small $\delta_\Sc$ suffices for small $\delta_\Img$, but the converse is the technique of the concealed cut: a large jump in world state hidden behind a rendering that cannot show it.

\section{Last Year at Marienbad: local consistency, global failure}

\emph{Last Year at Marienbad} (1961) is often described as one whose scenes, individually coherent, resist assembly into a single chronology. One reading in the present vocabulary treats each recalled scene as a local section of a sheaf of scene states over a cover of the film's story world, with each pair of overlapping recollections agreeing on what they share. If the family is compatible on overlaps and admits no global section, it is a contextual family in the sense of Chapter~\ref{ch:descent}. The reading makes a precise claim that can be argued against: that every pair, and indeed every small collection of scenes, admits a consistent state, while no state accounts for all. A critic who finds a consistent global chronology rejects the reading by exhibiting the global section. A critic who finds a pair of scenes that cannot both be true of one world rejects the claim of local consistency.

\section{Groundhog Day: appearance without history}

\emph{Groundhog Day} (1993) presents a day that repeats while one character retains knowledge across repetitions. As developed in \Cref{ex:groundhog}, the appearance is periodic while the state is not. The world clock obstruction of \Cref{prop:loop} applies to any attempt to give the film a global single-valued time with the image sequence as its only record. The obstruction disappears when the record of accumulated knowledge is placed in the history and the fold, which makes the world state non-periodic. This is the formal sense in which the film is about a history that the appearance does not show: two histories at the dawn of the first and tenth day have identical images and different admissible futures, which is \Cref{prop:appearance}.

The same example illustrates the preservation contracts of Chapter~\ref{ch:contracts}. A transformation that compresses the film by removing apparently repeated days would preserve event order only if every repetition were removable without effect, and fails to preserve the semantic fold, since the knowledge the protagonist accumulates is exactly the content of the repetitions.

\section{Rashomon: parallel accounts as distinct paths}

\emph{Rashomon} (1950) presents several accounts of one event. In the type-theoretic reading of \Cref{ex:rashomon}, the accounts are paths between the same two states of the world, and the film's concern is that they are not equal as paths even though they share endpoints. The reading separates the question of outcome, which is shared, from the question of how the outcome arose, which is contested. What a viewer must hold in mind is a family of paths and the claim, which the film neither confirms nor refutes, that no one of them is privileged.

\section{Perfect Days: a flat narrative potential}

\emph{Perfect Days} (2023) follows the repeated daily routine of a man in Tokyo, in a form that is repetitive without being a time loop. In the geometric vocabulary of Chapter~\ref{ch:scenario}, the narrative cost $U$ is nearly constant along the film's path, so the action functional is dominated by the perceptual term and critical paths are close to geodesics of the perceptual metric. The film's repetition is not a periodicity of the world clock but a recurrence of appearance with small, accumulating changes in state. Its interest, in this reading, lies in the history that the repeated appearances conceal, and in the small events that modify the fold without altering the image.

\section{A synthetic scene: one participant knows what the other does not}

A teaching scene that can be varied at will illustrates the machinery in one place. Two characters, $A$ and $B$, converse in a room. $A$ knows that the letter on the table is forged, and $B$ does not. The scene can be realized in several versions that differ in timing, soundtrack, cutting, and gesture.

\emph{Types.} The story state contains $\mathsf{Knows}(A,\mathsf{forged})$ and lacks $\mathsf{Knows}(B,\mathsf{forged})$. The event $\mathsf{reveal}(A,B)$ has the first as a precondition and, as an effect, adds the second. A version in which $B$ acts on the forgery before the reveal is ill typed.

\emph{State and appearance.} A shot of $A$ glancing at the letter has the same image whether or not $A$ knows. The state layer distinguishes the two cases, so by \Cref{prop:appearance} the image cannot determine the admissible futures, and the audience's inferences depend on what the history has established.

\emph{Transformation.} A compression that removes a pause before the reveal preserves event order. If the pause was an opportunity for the audience to infer $A$'s knowledge, it may fail a recovery or anticipation relation. A contract that lists only order permits the compression, and one that lists anticipation does not.

\emph{Equivalence.} Moving the reveal by $\Delta$ seconds changes comprehension for $\Delta$ above some threshold that depends on the viewer and task. By \Cref{prop:nontransitive}, a sequence of small shifts of the reveal can accumulate to a visible change although each step is below threshold.

\emph{Direction.} The two commands ``$A$ learns the letter is forged'' and ``$A$ confronts $B$'' do not commute (\Cref{ex:noncomm}). A director issuing them in the wrong order obtains an inadmissible version, and the system must either refuse or replan.

The scene is synthetic and its soundtrack and cutting can be varied without referring to commercial footage, which makes it suitable for experiments. Findings from such scenes are findings about the constructed variations, and they do not generalize automatically to existing films.

\section{The synthetic scene, analyzed by time}\label{sec:timeindexed}

The previous section applied each tool once. This section revisits the forged-letter scene segment by segment, using the specification of Section~\ref{sec:scene}, and for each segment records five things: what is \emph{observed} on screen and sound, how it is \emph{annotated} in the layers of Chapter~\ref{ch:layers}, the \emph{formal} object from earlier chapters that represents it, an \emph{alternative} edit and the relation it tests, and an \emph{evaluation} question that an experiment or reader could answer. Everything in the observation lines is stipulated by the specification, so the analysis is checkable against it. The evaluation lines are open.

\subsection*{$0$ to $6$ s: entrance}
\begin{description}
\item[Observed.] $A$ enters the kitchen. Ambient sound at $40$ dB.
\item[Annotated.] History: $\mathsf{enter}(A)$. State: $A$ knows the letter is forged. Appearance: nothing marks this. Audience: no knowledge of the forgery yet.
\item[Formal.] A seamless segment, a morphism in $\Path(\Sc)$; no seam. The state and appearance layers already differ here (\Cref{prop:appearance}).
\item[Alternative.] Cut the entrance in two shots. This introduces a seam and tests continuity in the sense of Chapter~\ref{ch:categories}.
\item[Evaluation.] Does a viewer perceive the second version as having a cut at the seam, given a match on $A$'s position?
\end{description}

\subsection*{$6$ to $18$ s: greeting}
\begin{description}
\item[Observed.] $B$ greets $A$ at $t=6$. Ordinary dialogue.
\item[Annotated.] History: $\mathsf{greet}(B,A)$. Neither character acts on the letter.
\item[Formal.] A word in the free monoid of shots. Cut density is a statistic of the word, not of the story (Chapter~\ref{ch:editing}).
\item[Alternative.] Replace three cuts by one: the cut rate drops, the events stay.
\item[Evaluation.] Is the perceived pace changed (the tolerance $\varepsilon$ of Chapter~\ref{ch:equivalence}) when the cut count changes and the events do not?
\end{description}

\subsection*{$18$ to $31$ s: the glance and the street}
\begin{description}
\item[Observed.] At $t=18$, $A$ glances at the letter. At $t=20$ to $22$ street noise at $64$ dB.
\item[Annotated.] The glance is an appearance with no history event; the street noise is a non-cue exceeding event.
\item[Formal.] The glance is the observable that supports $\mathsf{T}_{\mathrm{con}}$ over $\mathsf{T}_{\mathrm{int}}$ (Section~\ref{sec:tropeexample}); the street noise is a block of three exceeding steps (Section~\ref{sec:whisper}).
\item[Alternative.] Remove the glance (state unchanged, appearance changed), or lower the street noise by $1$ dB (occupation changed, no other relation).
\item[Evaluation.] Does removing the glance change which reading viewers report at $t=46$?
\end{description}

\subsection*{$31$ to $38$ s: the pause}
\begin{description}
\item[Observed.] A $7$ s pause with no dialogue and ambient sound only.
\item[Annotated.] Quiet interval. No history event.
\item[Formal.] The first moment where $\rho_{\mathrm{ant}}$ of Section~\ref{sec:altversions} bites: a quiet run of $7$ s against a threshold $\tau_0=5$ s. Parameter $p_1$ of Section~\ref{sec:perturb}.
\item[Alternative.] Shorten to $3$ s ($V_1$): order preserved, anticipation relation broken.
\item[Evaluation.] Between $3$ s and $7$ s, where does a viewer population's judgment of anticipation drop by more than $\varepsilon$, and is the dependence asymmetric as the model predicts?
\end{description}

\subsection*{$38$ to $58$ s: question, deflection, buzz}
\begin{description}
\item[Observed.] $B$ asks the question at $t=38$. $A$ deflects at $t=46$. $B$'s phone buzzes at $t=52$ at $54$ dB.
\item[Annotated.] History: $\mathsf{ask}$, $\mathsf{deflect}$, and the buzz event. The buzz is a causal cue.
\item[Formal.] The word $\mathsf{ask}\,\mathsf{deflect}$ is a prefix of a development of $\mathsf{T}_{\mathrm{con}}$, with the timing window checks of Section~\ref{sec:tropeexample}. The buzz fixes the cue limit $a_{\max}=4$ dB (\Cref{prop:cuesafe}).
\item[Alternative.] Reorder: place the audience disclosure at $t=18$ ($V_4$). The type of the question changes from mystery to irony.
\item[Evaluation.] Which of the two readings, $\mathsf{T}_{\mathrm{con}}$ or $\mathsf{T}_{\mathrm{int}}$, do viewers report at $t=46$, with and without the discriminating observation?
\end{description}

\subsection*{$58$ to $66$ s: the reveal and the reaction}
\begin{description}
\item[Observed.] $A$ reveals the forgery at $t=58$. A reaction shot of $B$ from $t=62$ to $66$.
\item[Annotated.] History: $\mathsf{reveal}$, then $\mathsf{react}$. The reveal adds $\mathsf{Knows}(B,\varphi)$. Audience: now informed.
\item[Formal.] The only step at which both the character-knowledge and audience-knowledge processes change. Montage surplus: the reaction shot's reading depends on the preceding line (Chapter~\ref{ch:editing}). Reaction duration is parameter $p_4$.
\item[Alternative.] Place the reaction before the reveal line, or lengthen it from $4$ s to $6$ s.
\item[Evaluation.] Is the reaction read as a response to the reveal in the reordered version? By how much does perceived salience change with duration?
\end{description}

\subsection*{$66$ to $78$ s: the aftermath}
\begin{description}
\item[Observed.] Dialogue and ambient only; a quiet run of $12$ steps.
\item[Annotated.] No history event occurs here; the segment carries the aftermath as appearance only.
\item[Formal.] Part of the $55$-step recovery interval between the street noise and the slam (Section~\ref{sec:whisper}), which is the longest quiet run in the scene and is untouched by any attenuation policy considered there.
\item[Alternative.] Insert an additional exchange ($+6$ s): changes running time and the placement of the slam.
\item[Evaluation.] Does the additional exchange change the perceived impact of the slam?
\end{description}

\subsection*{$78$ to $80$ s: the exit}
\begin{description}
\item[Observed.] $B$ leaves. Door slam at $72$ dB for $2$ steps.
\item[Annotated.] History: $\mathsf{exit}(B)$, a causal cue. State: $B$ is no longer present, so any further command that requires $B$ is inadmissible (Chapter~\ref{ch:linear}).
\item[Formal.] An exceeding block of two steps, and $\rho_{\mathrm{cue}}$, with $\rho_{\mathrm{sal}}(\kappa)$ for $\kappa>\mu$ in conflict with zero occupation (Section~\ref{sec:whisper}); a command \texttt{s d x} with its precondition (Section~\ref{sec:session}).
\item[Alternative.] $S^*$ (slam lowered by $9$ dB), or a soft slam of $51$ dB.
\item[Evaluation.] Do listeners prefer $S$ to $S^*$? Is the door slam still read as the end of the exchange at $63$ dB, and what salience margin $\kappa$ does the author intend?
\end{description}

\subsection*{$80$ to $96$ s: the music}
\begin{description}
\item[Observed.] Music enters at $t=80$ at $58$ dB and rises by $0.8$ dB per second to $70$ dB at $t=95$. The scene ends at $96$.
\item[Annotated.] Appearance only; no history event. A trailing exceeding block $t=87$ to $95$.
\item[Formal.] The largest block of occupation (nine of the $14$ exceeding steps). Music parameters $p_2,p_3$.
\item[Alternative.] Lower the music by $7$ dB (policy $S$) so that the block disappears.
\item[Evaluation.] What is lost: the modeled fit measure $q_2$ falls from $0.68$ to $0.27$. Is this reproduced by listeners?
\end{description}

The analysis shows how the tools divide the work. Only two events change the semantic state, the reveal and the exit. Pace statistics and sound levels describe appearance, tropes constrain state and audience knowledge together, and contracts name which of these a transformation may touch. The same segment appears in the analyses of several chapters, which is the purpose of carrying one scene through the book.

\section{A protocol for analyzing a real passage}\label{sec:protocol}

The film subsections of this chapter are readings, and none is a time-indexed analysis of footage. A reader who wants to carry out such an analysis on an actual film needs source material that this book does not supply, and should supply it with its provenance. The following protocol is the form that the synthetic analysis above takes when applied to a real passage. Section~\ref{sec:real} carries it out on a passage for which measurements exist.
\begin{enumerate}
\item \textbf{Source.} State the title, the edition or release, the runtime, and the player or medium. Timecodes belong to an edition and differ between releases, so they must be tied to one.
\item \textbf{Passage.} Give start and end timecodes, in a stated precision, for each segment analyzed.
\item \textbf{Observation.} For each segment, list what is seen and heard, in terms that another viewer can check without accepting the reading: the shots, the cuts, the dialogue, and the sound events, with timecodes.
\item \textbf{Annotation.} Assign each observation to a layer: history event, state fact, appearance, or audience knowledge. Mark which assignments are inferences.
\item \textbf{Claim type.} Mark each claim as an observation, a reading, or a hypothesis about perception, so that the evidence required differs correctly.
\item \textbf{Formal representation.} Name the object used (a path, a section of a sheaf, a schema instance, a command), and say which checks it passes. A representation that cannot fail a check is decoration.
\item \textbf{Alternative edit.} State a specific change, the relation it is meant to preserve, and the relation it breaks.
\item \textbf{Evaluation.} State the question a viewer study would answer, the population, and the tolerance.
\end{enumerate}
The protocol separates what can be verified from the film (items 1 to 3) from what is interpretive (4 to 6) and from what is empirical (7 and 8). A careful analysis keeps the three apart.

\section{A measured passage: \emph{Drive Through Fire}}\label{sec:real}

The synthetic scene is clear because it was built to be. To test whether that clarity made the framework look easier to apply than it is, the protocol of Section~\ref{sec:protocol} is applied here to a real passage for which measurements exist: the soundtrack analysis of the 2026 feature \emph{Drive Through Fire} in the author's monograph on acoustic occupation~\cite{flyxion2026occupation}. The numbers below come from two sources: Tables~7.3 to~7.5 and Section~7.5.7 of that monograph, and, for quantities the monograph does not report, a recomputation from its per-second table (\texttt{audio-analysis/results/seconds.csv} in the author's repository, commit \texttt{1c016ab}). The recomputation reproduces the monograph's nine occupation blocks exactly (the script \texttt{real\_passage.py} distributed with the source does this, and \texttt{verify.py} checks the transcribed tables for arithmetic consistency). The film itself, its picture, and its story were not available to this book. The recomputation shows that the tables follow from the per-second data, not that the per-second data follow from the film.

\paragraph{1. Source.} The 2026 release analysed in the monograph, runtime $1{:}27{:}02$. The level measure is ITU-R BS.1770 momentary loudness of the 5.1 mix, energy-averaged per second, in LUFS. The dialogue reference is $D=-34.6$ LUFS, the median level of the $1499$ seconds covered by subtitles ($799$ cues, applied offset $0.0$ s). Timecodes belong to this release.

\paragraph{2. Passage.} $1{:}00{:}39$ to $1{:}20{:}51$, a span of $1212$ s, bounded by the longest intervals without a qualifying recovery window.

\paragraph{3. Observation.} The monograph defines exceedance $E(t)=L(t)-D$, occupation as $E\ge\delta=10$ dB, and a recovery window as at least $5$ s at or below $D$. Terminology matters here: a \emph{recovery window} is a stretch of at least $5$ s at or below $D$, and what the monograph calls a recovery gap is the complementary interval \emph{without} a qualifying window, not a quiet interval. In the passage there are three qualifying windows, of $9$, $7$, and $13$ s, the last two separated by a single second above $D$.
\begin{center}
\small
\begin{tabular}{llr}
\toprule
interval & kind & seconds\\
\midrule
$1{:}00{:}39$ to $1{:}03{:}18$ & no qualifying window & $159$\\
$1{:}03{:}18$ to $1{:}03{:}27$ & recovery window & $9$\\
$1{:}03{:}27$ to $1{:}09{:}00$ & no qualifying window & $333$\\
$1{:}09{:}00$ to $1{:}09{:}07$ & recovery window & $7$\\
$1{:}09{:}07$ to $1{:}09{:}08$ & above $D$ & $1$\\
$1{:}09{:}08$ to $1{:}09{:}21$ & recovery window & $13$\\
$1{:}09{:}21$ to $1{:}20{:}51$ & no qualifying window & $690$\\
\bottomrule
\end{tabular}
\end{center}
The intervals sum to $1212$ s. The three windows total $29$ s, not the $30$ s stated in the monograph, whose figure is the length of the two short intervals between the three long ones ($9+21$ s) and so includes the $1$ s above $D$ and counts two windows where there are three. Counting every second at or below $D$, not only those in windows of at least $5$ s, there are $80$ such seconds in the passage, in $32$ separate runs, the longest $13$ s. The passage is therefore not free of quiet moments; its quiet moments are fragmented. These are acoustic opportunities for recovery under a stated rule, not measurements of a listener's recovery. Six occupation blocks (level-defined, merge gap $g=5$ s, minimum length $60$ s) lie inside the passage.
\begin{center}
\small
\begin{tabular}{llrrrrr}
\toprule
start & end & seconds & median $E$ & $Q_{95}(E)$ & throttle $T_b$ & subtitled s\\
\midrule
$1{:}00{:}48$ & $1{:}02{:}29$ & $101$ & $+16.2$ & $+19.1$ & $9.1$ & $10$\\
$1{:}04{:}24$ & $1{:}05{:}56$ & $92$ & $+16.8$ & $+20.7$ & $10.7$ & $0$\\
$1{:}06{:}46$ & $1{:}08{:}05$ & $79$ & $+16.9$ & $+22.8$ & $12.8$ & $10$\\
$1{:}09{:}24$ & $1{:}11{:}34$ & $130$ & $+12.5$ & $+16.5$ & $6.5$ & $48$\\
$1{:}13{:}45$ & $1{:}14{:}58$ & $73$ & $+15.5$ & $+18.8$ & $8.8$ & $13$\\
$1{:}15{:}09$ & $1{:}19{:}53$ & $284$ & $+16.0$ & $+20.4$ & $10.4$ & $30$\\
\bottomrule
\end{tabular}
\end{center}
Levels are in dB relative to $D$ and the throttle depth is $T_b=\max(0,Q_{95}(E\mid b)-\delta)$. The six blocks total $759$ s, which is $62.6\%$ of the $1212$ s of the span, the figure the monograph reports; before $1{:}00{:}39$ the share is $6.8\%$. Over the whole film, $30.8\%$ of audible seconds lie at least $10$ dB above $D$.

\paragraph{4. Annotation.} Everything above lies in the appearance layer of Chapter~\ref{ch:layers}, and specifically in sound. The monograph carries no shot boundaries and no history or state annotation for this passage: nothing says what happens in the story during any block. Subtitled seconds mark when speech occurs and not how loud it is, and the per-second level is the level of the whole mix, so a subtitled second inside a block records speech occurring in the loud material, not loud speech. The annotation step of the protocol therefore stops at the first layer, and every statement below about plot function would be an addition not supported by the source.

\paragraph{5. Claim types.} The table entries are observations (measured, in the monograph's grading). The monograph's statement that the loud plateau reflects compression or limiting is a hypothesis about production, graded as such there. The listener behaviour that motivated the study, turning the volume down for one to two minutes or more, is a single retrospective report without timestamps. It is the motivation of the study and not evidence in it.

\paragraph{6. Formal representation.} Three constructions of the earlier chapters apply, with differences that matter.

\emph{Exceedance and arrangement.} \Cref{prop:arrangement} says the histogram of levels does not determine the arrangement of exceedance. The real data show it concretely: the film-level figure of $30.8\%$ is the same whether occupation is spread evenly or concentrated, and here it is concentrated, with $13\%$ to $40\%$ of each ten-minute segment before $1{:}00{:}00$ and $62\%$ and $71\%$ in the next two. A further point is that the count of blocks depends on a parameter the raw definition of Section~\ref{sec:occupation} lacks. The monograph merges runs separated by at most $g=5$ s. The last two blocks are separated by $11$ s ($1{:}14{:}58$ to $1{:}15{:}09$). Rerunning the merge on the per-second occupied set gives six blocks inside the passage at $g=5$, five at $g=10$, and four at $g=11$ and $g=12$. The count does not fall only when the last two blocks merge, because merging also absorbs the short occupied episodes between blocks, which are not in Table~7.3: at $g=10$ the second and third blocks join across such episodes into one stretch from $1{:}03{:}44$ to $1{:}08{:}37$ (seconds $3824$ to $4117$), and at $g=11$ the last two blocks merge into one of $368$ s. The number of occupation blocks is therefore a statement about the exceedance set and a choice of $g$, and a report should give both.

\emph{Uniform attenuation.} The throttle depth $T_b$ is the uniform attenuation $a=T_b$ that brings the $95$th percentile of exceedance in block $b$ to the tolerance. By definition of the percentile, at least $95\%$ of the block's seconds then satisfy $E-a\le\delta$; the small residual fraction ($5.1\%$ to $6.5\%$ in the six blocks, recomputed from the per-second table) is built into the rule and is not an independent result. By \Cref{prop:mufflerizer} the exceedance sets are nested in $a$, so larger attenuations clear more. All such occupation figures use the \emph{fixed original} reference $D$. A uniform gain lowers dialogue and effects together, so measured against dialogue that has itself been lowered by $a$ the exceedance is unchanged, which is the listener's dilemma in the monograph. A reduction in occupation under a fixed reference is a statement about absolute level, and it says nothing about intelligibility or listening quality. The six values, $6.5$ to $12.8$ dB, are far above the $4$ dB cue-safe limit of the synthetic scene. That limit came from a stipulated floor on named cue sounds, and nothing in the source identifies cue sounds, so \Cref{prop:cuesafe} cannot be evaluated here. The quantity it needs, which sounds are the only evidence of a plot event, is the one the data lack.

\emph{Dialogue as a stem.} In the synthetic scene dialogue was a separate stem and was left unchanged. In the measurement it is not: the level is that of the whole mix. A global attenuation by $a$ lowers dialogue by $a$ as well, so $\rho_{\mathrm{dia}}$ fails in every block by between $6.5$ and $12.8$ dB. The stem-wise policy of Section~\ref{sec:whisper} needs a proxy stem. The monograph supplies one: subtitled seconds have a median centre-channel advantage of $16.0$ dB. The measured centre advantage shows how far this works, with a caution about what is measured. The seconds are subtitle-associated seconds, each a mixture of whatever sounds occur in that second, not isolated dialogue; channel location is not sound function. For subtitle-associated seconds over the whole film the median is $16.0$ dB, but \emph{inside the blocks} it is $+8.8$ dB at the median (quartiles $+3.9$ and $+12.0$, tenth percentile $-0.5$), against $+17.0$ dB outside them. Dialogue inside loud material is much less centre-dominant than dialogue elsewhere. Under the idealization that mix power is the sum of centre power $C$ and the rest $N$ with advantage $x=C-N$ in dB, removing $N$ entirely changes the mix level by
\[
10\log_{10}\bigl(1+10^{-x/10}\bigr)\ \mathrm{dB}\quad\text{(a decrease)},
\]
which is $0.11$ dB at $x=16$ but $0.54$ dB at $x=8.8$ and $1.5$ dB at $x=3.9$. Applied second by second to the $128$ subtitled seconds inside blocks ($8.5\%$ of all subtitled dialogue), the median decrease is $0.54$ dB and one in ten of these seconds falls by $3.3$ dB or more. The mix level of those seconds therefore contains a substantial non-centre part. Removing the non-centre channels might remove masking material, useful speech, or both, and the levels cannot say which. The reported changes therefore describe the channel mixture under the stated idealization, and dialogue preservation is left unresolved. (The idealization is approximate: the pipeline's RMS level of the mix is, on median, $3.0$ dB below the power sum of its centre and non-centre levels, and the BS.1770 weights are ignored.)

The same data bound what channel-wise attenuation can do for occupation. The passage has $817$ occupied seconds, $759$ of them in blocks. In the idealization, lowering the non-centre channels by $6$, $12$, or $20$ dB leaves $579$, $480$, or $441$ of the $759$ block seconds occupied, and removing them entirely leaves $435$ (about $57\%$). So even a perfect separation of the non-centre material would leave most of the blocks occupied: in $57.8\%$ of the passage's occupied seconds the non-centre channels are louder than the centre, yet the centre alone is still well above dialogue level. The loud material is not only in the surround and effects channels. This is why the monograph's downmix also compresses, and it shows that a stem-wise policy by channel is not by itself an occupation remedy for this passage.

\paragraph{7. Alternative edits.} Three transformations of increasing separation are on the monograph's own list: global gain, channel-based rebalancing with compression, and class-specific attenuation. Against the relations of Section~\ref{sec:altversions}:
\begin{center}
\small
\begin{tabular}{p{3.2cm}p{3.1cm}p{3.1cm}p{2.4cm}}
\toprule
 & occupation & $\rho_{\mathrm{dia}}$ & $\rho_{\mathrm{cue}}$\\
\midrule
global gain by $T_b$ & at least $95\%$ of block seconds cleared (computed) & fails by $T_b$ (computed) & not evaluable\\
channel rebalancing & idealized: removing all non-centre content leaves $435$ of $759$ block seconds occupied & median mix change $0.54$ dB in block speech, $\ge3.3$ dB in one in ten & not evaluable\\
class-specific & needs a detector with a stated error rate & depends on classifier & needs a cue list\\
\bottomrule
\end{tabular}
\end{center}
The compression stage of the channel rebalancing changes the mix everywhere and not only in blocks, as the monograph notes, so $\rho_{\mathrm{ord}}$-type relations on the whole film and not only on the passage would also have to be stated.

\paragraph{8. Evaluation.} The monograph specifies the study that would settle what the table cannot: logged volume behaviour under at least two playback conditions, one of them the dialogue-forward downmix, with comprehension questions after each viewing. In the vocabulary of this book, that study compares two versions that differ by a transformation whose contract is the list above, and it measures the preferred experience, the fourth claim of Section~\ref{sec:whisper}, which no computation here addresses.

\paragraph{What the real passage showed about the framework.}
\begin{itemize}
\item The occupation and recovery constructions transferred without change, and the monograph's definitions are refinements of them: a tolerance $\delta=10$ dB where the synthetic scene used $\mu=3$ dB, a merge gap, a minimum block length, and a recovery window of $5$ s at or below $D$. The propositions apply to the raw exceedance set, and statements about blocks must name the merge parameters.
\item The contract machinery was only half usable. $\rho_{\mathrm{dia}}$ could be assessed for global gain and, with the per-second table, for a channel policy, where it turned out that the dialogue inside blocks is far less centre-dominant than the film-wide figure suggests. $\rho_{\mathrm{cue}}$ and $\rho_{\mathrm{sal}}$ could not be assessed at all, because they depend on plot function, and the source records no plot.
\item The time-indexed analysis of Section~\ref{sec:timeindexed} had five columns. For this passage only observation and, partly, formal representation could be filled. Annotation, alternative edit, and evaluation needed information (shots, story events, listener data) that the source lacks. The synthetic scene looked easy to analyse in part because its events were stipulated together with their causal roles. Real passages do not come with that.
\item The response-model machinery of Chapter~\ref{ch:equivalence} had no data. One retrospective report is not a response model, so no sensitivity matrix or effective dimension can be estimated yet.
\item The monograph's grading of claims as established, measured, hypothesised, or rejected is the same discipline as Appendix~\ref{app:register}, and the transcribed numbers here are in its grade of measured.
\end{itemize}

\section{A shorter passage: speech timing against sound level}\label{sec:excerpt}

The previous section left the story and picture columns empty. This section takes a $130$ s excerpt inside the same passage, from $1{:}09{:}24$ to $1{:}11{:}34$ (seconds $4164$ to $4293$ of the per-second table), chosen because it lies wholly inside one occupation block and has a substantial number of subtitle-associated seconds. An earlier note gave this stretch by a clock reading that could be taken in two ways; the time here is hours, minutes and seconds from the start of the film. The section uses only the \emph{timing} of the subtitle cues, in the file \texttt{data/cues\_1h09\_1h11.csv}, and does not reproduce or interpret what is said. The subtitle-associated flag of the monograph marks a second when the cues cover at least half of it.

\paragraph{Observation.} All $130$ seconds lie in a block, and $106$ of them ($81.5\%$) are occupied under $\delta=10$ dB. The minimum exceedance over the excerpt is $5.3$ dB and the median $12.6$ dB, so no second of the excerpt is at or below the dialogue reference $D$, and no recovery window can occur. The flag marks $48$ seconds. They fall in short runs, with one long speech-free stretch: from $1{:}09{:}33$ the flag is absent for $41$ s, up to $1{:}10{:}13$. Table~\ref{tab:excerpt} gives ten-second segments.

\begin{table}[ht]
\centering\small
\caption{Excerpt $1{:}09{:}24$ to $1{:}11{:}34$ in ten-second segments, from the per-second table and the cue timings. ``Sub.'' is the number of subtitle-associated seconds, ``Occ.'' the number of occupied seconds, $E$ the median exceedance over $D=-34.6$ LUFS, and ``adv.'' the median centre advantage.}\label{tab:excerpt}
\begin{tabular}{lrrrr}
\toprule
segment start & Occ. & Sub. & median $E$ (dB) & median adv. (dB)\\
\midrule
$1{:}09{:}24$ & 8 & 3 & 10.9 & $-7.8$\\
$1{:}09{:}34$ & 10 & 0 & 15.1 & $-6.4$\\
$1{:}09{:}44$ & 9 & 0 & 14.7 & $-6.6$\\
$1{:}09{:}54$ & 7 & 0 & 12.4 & $-12.7$\\
$1{:}10{:}04$ & 10 & 0 & 13.2 & $-2.3$\\
$1{:}10{:}14$ & 10 & 10 & 12.9 & $10.4$\\
$1{:}10{:}24$ & 2 & 4 & 8.8 & $0.1$\\
$1{:}10{:}34$ & 8 & 5 & 13.9 & $9.0$\\
$1{:}10{:}44$ & 8 & 6 & 11.9 & $4.5$\\
$1{:}10{:}54$ & 9 & 6 & 11.8 & $6.2$\\
$1{:}11{:}04$ & 10 & 3 & 12.5 & $0.4$\\
$1{:}11{:}14$ & 8 & 8 & 11.2 & $8.2$\\
$1{:}11{:}24$ & 7 & 3 & 12.5 & $1.7$\\
\bottomrule
\end{tabular}
\end{table}

\paragraph{What the table shows.} Two features are visible without any interpretation of content. First, the exceedance does not follow speech density: the five segments with at most three flagged seconds have medians between $10.9$ and $15.1$ dB, and the speech-dense segments between $8.8$ and $13.9$ dB. Second, the centre advantage does: in the $41$ s speech-free stretch the segment medians are between $-12.7$ and $-6.4$ dB, so the non-centre channels dominate, and once speech begins the medians are mostly positive. Over the excerpt the median centre advantage is $+9.3$ dB in subtitle-associated seconds (three of the $48$ are negative) and $-4.8$ dB elsewhere. The acoustic organization of the excerpt is therefore a loud, non-centre-dominated stretch followed by a loud stretch in which speech is carried on top of material that remains well above $D$. The lowest exceedance (segment starting $1{:}10{:}24$, with only two occupied seconds) is a dip within the block, not a recovery, since it stays $5.3$ dB above $D$ at its minimum.

\paragraph{Which of the five columns can be filled.} Observation: yes, as above. Formal representation: partly. The excerpt is a path in the product $C\times A\times L\times E$ of Chapter~\ref{ch:scenario} whose acoustic coordinate is the exceedance sequence of Table~\ref{tab:excerpt}, and the occupation relation of Section~\ref{sec:occupation} holds on it for $106$ of $130$ seconds. Annotation, alternative edit and evaluation: no. They need what the source does not record: cuts, camera movement and visible action (for pacing); what is learned by whom (for narrative change); and which sounds are the only evidence of an event (for $\rho_{\mathrm{cue}}$). The next table lists, for this excerpt, exactly what a video excerpt would add. It is a worksheet, not a finding.

\begin{center}
\small
\begin{tabular}{p{3.6cm}p{4.2cm}p{4.2cm}}
\toprule
layer & question for the excerpt & source needed\\
\midrule
visual pacing & do cut rate and camera motion differ between the speech-free stretch and the speech-dense stretch? & shot boundaries, frame-differencing or manual log\\
narrative change & does the knowledge state (Chapter~\ref{ch:layers}) change at the start of the speech-dense stretch, or is it constant across both? & annotated story events with time stamps\\
acoustic organization & do occupation and exceedance track either of the above? & the table above (available)\\
\bottomrule
\end{tabular}
\end{center}

With these in hand the connection can be tested rather than asserted. If the shot-change rate is high in the speech-free stretch and the state layer is constant, the excerpt would exhibit visual and acoustic agitation without narrative change; if the state layer changes at the onset of speech while the exceedance does not, the acoustic organization is independent of the narrative change. Either outcome is a statement about one $130$ s excerpt, and neither bears on the film as a whole or on any viewer.


\section*{Exercises}

\begin{exercise}
Choose a film that uses concealed cuts and identify the seams. For each, estimate informally which of $\delta_\Sc$ and $\delta_\Img$ is small and state what kind of object conceals the discontinuity.
\end{exercise}

\begin{exercise}
For a film with an unreliable narrator, set up a sheaf-theoretic model with local sections given by what the narrator reports on each scene. State the condition under which the narrator's account is globally consistent and the condition under which it is not.
\end{exercise}

\begin{exercise}
Extend the synthetic scene by adding a third character $C$ who knows the letter is genuine. Write the story state, the typing of the reveals, and the number of valid orderings of three reveal events.
\end{exercise}

\begin{exercise}
Carry out the protocol of Section~\ref{sec:protocol} on a passage of a film of your choice, of at most two minutes, supplying edition and timecodes. Mark each claim as observation, reading, or hypothesis, and state one alternative edit with the relation it preserves and the relation it breaks.
\end{exercise}

\begin{exercise}
In the time-indexed analysis of Section~\ref{sec:timeindexed}, identify the segments whose annotation would change if the audience were told at $t=18$ that $A$ knows ($V_4$). Which formal objects in the table change, and which are unchanged?
\end{exercise}

\chapter{Open Problems}
\label{ch:open}

The book has proposed a vocabulary and proved what follows from it. It has also, repeatedly, marked claims as principles and left questions to experiment. This chapter collects the principal open problems, sorted by the kind of work needed to settle them. Each is stated so that a result would be recognizable.

\section{Mathematical problems}

\emph{A sheaf-theoretic account of editing.} Chapter~\ref{ch:sheaves} treated the sheaf of takes over a time interval, and Chapter~\ref{ch:descent} treated coverage over a scene. A unified treatment would place both on a single site whose objects are spacetime regions, with a sheaf of scene states and a morphism of sheaves into images. The open problem is to characterize the image of rendering in the category of sheaves in terms of cohomological conditions, so that a statement of the form ``this family of images lifts to a consistent world'' has a cohomological criterion.

\emph{Higher obstructions.} The class in $H^1$ of Chapter~\ref{ch:cohomology} detects the failure of a single correction to be explained by local data. Obstructions to the existence of descent data in the stronger sense of Chapter~\ref{ch:descent} live in higher cohomology or in nonabelian cohomology with coefficients in groups of isomorphisms of the local objects. Whether such higher obstructions have natural filmic interpretations, beyond identity-tracking across several views, is unsettled.

\emph{Typed generation.} The Galois connection of Chapter~\ref{ch:scripts} and the typing of Chapter~\ref{ch:types} describe satisfiability. The efficient generation of films from a rich type system, with decidable checking and tractable sampling, is a problem at the boundary of programming languages and probabilistic inference. For the linear fragment of Chapter~\ref{ch:linear}, the existence of balanced runs is a Petri-net reachability problem, and the cost of sampling uniformly from balanced runs is not addressed here.

\emph{Quantitative preservation.} \Cref{prop:compose-contracts} gives a worst-case bound on composite drift. Under distributional assumptions on the viewer's dissimilarity, one expects tighter bounds, and a theory of preservation for stochastic transformations would give them.

\section{Empirical problems}

\emph{Conditional dimension.} Principle~\ref{pr:conditional} says that the number of perceptually independent variations is a conditional quantity. The experiment is to estimate the singular spectrum of the sensitivity matrix for several scenes, tasks, and viewer groups, and to ask whether the spectra are similar enough across scenes to justify any common figure. A rapidly decaying spectrum for most scenes would support the cue-repertoire conjecture of Chapter~\ref{ch:equivalence}, and a flat spectrum would count against it.

\emph{The geodesic principle.} Principle~\ref{pr:geodesic} asserts that natural camera movements are critical points of an action with a perceptual metric. The test is to fit the metric to viewer judgments of naturalness for a set of synthetic camera paths and to ask whether the fitted metric makes held-out judgments predictable. A failure of predictive power for every metric in the Fisher--Rao family would count against the principle.

\emph{Recovery and occupation.} The temporal organization of acoustic occupation, formalized in Chapter~\ref{ch:contracts}, supports the hypothesis that arrangements with equal level histograms but different block structure differ in listener fatigue and comprehension. Whether they do, and whether any benefit of recovery intervals depends on their length and placement, are questions for controlled listening studies with the response variables of Chapter~\ref{ch:equivalence} specified in advance.

\emph{Authority of the history.} Principle~\ref{pr:authority} states that in a persistent generative world the history must be authoritative at commit. A systematic comparison of generative systems with and without an explicit history layer, on tasks requiring long-range consistency, would test whether the failures attributed to its absence are in fact explained by it.

\section{Design problems}

\emph{Interfaces for hierarchical direction.} Chapter~\ref{ch:direction} showed that prefix-free command sets parse uniquely and that commands generally do not commute. The design of notations that expose conflicts, preview realizations at the user's level of abstraction, and support undo without silently invalidating later commitments is an open problem in human-computer interaction.

\emph{Presentation contracts.} The same film may be offered under different presentation contracts, each preserving different things: theatrical form, quiet home viewing, analytical study. How to let a viewer move between them without losing place or narrative context is a problem about synchronization of several versions, which is a cover of the film by its presentations and, in the language of Chapter~\ref{ch:selective}, a problem of branching with shared segments.

\section{Limits of the approach}

The formalism treats a film as a path, a word, or a family of local data, and a viewer as a source of dissimilarity judgments. These are idealizations, and they exclude much of what criticism attends to. Interpretation is contestable even when the measurements are reproducible. A theorem about a model of a film is a theorem about the model, and the book's principles are the places where it ventures a claim about film itself, together with a statement of what would count against the claim. If the programme is to be useful, it will be because the models sharpen disagreements into questions that evidence can answer, and because they make clear what a generative system must be told if it is to avoid errors that no amount of rendering quality can repair.


\appendix
\chapter{Category Theory in Brief}
\label{app:cat}

This appendix collects the definitions used in the main text. Standard references are~\cite{maclane1998,awodey2010}.

\section*{Categories and functors}

A \emph{category} $\mathbf{C}$ consists of a class of objects, for each ordered pair of objects $x,y$ a set $\Hom(x,y)$ of morphisms, an associative composition $\circ:\Hom(y,z)\times\Hom(x,y)\to\Hom(x,z)$, and for each object an identity morphism $\id_x$ satisfying $f\circ\id_x=f=\id_y\circ f$. A category with one object is the same as a monoid.

A \emph{functor} $F:\mathbf{C}\to\mathbf{D}$ assigns to each object $x$ an object $Fx$ and to each morphism $f:x\to y$ a morphism $Ff:Fx\to Fy$ with $F(g\circ f)=Fg\circ Ff$ and $F\id_x=\id_{Fx}$. It is \emph{faithful} if it is injective on each $\Hom(x,y)$ and \emph{full} if it is surjective on each. A contravariant functor reverses the direction of morphisms and is a functor $\mathbf{C}^{\op}\to\mathbf{D}$.

A \emph{natural transformation} $\eta:F\Rightarrow G$ between functors $F,G:\mathbf{C}\to\mathbf{D}$ assigns to each object $x$ a morphism $\eta_x:Fx\to Gx$ such that $\eta_y\circ Ff=Gf\circ\eta_x$ for every $f:x\to y$.

\section*{Limits and fiber products}

Given functors $F:\mathbf{A}\to\mathbf{C}$ and $G:\mathbf{B}\to\mathbf{C}$, the \emph{fiber product} $\mathbf{A}\times_{\mathbf{C}}\mathbf{B}$ has as objects pairs $(a,b)$ with $Fa=Gb$ and as morphisms pairs $(f,g)$ with $Ff=Gg$. It is the limit of the cospan $\mathbf{A}\to\mathbf{C}\leftarrow\mathbf{B}$ in the category of small categories.

\section*{Monoidal categories}

A \emph{monoidal category} is a category with a bifunctor $\otimes$ and a unit object $\mathbf{1}$ together with natural isomorphisms for associativity and unit, satisfying the pentagon and triangle coherence conditions. The \emph{interchange law} $(g\circ f)\otimes(g'\circ f')=(g\otimes g')\circ(f\otimes f')$ is part of the statement that $\otimes$ is a functor on the product category. A \emph{lax monoidal functor} $F$ between monoidal categories comes with a morphism $Fx\otimes Fy\to F(x\otimes y)$ natural in $x,y$ and a morphism $\mathbf{1}\to F\mathbf{1}$ satisfying coherence. When the morphisms are identities or isomorphisms the functor is \emph{strong}.

\section*{Adjunctions and Galois connections}

Functors $F:\mathbf{C}\to\mathbf{D}$ and $G:\mathbf{D}\to\mathbf{C}$ are \emph{adjoint}, written $F\dashv G$, if there is a bijection $\Hom_{\mathbf{D}}(Fx,y)\cong\Hom_{\mathbf{C}}(x,Gy)$ natural in $x$ and $y$. When $\mathbf{C}$ and $\mathbf{D}$ are posets, this says $Fx\le y\iff x\le Gy$, which is a monotone Galois connection. The antitone connection of Chapter~\ref{ch:scripts} is the same condition with one of the orders reversed.

\chapter{Sheaves and Cohomology in Brief}
\label{app:sheaf}

Standard references are~\cite{maclane1992,bott1982}.

\section*{Sites and sheaves}

For a topological space $X$, a presheaf of sets is a functor $\Open(X)^{\op}\to\Set$, and a sheaf is a presheaf satisfying the locality and gluing axioms of Chapter~\ref{ch:sheaves}. A sheaf of abelian groups takes values in $\Ab$. The sheaf condition is equivalent to the statement that, for every open cover $\{U_i\}$ of $U$, the diagram
\[
F(U)\longrightarrow\prod_iF(U_i)\rightrightarrows\prod_{i,j}F(U_{ij})
\]
is an equalizer, where the two parallel maps restrict from $U_i$ and from $U_j$ to $U_{ij}$.

A \emph{morphism of sheaves} is a natural transformation. For a continuous map $\Ren:\Sc\to\Img$ and the sheaves of continuous maps into $\Sc$ and $\Img$ over a base $T$, post-composition with $\Ren$ defines a morphism of sheaves.

\section*{Cech cohomology in low degree}

For an open cover $\mathcal{U}=(U_i)_{i\in I}$ with $I$ totally ordered and a sheaf of abelian groups $A$, the Cech complex is
\[
C^0(\mathcal{U};A)\xrightarrow{d^0}C^1(\mathcal{U};A)\xrightarrow{d^1}C^2(\mathcal{U};A)\to\cdots
\]
with the differentials given in Chapter~\ref{ch:cohomology}. The identity $d^1d^0=0$ is verified directly: for $a\in C^0$,
\[
(d^1d^0a)_{ijk}=(a_k-a_j)-(a_k-a_i)+(a_j-a_i)=0,
\]
with all terms understood as restricted to $U_{ijk}$. The group $H^1(\mathcal{U};A)$ is the quotient $\ker d^1/\operatorname{im}d^0$. A cover is called \emph{acyclic for $A$} if $H^p(U_{i_0\cdots i_k};A)=0$ for all $p\ge1$ on every finite intersection; for such a cover the Cech groups agree with the sheaf cohomology of the space (Leray's theorem).

\section*{Locally constant sheaves}

The constant sheaf $\underline{A}$ assigns to $U$ the group of locally constant functions $U\to A$. On a connected open set its sections are the constants. On a disjoint union of connected components its sections are an arbitrary choice of constant per component, which is how the calculations of Chapter~\ref{ch:cohomology} proceed.

\chapter{Type Theory in Brief}
\label{app:types}

Standard references are~\cite{martinlof1984,hottbook2013,girard1987}.

\section*{Judgments and contexts}

The basic judgments are $\Gamma\vdash A\ \mathsf{type}$, $\Gamma\vdash a:A$, and, for equality, $\Gamma\vdash a\equiv b:A$. A context $\Gamma=(x_1:A_1,x_2:A_2,\dots,x_n:A_n)$ is a list in which each $A_k$ is a type in the preceding context. Contexts support \emph{weakening}, the addition of an unused hypothesis, and \emph{substitution}, the replacement of a variable by a term of its type.

\section*{Dependent sums and products}

The dependent product $\Pi_{x:A}B(x)$ has introduction by $\lambda$-abstraction and elimination by application. The dependent sum $\Sigma_{x:A}B(x)$ has introduction by pairs $(a,b)$ with $b:B(a)$ and elimination by projections. When $B$ does not depend on $x$ these reduce to function and product types.

\section*{Identity types}

For $a,b:A$ the identity type $a=_Ab$ has an introduction rule $\mathsf{refl}_a:a=_Aa$ and an elimination rule, path induction. In the homotopy interpretation, a type is a space, a term is a point, and a term of $a=_Ab$ is a path from $a$ to $b$. A type is a \emph{set} if any two parallel paths are equal.

\section*{Linear logic}

In intuitionistic linear logic a context is split into an unrestricted part $\Gamma$ and a linear part $\Delta$. The linear implication $A\multimap B$ has introduction by a function that consumes exactly one $A$ and elimination by application to a term built from a disjoint part of the linear context. Weakening and contraction are admissible on $\Gamma$ and not on $\Delta$. The multiset semantics used in Chapter~\ref{ch:linear} is the Petri-net reading of the multiplicative fragment with atomic propositions as places.

\chapter{Register of Claims and Their Status}
\label{app:register}

The book mixes three kinds of statement: results proved or cited, definitions and stipulated examples, and hypotheses about perception and craft. The register below lists the hypotheses, with the part of each that is mathematics, the part that is empirical, and what would count against it. Every counting and numerical claim of the worked examples (the counts of edits, orderings, and realizations, the occupation tables, the singular values, and the remainder bounds) is reproduced by the script \texttt{verify.py} distributed with the source, computed independently of the text; a disagreement between script and text indicates an error in one of them. Numerical values in worked examples (the forged-letter scene, the response model, the sound levels) are stipulated for illustration and are not measurements.

\begin{center}
\small
\begin{tabular}{p{2.6cm}p{3.2cm}p{3.5cm}p{3.6cm}}
\toprule
principle & mathematical part & empirical part & would count against it\\
\midrule
Geodesic principle (\Cref{pr:geodesic}) & Euler--Lagrange equation of \Cref{prop:EL} & that natural camera motion approximately extremizes such an action & natural paths failing the equation for every metric in the family\\[4pt]
Obstruction principle (\Cref{pr:obstruction}) & equivalence of consistency with a vanishing class, once corrections lie in an abelian sheaf & that a film's continuity constraints can be modeled that way & a film whose constraints do not fit, or a repair by local edits of a nonzero class\\[4pt]
Authority principle (\Cref{pr:authority}) & none beyond the layer definitions & that frame-only systems fail because they lack a history layer & failures of such systems unrelated to missing history, or success without the layer\\[4pt]
Shared-segment principle (\Cref{pr:shared}) & typing in an intersection of contexts & that authors can state the contexts of a branching film & shared segments that are well typed on some branches and still acceptable on all\\[4pt]
Conditional dimension (\Cref{pr:conditional}) & first-order effective dimension (\Cref{prop:effdim}) and its remainder (\Cref{prop:taylor}) & that the number of independent variations depends on scene, viewer, task, tolerance, and budget & estimates that are stable across scenes and populations to a degree justifying a single figure\\
\bottomrule
\end{tabular}
\end{center}

Two further conjectures are stated in the text without being named as principles. The conjecture of Chapter~\ref{ch:equivalence} that $d_{\mathrm{eff}}$ is small relative to the number of controls near typical production settings is empirical, and a slowly decaying singular spectrum would count against it. The claim of Chapter~\ref{ch:contracts} that listener preference is not determined by reduced occupation, preserved intelligibility, and preserved cues is also empirical, and is stated as a warning against inferring the fourth from the first three.


\backmatter
\chapter{Notation}
\label{app:notation}

\noindent
$\Sc$ scenario space; $\Ren:\Sc\to\Img$ rendering; $\Img$ image space; $\Path(\Sc)$ Moore path category; $\dur$ duration; $\Sigma^*$ free monoid of shots; $\shuffle$ shuffle product; $\Open(X)$ poset of open sets; $\underline{A}$ constant sheaf; $H^p(\mathcal{U};A)$ Cech cohomology; $E$ events; $\Hist$ histories; $\Adm$ admissible histories; $M$ state fold; $\sim$ future equivalence; $P_A$ conditioning of $P$ on $A$; $\KL$ Kullback--Leibler divergence; $d_\theta$ dissimilarity in situation $\theta$; $\approx_\varepsilon$ tolerance relation; $d_{\mathrm{eff}}$ effective dimension; $\Th,\Mod$ theory and model operators.

\begin{thebibliography}{99}

\bibitem{abramsky2011} S.~Abramsky and A.~Brandenburger. The sheaf-theoretic structure of non-locality and contextuality. \emph{New Journal of Physics}, 13:113036, 2011.

\bibitem{amari2016} S.~Amari. \emph{Information Geometry and Its Applications}. Springer, 2016.

\bibitem{awodey2010} S.~Awodey. \emph{Category Theory}. 2nd ed. Oxford University Press, 2010.

\bibitem{bazin1967} A.~Bazin. \emph{What Is Cinema?}, vol.~1. Translated by H.~Gray. University of California Press, 1967.

\bibitem{bordwell1985} D.~Bordwell. \emph{Narration in the Fiction Film}. University of Wisconsin Press, 1985.

\bibitem{bott1982} R.~Bott and L.~W.~Tu. \emph{Differential Forms in Algebraic Topology}. Springer, 1982.

\bibitem{brown1982} K.~S.~Brown. \emph{Cohomology of Groups}. Springer, 1982.

\bibitem{bruce2024genie} J.~Bruce et al. Genie: Generative interactive environments. In \emph{Proceedings of the 41st International Conference on Machine Learning}, 2024.

\bibitem{chandy1985} K.~M.~Chandy and L.~Lamport. Distributed snapshots: Determining global states of distributed systems. \emph{ACM Transactions on Computer Systems}, 3(1):63--75, 1985.

\bibitem{deleuze1986} G.~Deleuze. \emph{Cinema 1: The Movement-Image}. University of Minnesota Press, 1986.

\bibitem{deleuze1989} G.~Deleuze. \emph{Cinema 2: The Time-Image}. University of Minnesota Press, 1989.

\bibitem{eisenstein1949} S.~Eisenstein. \emph{Film Form: Essays in Film Theory}. Edited and translated by J.~Leyda. Harcourt, Brace, 1949.

\bibitem{flyxion2026occupation} Flyxion. \emph{Acoustic Occupation: Salience, Recovery, and the Temporal Organization of Film Sound}. Working manuscript, September 2026.

\bibitem{girard1987} J.-Y.~Girard. Linear logic. \emph{Theoretical Computer Science}, 50(1):1--102, 1987.

\bibitem{ha2018worldmodels} D.~Ha and J.~Schmidhuber. World models. arXiv:1803.10122, 2018.

\bibitem{hafner2020dreamer} D.~Hafner, T.~Lillicrap, J.~Ba, and M.~Norouzi. Dream to control: Learning behaviors by latent imagination. In \emph{International Conference on Learning Representations}, 2020.

\bibitem{hopcroft2006} J.~E.~Hopcroft, R.~Motwani, and J.~D.~Ullman. \emph{Introduction to Automata Theory, Languages, and Computation}. 3rd ed. Pearson, 2006.

\bibitem{hottbook2013} The Univalent Foundations Program. \emph{Homotopy Type Theory: Univalent Foundations of Mathematics}. Institute for Advanced Study, 2013.

\bibitem{lamport1978} L.~Lamport. Time, clocks, and the ordering of events in a distributed system. \emph{Communications of the ACM}, 21(7):558--565, 1978.

\bibitem{lee2013} J.~M.~Lee. \emph{Introduction to Smooth Manifolds}. 2nd ed. Springer, 2013.

\bibitem{luce1956} R.~D.~Luce. Semiorders and a theory of utility discrimination. \emph{Econometrica}, 24(2):178--191, 1956.

\bibitem{maclane1998} S.~Mac~Lane. \emph{Categories for the Working Mathematician}. 2nd ed. Springer, 1998.

\bibitem{maclane1992} S.~Mac~Lane and I.~Moerdijk. \emph{Sheaves in Geometry and Logic}. Springer, 1992.

\bibitem{martinlof1984} P.~Martin-L\"of. \emph{Intuitionistic Type Theory}. Bibliopolis, 1984.

\bibitem{metz1974} C.~Metz. \emph{Film Language: A Semiotics of the Cinema}. Translated by M.~Taylor. Oxford University Press, 1974.

\bibitem{penrose1992} R.~Penrose. On the cohomology of impossible figures. \emph{Leonardo}, 25(3/4):245--247, 1992.

\bibitem{selinger2011} P.~Selinger. A survey of graphical languages for monoidal categories. In \emph{New Structures for Physics}, Lecture Notes in Physics 813, pages 289--355. Springer, 2011.

\bibitem{shalizi2001} C.~R.~Shalizi and J.~P.~Crutchfield. Computational mechanics: Pattern and prediction, structure and simplicity. \emph{Journal of Statistical Physics}, 104:817--879, 2001.

\bibitem{valevski2024gamengen} D.~Valevski, Y.~Leviathan, M.~Arar, and S.~Fruchter. Diffusion models are real-time game engines. arXiv:2408.14837, 2024.

\end{thebibliography}


\end{document}
